% Lesson 20 — Reading: Texts on managing personal finances — Anglais 1ère A — Cours signé OnBuch+
% Programme officiel MINESEC. Compiler avec : tectonic ENA20-reading-texts-on-managing-personal-finances.tex
\newcommand{\DOCMATIERE}{Anglais}
\newcommand{\DOCNIVEAU}{1ère A}
\newcommand{\DOCMODULE}{Chapter 2 : Economic Life and Occupations}
\newcommand{\DOCLECON}{ENA20}
\newcommand{\DOCTITRE}{Lesson 20 — Reading: Texts on managing personal finances}
% OnBuch+ — préambule « cours signés » ENRICHI (compilé avec tectonic / XeLaTeX)
% Système visuel « Pop » d'OnBuch+ : fond crème (jamais blanc), bordures encre
% épaisses, ombres « dures » décalées, titres Archivo Black en capitales.
% Version MATHÉMATIQUES : graphes (pgfplots/tikz), boîtes pédagogiques
% (définition, propriété, méthode, attention, le savais-tu, exercice résolu,
% exercice, TP, bilan), page de garde et signature OnBuch+ ancrée en bas.
%
% Macros attendues dans main.tex AVANT \input{preamble} :
%   \DOCMATIERE  \DOCNIVEAU  \DOCMODULE  \DOCLECON (numéro)  \DOCTITRE
\documentclass[11pt]{article}

\usepackage{fontspec}
\usepackage[english,french]{babel} % français = langue principale ; l'anglais via \eng{..} / textebox
\usepackage[a4paper,margin=2cm,top=3.4cm,bottom=2.8cm,headheight=1.4cm,footskip=1.3cm]{geometry}
\usepackage{amsmath,amssymb,amsfonts}
\usepackage{mathtools} % \xmapsto, \coloneqq... souvent utilisés par les modèles
\usepackage{cancel} % simplifications barrées (bilans de forces)
% \abs{z} (module, valeur absolue) et \norm{u} (norme), utilisés par les modèles
\providecommand{\abs}{}\let\abs\relax\DeclarePairedDelimiter{\abs}{\lvert}{\rvert}
\providecommand{\norm}{}\let\norm\relax\DeclarePairedDelimiter{\norm}{\lVert}{\rVert}
\usepackage[table]{xcolor} % [table] : fournit aussi \rowcolors (alternance de lignes)
\usepackage{pagecolor}
\usepackage{tikz}
\usetikzlibrary{arrows.meta,positioning,calc,shapes.geometric,shapes.misc,shapes.arrows,decorations.pathmorphing,decorations.pathreplacing,decorations.markings,patterns,shadows,mindmap,trees,fit,backgrounds,matrix,intersections,angles,quotes}
\usepackage{pgfplots}
\usepackage[version=4]{mhchem} % équations nucléaires \ce{^{235}_{92}U}
\usepackage[european,siunitx]{circuitikz} % schémas électriques
\pgfplotsset{compat=1.18}
\usepgfplotslibrary{fillbetween,groupplots} % aires entre courbes (calcul intégral)
\usepackage{enumitem}
\usepackage{fancyhdr}
% [most] charge toutes les bibliothèques tcolorbox (listings, minted, raster,
% documentation...) et gonfle inutilement la mémoire TeX sur un document long
% (~300 boîtes) : on ne charge que ce qui est réellement utilisé.
\usepackage[skins,breakable]{tcolorbox}
\usepackage{graphicx}
\usepackage{colortbl}
\usepackage{booktabs}
\usepackage{tabularx}
\usepackage{diagbox} % case d'en-tête diagonale des tableaux à double entrée
% Colonnes extensibles centrée / gauche / droite (C, L, R), souvent utilisées par les modèles
\newcolumntype{C}{>{\centering\arraybackslash}X}
\newcolumntype{L}{>{\raggedright\arraybackslash}X}
\newcolumntype{R}{>{\raggedleft\arraybackslash}X}
\usepackage{array}
\usepackage{multicol}
\usepackage{multirow}
\usepackage{siunitx}
% parse-numbers=false : accepte aussi \SI{2,5 \times 10^{-3}}{...}
\sisetup{locale=FR,output-decimal-marker={,},group-minimum-digits=4,parse-numbers=false}
\DeclareSIUnit\ohm{\ensuremath{\symup{\Omega}}} % U+2126 absent de la police
\DeclareSIUnit\division{div} % réglages d'oscilloscope (ms/div, V/div)
\DeclareSIUnit\tr{tr} % tours (vitesse de rotation, ex. \SI{1500}{\tr\per\minute})
\DeclareSIUnit\molar{M} % concentration molaire (ex. \SI{3}{\molar}, \SI{1}{\micro\molar})
\DeclareSIUnit\an{an} % année (le modèle confond parfois avec \year, un compteur TeX) — ex. \SI{5}{\kilo\meter\per\an}
\DeclareSIUnit\FCFA{FCFA} % franc CFA (ex. \SI{500}{\FCFA\per\kilo\gram})
\DeclareSIUnit\degreeBrix{\textdegree Brix} % teneur en sucre d'un sirop/jus (réfractométrie)
\DeclareSIUnit\month{mois} % le modèle utilise parfois \month, un compteur TeX, comme unité de durée
\usepackage{qrcode}
% \degree : commande fréquemment utilisée par les modèles pour un angle ou une
% température en dehors de \SI{}{} (non fournie par les paquets chargés).
\providecommand{\degree}{^{\circ}}
% \qed (fin de démonstration), utilisé par les modèles sans amsthm
\providecommand{\qed}{\hfill$\square$}
% \barre : double barre des valeurs interdites dans un tableau de variations (tkz-tab)
\providecommand{\barre}{\ensuremath{\|}}
% environnement proof (amsthm non chargé) utilisé par les modèles
\ifdefined\proof\else\newenvironment{proof}[1][Démonstration]{\par\noindent\textit{#1.}\ }{\qed\par}\fi
\usepackage{lastpage}
\usepackage{hyperref}
\hypersetup{hidelinks}

% ============================================================
% Palette « Pop » OnBuch+ — mêmes hex que la classe P (plus_ui.dart)
% ============================================================
\definecolor{poporange}{HTML}{F59321}
\definecolor{popdark}{HTML}{E07A0C}
\definecolor{popcream}{HTML}{FFF6E8}
\definecolor{popink}{HTML}{1C1714}
\definecolor{poppurple}{HTML}{7A5AE0}
\definecolor{popgreen}{HTML}{2E9E6A}
\definecolor{poppink}{HTML}{E0457B}
\definecolor{popblue}{HTML}{2F7DE1}
\definecolor{popgold}{HTML}{C98A12}
\definecolor{popmuted}{HTML}{8A7F76}
\definecolor{popline}{HTML}{EADFD2}
\definecolor{popgray}{HTML}{8A7F76}
\colorlet{popgrey}{popgray}
% Alias tolérants (le modèle nomme parfois les couleurs différemment)
\colorlet{popwhite}{white}
\colorlet{popblack}{popink}
\colorlet{popred}{poppink}
\colorlet{popyellow}{popgold}
% Teintes claires (fonds de tableaux/encadrés) — le modèle nomme parfois
% les couleurs avec un suffixe L (ex. popblueL) sans les avoir définies.
\colorlet{poporangeL}{poporange!20}
\colorlet{popdarkL}{popdark!20}
\colorlet{poppurpleL}{poppurple!20}
\colorlet{popgreenL}{popgreen!20}
\colorlet{poppinkL}{poppink!20}
\colorlet{popblueL}{popblue!20}
\colorlet{popgoldL}{popgold!20}
\colorlet{popredL}{popred!20}
\colorlet{popgrayL}{popgray!20}
% Teintes claires pour fonds de tableaux / zones
\colorlet{popblueL}{popblue!10!white}
\colorlet{poporangeL}{poporange!14!white}
\colorlet{popgreenL}{popgreen!12!white}
\colorlet{poppurpleL}{poppurple!12!white}
\colorlet{poppinkL}{poppink!12!white}
\colorlet{popgoldL}{popgold!14!white}

\pagecolor{popcream}

% ============================================================
% Typographie
% ============================================================
\defaultfontfeatures{Ligatures=TeX}
\setmainfont{PlusJakartaSans-Medium.ttf}[Path=../../../fonts/,BoldFont=PlusJakartaSans-Bold.ttf]
% Maths en sans-serif (Fira Math) avec chiffres et lettres droites de Plus
% Jakarta, pour que les formules se fondent dans le texte.
\usepackage[math-style=ISO,bold-style=ISO]{unicode-math}
\setmathfont{FiraMath-Regular.otf}[Path=../../../fonts/]
\setmathfont{PlusJakartaSans-Medium.ttf}[Path=../../../fonts/,range={up/num,up/{Latin,latin},\mathup}]
\usepackage{icomma} % virgule décimale sans espace en mode math
% Caractères mathématiques Unicode absents des polices de texte (Plus Jakarta,
% Archivo Black) : rendus par leur équivalent mathématique, en texte comme en
% formule — sinon ils s'affichent en carré vide (ex. « Arithmétique dans ℤ »).
\usepackage{newunicodechar}
\newunicodechar{ℝ}{\ensuremath{\mathbb{R}}}
\newunicodechar{ℤ}{\ensuremath{\mathbb{Z}}}
\newunicodechar{ℂ}{\ensuremath{\mathbb{C}}}
\newunicodechar{ℕ}{\ensuremath{\mathbb{N}}}
\newunicodechar{ℚ}{\ensuremath{\mathbb{Q}}}
\newunicodechar{𝔻}{\ensuremath{\mathbb{D}}}
\newunicodechar{α}{\ensuremath{\alpha}}
\newunicodechar{β}{\ensuremath{\beta}}
\newunicodechar{γ}{\ensuremath{\gamma}}
\newunicodechar{δ}{\ensuremath{\delta}}
\newunicodechar{ε}{\ensuremath{\varepsilon}}
\newunicodechar{θ}{\ensuremath{\theta}}
\newunicodechar{λ}{\ensuremath{\lambda}}
\newunicodechar{μ}{\ensuremath{\mu}}
\newunicodechar{σ}{\ensuremath{\sigma}}
\newunicodechar{τ}{\ensuremath{\tau}}
\newunicodechar{φ}{\ensuremath{\varphi}}
\newunicodechar{ψ}{\ensuremath{\psi}}
\newunicodechar{ω}{\ensuremath{\omega}}
\newunicodechar{Σ}{\ensuremath{\Sigma}}
% majuscules produites par \MakeUppercase dans les titres (α-aminés -> Α)
\newunicodechar{Α}{\ensuremath{\alpha}}
\newunicodechar{Β}{\ensuremath{\beta}}
\newunicodechar{Γ}{\ensuremath{\Gamma}}
\newunicodechar{Δ}{\ensuremath{\Delta}}
\newunicodechar{Ε}{\ensuremath{\varepsilon}}
\newunicodechar{Θ}{\ensuremath{\Theta}}
\newunicodechar{Λ}{\ensuremath{\Lambda}}
\newunicodechar{Μ}{\ensuremath{\mu}}
\newunicodechar{Τ}{\ensuremath{\tau}}
\newunicodechar{Φ}{\ensuremath{\Phi}}
\newunicodechar{Ψ}{\ensuremath{\Psi}}
\newunicodechar{Ω}{\ensuremath{\Omega}}
\newunicodechar{↦}{\ensuremath{\mapsto}}
\newunicodechar{∈}{\ensuremath{\in}}
\newunicodechar{∉}{\ensuremath{\notin}}
\newunicodechar{∧}{\ensuremath{\wedge}}
\newunicodechar{⇔}{\ensuremath{\Leftrightarrow}}
\newunicodechar{⟺}{\ensuremath{\Longleftrightarrow}}
\newunicodechar{⇒}{\ensuremath{\Rightarrow}}
\newunicodechar{⟹}{\ensuremath{\Longrightarrow}}
\newunicodechar{∘}{\ensuremath{\circ}}
\newunicodechar{∪}{\ensuremath{\cup}}
\newunicodechar{∩}{\ensuremath{\cap}}
\newunicodechar{≡}{\ensuremath{\equiv}}
\newunicodechar{⊂}{\ensuremath{\subset}}
\newunicodechar{⊥}{\ensuremath{\perp}}
\newunicodechar{∃}{\ensuremath{\exists}}
\newunicodechar{∀}{\ensuremath{\forall}}
\newunicodechar{∛}{\ensuremath{\sqrt[3]{\,}}}
\newunicodechar{∜}{\ensuremath{\sqrt[4]{\,}}}
\newunicodechar{⃗}{\ensuremath{{}^{\to}}}
\newunicodechar{ⁿ}{\ifmmode^{n}\else\textsuperscript{\ensuremath{n}}\fi}
\newunicodechar{ˣ}{\ifmmode^{x}\else\textsuperscript{\ensuremath{x}}\fi}
\newunicodechar{ᵇ}{\ifmmode^{b}\else\textsuperscript{\ensuremath{b}}\fi}
\newunicodechar{ʳ}{\ifmmode^{r}\else\textsuperscript{\ensuremath{r}}\fi}
\newunicodechar{ᵘ}{\ifmmode^{u}\else\textsuperscript{\ensuremath{u}}\fi}
\newunicodechar{ʸ}{\ifmmode^{y}\else\textsuperscript{\ensuremath{y}}\fi}
\newunicodechar{ᵃ}{\ifmmode^{a}\else\textsuperscript{\ensuremath{a}}\fi}
\newunicodechar{ᵉ}{\ifmmode^{e}\else\textsuperscript{\ensuremath{e}}\fi}
\newunicodechar{ᵏ}{\ifmmode^{k}\else\textsuperscript{\ensuremath{k}}\fi}
\newunicodechar{ᵖ}{\ifmmode^{p}\else\textsuperscript{\ensuremath{p}}\fi}
\newunicodechar{ᴺ}{\ifmmode^{N}\else\textsuperscript{\ensuremath{N}}\fi}
\newunicodechar{ᶜ}{\ifmmode^{c}\else\textsuperscript{\ensuremath{c}}\fi}
\newunicodechar{ʰ}{\ifmmode^{h}\else\textsuperscript{\ensuremath{h}}\fi}
\newunicodechar{ᵠ}{\ifmmode^{q}\else\textsuperscript{\ensuremath{q}}\fi}
\newunicodechar{ₙ}{\ifmmode_{n}\else\textsubscript{\ensuremath{n}}\fi}
\newunicodechar{ₐ}{\ifmmode_{a}\else\textsubscript{\ensuremath{a}}\fi}
\newunicodechar{ᵢ}{\ifmmode_{i}\else\textsubscript{\ensuremath{i}}\fi}
\newunicodechar{ᵣ}{\ifmmode_{r}\else\textsubscript{\ensuremath{r}}\fi}
\newunicodechar{ₖ}{\ifmmode_{k}\else\textsubscript{\ensuremath{k}}\fi}
\newunicodechar{ᵦ}{\ifmmode_{\beta}\else\textsubscript{\ensuremath{\beta}}\fi}
\newunicodechar{ₓ}{\ifmmode_{x}\else\textsubscript{\ensuremath{x}}\fi}
\newfontfamily\popdisplay{ArchivoBlack-Regular.ttf}[Path=../../../fonts/]
\newfontfamily\poplogo{SpaceGrotesk-Bold.ttf}[Path=../../../fonts/]
\newfontfamily\popsemi{PlusJakartaSans-SemiBold.ttf}[Path=../../../fonts/]
\newfontfamily\popxbold{PlusJakartaSans-ExtraBold.ttf}[Path=../../../fonts/]
\newfontfamily\popxboldls{PlusJakartaSans-ExtraBold.ttf}[Path=../../../fonts/,LetterSpace=3.0]

\setlist[enumerate]{leftmargin=1.7em,itemsep=5pt,topsep=4pt}
\setlist[itemize]{leftmargin=1.7em,itemsep=4pt,topsep=4pt,label={\color{poporange}\textbullet}}
\setlength{\parindent}{0pt}
\setlength{\parskip}{5pt}
\renewcommand{\arraystretch}{1.25}

% pgfplots : style Pop par défaut
\pgfplotsset{
  every axis/.append style={axis line style={popink,line width=1pt},tick style={popink},
    grid style={popline,line width=0.5pt},label style={font=\small},tick label style={font=\footnotesize}},
  popcurve/.style={poporange,line width=1.8pt,no markers,smooth},
}
% Même style disponible dans un \draw TikZ simple (hors pgfplots)
\tikzset{popcurve/.style={poporange,line width=1.8pt}}
\tikzset{popgrid/.style={popline,very thin}}
\colorlet{popcurve}{poporange} % le nom du style est parfois utilisé comme couleur

% ============================================================
% Pill / squiggle
% ============================================================
\newcommand{\poppill}[3]{%
  \tikz[baseline=-0.55ex]\node[draw=popink,line width=1.2pt,rounded corners=2.6pt,%
    fill=#2,inner xsep=6.5pt,inner ysep=3pt]{%
    \popxboldls\fontsize{7}{7}\selectfont\color{#3}\MakeUppercase{#1}};%
}
\newcommand{\popsquiggle}[1]{%
  \tikz[baseline=-0.4ex]{
    \draw[#1,line width=2.6pt,line cap=round]
      plot [smooth] coordinates {(0,0.09) (0.34,0.01) (0.68,0.21) (1.02,0.01) (1.36,0.10)};
  }%
}
% Mot-clé surligné (marqueur orange sous le texte)
\newcommand{\cle}[1]{{\popxbold\color{popdark}#1}}

% ============================================================
% En-tête / pied
% ============================================================
\pagestyle{fancy}
\fancyhf{}
\renewcommand{\headrulewidth}{0pt}
\renewcommand{\footrulewidth}{0pt}
\lhead{\raisebox{-0.15ex}{\poplogo\fontsize{13}{13}\selectfont\color{popink}OnBuch\color{poporange}+}}
\rhead{\poppill{\DOCMATIERE\ \textbullet\ \DOCNIVEAU\ \textbullet\ Leçon \DOCLECON}{white}{popink}}
\lfoot{\popxbold\fontsize{7.6}{7.6}\selectfont\color{popmuted}\MakeUppercase{Cours signé OnBuch+}\ \textbullet\ {\popsemi\DOCTITRE}}
\rfoot{\tikz[baseline=-0.6ex]\node[rounded corners=8.5pt,fill=popink,minimum height=17pt,minimum width=17pt,inner xsep=5pt,inner sep=0]{\popxbold\fontsize{8}{8}\selectfont\color{popcream}\thepage\,/\,\pageref*{LastPage}};}

% ============================================================
% Titres
% ============================================================
\newcommand{\chaptitre}[2]{%
  \begin{tcolorbox}[enhanced,colback=poporange,colframe=popink,boxrule=2.6pt,arc=11pt,
      drop shadow=popink,left=15pt,right=15pt,top=12pt,bottom=11pt,before skip=3pt,after skip=18pt]
    \poppill{#2}{popink}{popcream}\par\vspace{9pt}
    {\raggedright\hyphenpenalty=10000\exhyphenpenalty=10000\popdisplay\fontsize{22}{24}\selectfont\color{popcream}\MakeUppercase{#1}\par}
  \end{tcolorbox}
}
% Section numérotée : pastille orange avec le numéro + titre Archivo
\newcounter{popsec}
\newcommand{\coursec}[1]{%
  \stepcounter{popsec}\setcounter{popsub}{0}%
  \par\vspace{16pt}\noindent
  \begin{minipage}{\linewidth}
  \tikz[baseline=-0.7ex]\node[draw=popink,line width=1.6pt,fill=poporange,rounded corners=4pt,minimum size=22pt,inner sep=0]{\popdisplay\fontsize{12}{12}\selectfont\color{popcream}\thepopsec};%
  \hspace{8pt}{\popdisplay\fontsize{15}{17}\selectfont\color{popink}\MakeUppercase{#1}}\par
  \vspace{4pt}\noindent{\color{poporange}\rule{36pt}{3.6pt}}
  \end{minipage}\par\nopagebreak\vspace{8pt}
}
\newcounter{popsub}
\newcommand{\courssub}[1]{%
  \stepcounter{popsub}%
  \par\vspace{9pt}\noindent{\popxbold\fontsize{11.5}{13}\selectfont\color{popink}{\color{poporange}\thepopsec.\thepopsub}\hspace{6pt}#1}\par\nopagebreak\vspace{4pt}
}

% ============================================================
% Boîtes pédagogiques — même recette : fond blanc, bordure épaisse colorée,
% ombre dure encre, badge plein. Titre optionnel [#1].
% ============================================================
\tcbset{popbase/.style={breakable,enhanced,colback=white,boxrule=2.3pt,arc=9pt,
  shadow={3pt}{-3pt}{0pt}{popink},left=14pt,right=14pt,top=11pt,bottom=11pt,before skip=11pt,after skip=15pt,
  fontupper=\popsemi\fontsize{10.6}{14.5}\selectfont\color{popink},
  fontlower=\popsemi\fontsize{10.6}{14.5}\selectfont\color{popink}}}
% Coche dessinée (le glyphe ✓ n'existe pas dans les polices du document)
\newcommand{\popcheck}[1]{\tikz[baseline=-0.5ex]\draw[#1,line width=1.6pt,line cap=round,line join=round](-0.13,0.0)--(-0.04,-0.09)--(0.13,0.1);}
\providecommand{\coche}{\popcheck{popgreen}}
\providecommand{\resultat}[1]{\par\smallskip\noindent\fcolorbox{popgreen}{popgreenL}{\parbox{\dimexpr\linewidth-2\fboxsep-2\fboxrule\relax}{\textbf{Résultat.} #1}}\par\smallskip}
\newcommand{\popboxhead}[3]{\poppill{#1}{#2}{white}\ifx\relax#3\relax\else\hspace{6pt}{\popxbold\fontsize{10.6}{12}\selectfont\color{popink}#3}\fi\par\vspace{7pt}}

\newtcolorbox{aretenir}[1][]{popbase,colframe=popgreen,before upper={\popboxhead{À retenir}{popgreen}{#1}}}
% Texte en anglais (document de lecture, dialogue, extrait) : cadre
% d'étude. Un titre optionnel (source, auteur) s'affiche dans l'en-tête.
\newtcolorbox{textebox}[1][]{popbase,colframe=popblue,before upper={\popboxhead{Text}{popblue}{#1}\begin{otherlanguage}{english}},after upper={\end{otherlanguage}}}
% Marques ✓ / ✗ (forme correcte / fautive) souvent utilisées par les modèles
\providecommand{\cross}{\textcolor{poppink}{\ensuremath{\times}}}
\providecommand{\xmark}{\cross}\providecommand{\crossmark}{\cross}\providecommand{\cmark}{\ensuremath{\checkmark}}
% Ligne de réponse / justification à compléter (inventée par les modèles)
\providecommand{\justification}[1]{\par\noindent\textit{Justification~:} \dotfill\par}
% \textipa (package tipa non chargé) : la police ne couvre pas l'API -> simple passage
\providecommand{\textipa}[1]{#1}
% Soulignement (ulem non chargé) : \uline, \ul
\providecommand{\uline}[1]{\underline{#1}}\providecommand{\ul}[1]{\underline{#1}}
% \glossaire{\item ...} inventée par les modèles : liste de glossaire
\providecommand{\glossaire}[1]{\begin{itemize}#1\end{itemize}}
% \alert (beamer) inventée par les modèles : texte mis en évidence
\providecommand{\alert}[1]{\textbf{\textcolor{poppink}{#1}}}
% variantes ASCII d'ordinaux écrites en macro
\providecommand{\iere}{\textsuperscript{re}}\providecommand{\ieme}{\textsuperscript{e}}\providecommand{\ere}{\textsuperscript{re}}\providecommand{\eme}{\textsuperscript{e}}\providecommand{\er}{\textsuperscript{er}}\providecommand{\ie}{\textsuperscript{e}}
% Ordinaux français écrits en macro par les modèles : 1\ière, 3\ième
\newcommand{\ière}{\textsuperscript{re}}\newcommand{\ième}{\textsuperscript{e}}
% \exemple (inventée par les modèles) : étiquette « Exemple : »
\providecommand{\exemple}{\textbf{Exemple~:}\ }
% \square utilisable aussi hors mode math (cases à cocher)
\AtBeginDocument{\let\squareMath\square\renewcommand{\square}{\ensuremath{\squareMath}}}
% Mot ou phrase en anglais à l'intérieur d'un paragraphe français
\newcommand{\eng}[1]{\ifmmode\text{\textit{#1}}\else\begingroup\selectlanguage{english}\itshape #1\endgroup\fi}
\let\esp\eng % alias toléré
% flèches d'intonation inventées par les modèles
\providecommand{\textsearrow}{}\renewcommand{\textsearrow}{\ensuremath{\searrow}}\providecommand{\textnearrow}{}\renewcommand{\textnearrow}{\ensuremath{\nearrow}}
% \fr{..} : mot français (inventé par les modèles)
\providecommand{\fr}[1]{\textit{#1}}
\newtcolorbox{exemplebox}[1][]{popbase,colframe=poporange,before upper={\popboxhead{Exemple}{poporange}{#1}}}
\newtcolorbox{definition}[1][]{popbase,colframe=popblue,before upper={\popboxhead{Définition}{popblue}{#1}}}
\newtcolorbox{propriete}[1][]{popbase,colframe=poppurple,before upper={\popboxhead{Propriété}{poppurple}{#1}}}
\newtcolorbox{methode}[1][]{popbase,colframe=popgold,before upper={\popboxhead{Méthode}{popgold}{#1}}}
\newtcolorbox{attention}[1][]{popbase,colframe=poppink,before upper={\popboxhead{Attention}{poppink}{#1}}}
\newtcolorbox{experience}[1][]{popbase,colframe=popblue,colback=popblueL,before upper={\popboxhead{Expérience}{popblue}{#1}}}
% « Le savais-tu ? » : carte violette pleine (contexte, culture, Cameroun)
\newtcolorbox{savaistu}[1][]{popbase,colback=poppurple,colframe=popink,
  fontupper=\popsemi\fontsize{10.4}{14.2}\selectfont\color{white},
  before upper={\poppill{Le savais-tu\,?}{white}{poppurple}\ifx\relax#1\relax\else\hspace{6pt}{\popxbold\color{white}#1}\fi\par\vspace{7pt}}}
% Exercice résolu : énoncé en haut, \tcblower puis la solution sur fond crème
\newcounter{popexo}\newcounter{popexr}
\newtcolorbox{exoresolu}[1][]{popbase,colframe=poporange,colbacklower=popcream,
  segmentation style={popink,line width=1.2pt,dashed},
  before upper={\stepcounter{popexr}\popboxhead{Exercice résolu \thepopexr}{poporange}{#1}},
  before lower={\poppill{Solution}{popink}{popcream}\par\vspace{6pt}}}
% Exercice (non résolu en place) — la correction est dans la partie Corrigés
\newtcolorbox{exercice}[1][]{popbase,colframe=popink,
  before upper={\stepcounter{popexo}\popboxhead{Exercice \thepopexo}{popink}{#1}}}
% Corrigé
\newtcolorbox{corrige}[1][]{popbase,colframe=popgreen,colback=popgreenL,
  before upper={\popboxhead{Corrigé}{popgreen}{#1}}}
% Objectifs / prérequis / bilan
\newtcolorbox{objectifs}{popbase,colframe=popink,colback=poporangeL,
  before upper={\popboxhead{Objectifs de la leçon}{popink}{}}}
\newtcolorbox{prerequis}{popbase,colframe=popink,colback=popblueL,
  before upper={\popboxhead{Prérequis}{popblue}{}}}
\newtcolorbox{bilan}{popbase,colframe=popink,colback=white,boxrule=2.8pt,
  before upper={\popboxhead{Fiche bilan}{poporange}{L'essentiel à connaître pour l'examen}}}
% Figure encadrée avec légende
\newtcolorbox{popfigure}[1][]{enhanced,breakable=false,colback=white,colframe=popline,boxrule=1.4pt,arc=8pt,
  left=10pt,right=10pt,top=10pt,bottom=8pt,before skip=10pt,after skip=12pt,halign=center,
  lower separated=false,halign lower=center,
  fontlower=\popsemi\fontsize{9}{11.5}\selectfont\color{popmuted},
  #1}
\newcommand{\legende}[1]{\par\vspace{6pt}{\popsemi\fontsize{9}{11.5}\selectfont\color{popmuted}#1}}

% ============================================================
% Signature OnBuch+
% ============================================================
\newcommand{\poplogoinline}[1]{{\poplogo\fontsize{#1}{#1}\selectfont\color{popink}OnBuch\color{poporange}+}}
\newcommand{\signatureonbuch}{%
  \begin{tcolorbox}[enhanced,colback=popink,colframe=popink,boxrule=2.2pt,arc=10pt,
      drop shadow=poporange,left=14pt,right=14pt,top=10pt,bottom=10pt,before skip=0pt,after skip=0pt]
    \tikz[baseline=-0.7ex]\node[circle,fill=popgreen,draw=popcream,line width=1.3pt,minimum size=22pt,inner sep=0]{\popcheck{white}};%
    \hspace{9pt}\begin{minipage}[c]{0.62\linewidth}
      {\popxbold\fontsize{12}{13}\selectfont\color{popcream}Signé\ \textbullet\ {\poplogo OnBuch\color{poporange}+}}\par\vspace{2pt}
      {\raggedright\popsemi\fontsize{8}{10.5}\selectfont\color{popline}\MakeUppercase{Équipe pédagogique OnBuch+}\\\MakeUppercase{Conforme au programme officiel MINESEC}\par}
    \end{minipage}\hfill
    {\poplogo\fontsize{22}{22}\selectfont\color{popcream}OnBuch\color{poporange}+}
  \end{tcolorbox}%
}

% ============================================================
% Page de garde
% #1 titre · #2 matière · #3 niveau · #4 module · #5 n° leçon · #6 durée ·
% #7 description / objectifs (texte ou itemize)
% ============================================================
\newcommand{\pagegarde}[7]{%
  \thispagestyle{empty}
  \newgeometry{a4paper,margin=1.7cm,top=1.6cm,bottom=1.4cm}%
  \begin{tikzpicture}[remember picture,overlay]
    \fill[poporange!18] ([xshift=-3.2cm,yshift=-3.6cm]current page.north east) circle (4.6cm);
    \fill[poppurple!14] ([xshift=2.4cm,yshift=7.6cm]current page.south west) circle (3.8cm);
  \end{tikzpicture}%
  {\poplogo\fontsize{30}{30}\selectfont\color{popink}OnBuch\color{poporange}+}\hfill
  \raisebox{0.6ex}{\poppill{Cours signé \textbullet\ Premium}{popink}{popcream}}\par
  \vspace{1pt}\popsquiggle{poppurple}\par
  \vspace{8pt}
  \begin{tcolorbox}[enhanced,colback=poporange,colframe=popink,boxrule=2.8pt,arc=13pt,
      drop shadow=popink,left=18pt,right=18pt,top=12pt,bottom=13pt,after skip=10pt]
    \poppill{#2 \textbullet\ #3}{popink}{popcream}\hspace{5pt}\poppill{#4}{white}{popink}\par\vspace{8pt}
    {\popdisplay\fontsize{14}{14}\selectfont\color{popink}LEÇON #5}\par\vspace{4pt}
    {\raggedright\hyphenpenalty=10000\exhyphenpenalty=10000\popdisplay\fontsize{25}{27}\selectfont\color{popcream}\MakeUppercase{#1}\par}
    \vspace{8pt}
    {\popxbold\fontsize{9.5}{9.5}\selectfont\color{popink}\MakeUppercase{Durée conseillée : #6}}
  \end{tcolorbox}
  \begin{tcolorbox}[enhanced,colback=white,colframe=popink,boxrule=2.2pt,arc=10pt,
      drop shadow=popink,left=16pt,right=16pt,top=10pt,bottom=10pt,after skip=8pt]
    \poppill{Dans cette leçon}{poppurple}{white}\par\vspace{5pt}
    \popsemi\fontsize{9.3}{12.6}\selectfont\color{popink}#7
  \end{tcolorbox}
  \noindent\begin{minipage}[t]{0.487\linewidth}
    \begin{tcolorbox}[enhanced,colback=popgreen,colframe=popink,boxrule=2.2pt,arc=10pt,
        drop shadow=popink,left=11pt,right=11pt,top=10pt,bottom=10pt,height=3.1cm,valign=center]
      \tcbox[colback=white,colframe=popink,boxrule=1.3pt,arc=5pt,boxsep=0pt,nobeforeafter,
        left=3pt,right=3pt,top=3pt,bottom=3pt,valign=center]{\qrcode[height=1.9cm]{https://play.google.com/store/apps/details?id=com.luvvix.onbuch&hl=fr}}%
      \hspace{8pt}\begin{minipage}[c]{0.5\linewidth}
        {\popdisplay\fontsize{11}{12.5}\selectfont\color{white}\MakeUppercase{Télécharge OnBuch}}\par\vspace{4pt}
        {\raggedright\popsemi\fontsize{8.4}{11}\selectfont\color{white}Gratuit \textbullet\ Google Play\\Tuteur IA, annales, concours, résultats\par}
      \end{minipage}
    \end{tcolorbox}
  \end{minipage}\hfill
  \begin{minipage}[t]{0.487\linewidth}
    \begin{tcolorbox}[enhanced,colback=poppurple,colframe=popink,boxrule=2.2pt,arc=10pt,
        drop shadow=popink,left=11pt,right=11pt,top=10pt,bottom=10pt,height=3.1cm,valign=center]
      \tikz[baseline=-0.7ex]\node[circle,fill=white,draw=popink,line width=1.3pt,minimum size=1.6cm,inner sep=0]{\popdisplay\fontsize{24}{24}\selectfont\color{poppurple}+};%
      \hspace{8pt}\begin{minipage}[c]{0.6\linewidth}
        {\popdisplay\fontsize{11}{12.5}\selectfont\color{white}\MakeUppercase{Passe à OnBuch+}}\par\vspace{4pt}
        {\raggedright\popsemi\fontsize{8.4}{11}\selectfont\color{white}Tous les cours signés, Tuteur IA illimité, fiches PDF — onglet OnBuch+ de l'app\par}
      \end{minipage}
    \end{tcolorbox}
  \end{minipage}
  \vspace{8pt}
  \popcontenu
  \vfill
  \signatureonbuch
  \restoregeometry
  \newpage
}

% Rangée « Ce cours contient » (page de garde)
\newcommand{\poptile}[3]{% #1 couleur #2 gros texte #3 légende
  \begin{tcolorbox}[enhanced,colback=#1,colframe=popink,boxrule=2pt,arc=9pt,shadow={2.5pt}{-2.5pt}{0pt}{popink},
      left=6pt,right=6pt,top=9pt,bottom=9pt,halign=center,valign=center,height=1.75cm,width=\linewidth,nobeforeafter]
    {\popdisplay\fontsize{12.5}{13}\selectfont\color{white}\MakeUppercase{#2}}\par\vspace{3pt}
    {\centering\popsemi\fontsize{7.8}{9.5}\selectfont\color{white}#3\par}
  \end{tcolorbox}}
\newcommand{\popcontenu}{%
  \par\noindent{\popxboldls\fontsize{7.5}{7.5}\selectfont\color{popmuted}CE COURS SIGNÉ CONTIENT}\par\vspace{7pt}
  \noindent\begin{minipage}[t]{0.235\linewidth}\poptile{popblue}{Cours}{complet, illustré, pas à pas}\end{minipage}\hfill
  \begin{minipage}[t]{0.235\linewidth}\poptile{poporange}{Méthodes}{et exercices résolus}\end{minipage}\hfill
  \begin{minipage}[t]{0.235\linewidth}\poptile{poppink}{Exercices}{type examen + corrigés}\end{minipage}\hfill
  \begin{minipage}[t]{0.235\linewidth}\poptile{popgreen}{Fiche bilan}{l'essentiel à retenir}\end{minipage}\par}

% Sommaire
\newtcolorbox{sommaire}{enhanced,colback=white,colframe=popink,boxrule=2.2pt,arc=10pt,
  drop shadow=popink,left=18pt,right=18pt,top=13pt,bottom=13pt,before skip=6pt,after skip=20pt,
  before upper={\poppill{Sommaire}{poppurple}{white}\par\vspace{9pt}\popsemi\fontsize{10.6}{14.5}\selectfont\color{popink}}}

% Signature de fin de cours — poussée tout en bas de la dernière page
\newcommand{\findecours}{%
  \par\vfill\nopagebreak
  \begin{center}\popsquiggle{poporange}\end{center}\vspace{4pt}
  \signatureonbuch
}
\AtBeginDocument{\def\llbracket{\mathopen{[\mkern-3.5mu[}}\def\rrbracket{\mathclose{]\mkern-3.5mu]}}} % intervalles d'entiers (glyphe ⟦ absent de la police)
% Glyphes absents de Fira Math : redéfinis à partir de glyphes disponibles
\AtBeginDocument{%
  \def\setminus{\mathbin{\backslash}}\let\smallsetminus\setminus
  \def\vdots{\mathord{\vbox{\baselineskip3.2pt\lineskiplimit0pt\kern4pt\hbox{.}\hbox{.}\hbox{.}}}}%
  \def\ddots{\mathinner{\mkern1mu\raise6.5pt\hbox{.}\mkern2mu\raise3.6pt\hbox{.}\mkern2mu\raise0.7pt\hbox{.}\mkern1mu}}%
  \def\triangle{\mathord{\scalebox{0.9}{$\symup{\Delta}$}}}%
  \def\nabla{\mathord{\raisebox{\depth}{\scalebox{1}[-1]{$\symup{\Delta}$}}}}%
  \def\checkmark{\popcheck{popdark}}%
  \def\star{\mathbin{\ast}}\def\bigstar{\mathord{\ast}}%
  \def\diamond{\mathbin{\scalebox{0.6}{$\Diamond$}}}%
  \def\bigcup{\mathop{\mathchoice{\raisebox{-0.3em}{\scalebox{1.7}{$\cup$}}}{\scalebox{1.25}{$\cup$}}{\cup}{\cup}}\displaylimits}%
  \def\bigcap{\mathop{\mathchoice{\raisebox{-0.3em}{\scalebox{1.7}{$\cap$}}}{\scalebox{1.25}{$\cap$}}{\cap}{\cap}}\displaylimits}%
  \def\lesseqqgtr{\mathrel{\lessgtr}}\def\bigoplus{\mathop{\oplus}\displaylimits}%
}
\providecommand{\ssi}{si et seulement si} % abréviation fréquente
\AtBeginDocument{\providecommand{\wideparen}{\overparen}} % arc de cercle

\begin{document}

\pagegarde{Lesson 20 — Reading: Texts on managing personal finances}{Anglais}{1ère A}{Chapter 2 : Economic Life and Occupations}{ENA20}{2 h}{Cette leçon de compréhension écrite propose un article original sur la gestion des finances personnelles (budget, épargne, dépenses). Elle entraîne les élèves aux stratégies de lecture (skimming, scanning, déduction du sens) et enrichit leur lexique économique.
\vspace{4pt}
\begin{multicols}{2}\small
\begin{itemize}
\item À la fin de cette leçon, je sais repérer l'idée générale d'un texte grâce au skimming.
\item À la fin de cette leçon, je sais localiser une information précise grâce au scanning.
\item À la fin de cette leçon, je sais déduire le sens d'un mot inconnu à partir du contexte.
\item À la fin de cette leçon, je sais utiliser le lexique de l'argent : income, expenses, budget, savings, debt.
\item À la fin de cette leçon, je sais répondre à des questions de compréhension (vrai/faux, QCM, réponses courtes).
\item À la fin de cette leçon, je sais distinguer les besoins (needs) des envies (wants) dans un texte.
\end{itemize}\end{multicols}}

\begin{sommaire}
\begin{multicols}{2}
\begin{enumerate}
\item Warm-up and pre-reading
\item Reading text with glossary
\item Comprehension activities
\item Thematic vocabulary
\item Reading method and worked example
\item Guided production: summary and opinion
\item Activité d'intégration
\item Méthodes et exercices résolus
\item Exercices
\item Corrigés des exercices
\item Fiche bilan
\end{enumerate}
\end{multicols}
\end{sommaire}

\begin{objectifs}
\begin{itemize}
\item À la fin de cette leçon, je sais repérer l'idée générale d'un texte grâce au skimming.
\item À la fin de cette leçon, je sais localiser une information précise grâce au scanning.
\item À la fin de cette leçon, je sais déduire le sens d'un mot inconnu à partir du contexte.
\item À la fin de cette leçon, je sais utiliser le lexique de l'argent : income, expenses, budget, savings, debt.
\item À la fin de cette leçon, je sais répondre à des questions de compréhension (vrai/faux, QCM, réponses courtes).
\item À la fin de cette leçon, je sais distinguer les besoins (needs) des envies (wants) dans un texte.
\item À la fin de cette leçon, je sais résumer un texte en quelques phrases avec mes propres mots.
\item À la fin de cette leçon, je sais donner mon opinion sur la gestion de l'argent en anglais.
\end{itemize}
\end{objectifs}

\begin{prerequis}
\begin{itemize}
\item Vocabulaire de base de l'argent vu en Seconde : money, buy, sell, price, pay
\item Le présent simple pour exprimer des habitudes et des vérités générales
\item Les modaux should et must pour donner des conseils
\item Les connecteurs simples : and, but, because, so
\item La lecture de textes courts avec questions de compréhension (Seconde)
\end{itemize}
\end{prerequis}

\newpage

\coursec{Warm-up and pre-reading}

\begin{textebox}[A true-to-life situation — “Where did the money go?”]
Aminatou is a Form Five student at a secondary school in Buea. Last month, her uncle gave her 10,000 FCFA for her birthday. She was delighted. On Monday she bought a new phone case for 2,500 FCFA. On Wednesday she treated her friends to puff-puff and soft drinks at the school gate — 1,500 FCFA. On Saturday she paid for a bendskin ride, bought data for her phone and a pair of earrings at the market. Two weeks later, the money was gone, and Aminatou could not say exactly where it went. She had not planned anything. She had never made a budget.

Aminatou is not alone. Many teenagers in Cameroon, Nigeria and the United Kingdom receive pocket money every week or every month, but very few of them know how to manage it. They spend first and think later. Yet managing money is a skill — and like every skill, it can be learned. Financial experts give simple advice: know your income (the money that comes in), know your expenses (the money that goes out), make a budget (a plan for your money), and try to keep some savings (money you do not spend). If you spend more than you earn, you fall into debt — you owe money to somebody.

Today, we are going to read a text about managing personal finances. Before we read, we will activate what we already know, learn five key words, and discover three powerful reading strategies: skimming, scanning and guessing meaning from context.
\end{textebox}

\textbf{Glossaire des mots difficiles :}

\begin{center}
\begin{tabularx}{\linewidth}{|l|l|X|}
\hline
\rowcolor{popblueL} \textbf{Mot anglais} & \textbf{Nature} & \textbf{Traduction française} \\
\hline
pocket money & nom & l'argent de poche \\
\hline
to treat (somebody) & verbe & offrir (à quelqu'un), régaler \\
\hline
to manage & verbe & gérer \\
\hline
to owe & verbe & devoir (de l'argent à quelqu'un) \\
\hline
advice & nom (indénombrable) & des conseils \\
\hline
skill & nom & compétence, savoir-faire \\
\hline
\end{tabularx}
\end{center}

\emph{Problématique de la leçon :} \eng{How can we read an English text about money efficiently, and how can we talk about our own finances in English?} — Comment lire efficacement un texte anglais sur l'argent, et comment parler de nos propres finances en anglais ?

\courssub{Opening discussion: let's talk about money!}

Avant de lire, parlons ! Le professeur pose deux questions à la classe entière. Réponds en anglais, avec des phrases complètes. Voici des modèles pour t'aider :

\begin{itemize}
\item \eng{How much pocket money do you get?} → \eng{I get 500 FCFA a week from my mother.} / \eng{I don't get regular pocket money, but my father gives me money for transport.} (« Je reçois 500 FCFA par semaine de ma mère. » / « Je ne reçois pas d'argent de poche régulier, mais mon père me donne de l'argent pour le transport. »)
\item \eng{What do you spend it on?} → \eng{I spend it on food, phone credit and transport.} / \eng{I save part of it to buy books.} (« Je le dépense en nourriture, crédit téléphonique et transport. » / « J'en épargne une partie pour acheter des livres. »)
\end{itemize}

\begin{attention}[Les structures utiles : spend, get, save]
Note bien les constructions :
\begin{itemize}
\item \eng{to get money from somebody} : \eng{I get money \textbf{from} my aunt.} (préposition \cle{from}, pas « of »).
\item \eng{to spend money \textbf{on} something} : \eng{She spends money \textbf{on} clothes.} — jamais « spend money for clothes ».
\item \eng{to spend time/money \textbf{doing} something} : \eng{He spent 2,000 FCFA \textbf{buying} data.} — le verbe suivant est en \cle{-ing}.
\item \eng{to save money \textbf{for} something} : \eng{I am saving money \textbf{for} a bicycle.}
\end{itemize}
Erreur typique du francophone : dire « I spend my money \emph{in} food » par calque de « dépenser \emph{en} nourriture ». En anglais : \eng{spend money \textbf{on} food}.
\end{attention}

\begin{experience}[Class survey — two minutes]
En binômes, pose les deux questions à ton voisin et note sa réponse. Puis le professeur fait un sondage rapide au tableau : \eng{Who spends most of their money on food? On phone credit? Who saves money?} Utilise : \eng{Most students spend their money on...} / \eng{Only two students save money.} (« La plupart des élèves dépensent leur argent en... » / « Seulement deux élèves épargnent. »)
\end{experience}

\courssub{Brainstorming: words we already know}

Fermons les yeux une seconde : quels mots anglais liés à l'argent connais-tu déjà ? En classe, le professeur écrit \cle{MONEY} au centre du tableau et la classe propose des mots. Voici le résultat probable, organisé en carte mentale :

\begin{popfigure}
\begin{center}
\begin{tikzpicture}[mindmap, grow cyclic, every node/.style={concept, text=white}, root concept/.append style={concept color=popblue, font=\Large\bfseries}, level 1/.append style={level distance=3.4cm, sibling angle=90}, level 2/.append style={level distance=2.2cm, sibling angle=45}]
\node[root concept]{MONEY}
  child[concept color=popgreen]{ node {income} 
      child { node {salary} }
      child { node {pocket money} }
  }
  child[concept color=poporange]{ node {expenses}
      child { node {food} }
      child { node {transport} }
  }
  child[concept color=poppurple]{ node {savings}
      child { node {bank} }
      child { node {tontine} }
  }
  child[concept color=poppink]{ node {debt}
      child { node {to borrow} }
      child { node {to owe} }
  };
\end{tikzpicture}
\end{center}
\legende{Carte mentale du vocabulaire de l'argent : quatre grandes familles de mots.}
\end{popfigure}

\begin{center}
\begin{tabularx}{\linewidth}{|X|X|X|}
\hline
\rowcolor{popblueL} \textbf{English} & \textbf{Français} & \textbf{Exemple anglais} \\
\hline
money (indénombrable) & l'argent & \eng{Money is important, but it is not everything.} \\
\hline
price & le prix & \eng{What is the price of this bag?} \\
\hline
bank & la banque & \eng{My father keeps his money in a bank.} \\
\hline
market & le marché & \eng{She buys vegetables at the market in Douala.} \\
\hline
to buy / to sell & acheter / vendre & \eng{They sell cocoa and buy rice.} \\
\hline
cheap / expensive & bon marché / cher & \eng{Bendskins are cheap; taxis are expensive.} \\
\hline
\end{tabularx}
\end{center}

\begin{attention}[Faux amis et pièges autour de « money »]
\begin{itemize}
\item \eng{money} est \textbf{indénombrable} : on dit \eng{money \textbf{is}} (jamais « moneys » ni « a money »). Pour compter : \eng{a coin} (une pièce), \eng{a note / a bill} (un billet).
\item \eng{expensive} ≠ « expensif » : le mot français correct est « cher ». Et attention : un prix n'est jamais « expensive » ! On dit \eng{The price is \textbf{high}} / \eng{The bag is \textbf{expensive}}. (« Le prix est élevé » / « Le sac est cher ».)
\item \eng{to earn money} = gagner de l'argent (travail) ; \eng{to win money} = gagner de l'argent (au jeu, à la loterie). Ne confonds pas : \eng{She \textbf{earns} 50,000 FCFA a month.} / \eng{He \textbf{won} money in a lottery.}
\end{itemize}
\end{attention}

\courssub{Predicting: what will the text be about?}

Bonne nouvelle : un bon lecteur commence à lire \emph{avant} de lire. Le titre du texte que nous allons étudier est : \eng{“Managing Personal Finances: A Guide for Young People”} (« Gérer ses finances personnelles : un guide pour les jeunes »). À partir de ce titre seul, que peux-tu prédire ?

\begin{methode}[How to predict the content of a text — trois étapes]
\begin{enumerate}
\item \textbf{Read the title.} Souligne les mots importants. Ici : \eng{managing}, \eng{personal finances}, \eng{young people}. Demande-toi : \eng{What do I already know about this topic?}
\item \textbf{Look at the images.} Une photo d'un carnet, d'une tirelire ou d'un marché te donne des indices sur le contenu et le ton (sérieux ? humoristique ? pratique ?).
\item \textbf{Read the first sentence of each paragraph.} En anglais, la première phrase d'un paragraphe annonce souvent son idée principale (la \emph{topic sentence}). En lisant seulement ces phrases, tu obtiens le plan du texte.
\end{enumerate}
Ensuite, formule une prédiction complète : \eng{I think the text will explain how teenagers can plan their money and avoid debt.} (« Je pense que le texte expliquera comment les adolescents peuvent planifier leur argent et éviter les dettes. »)
\end{methode}

\begin{exoresolu}[Predicting from a title]
Énoncé : le titre d'un article est \eng{“Why Every Student Needs a Budget”}. Sans lire l'article, écris deux prédictions sur son contenu, en anglais.
\tcblower
Solution détaillée : j'applique la méthode en trois étapes.
\begin{enumerate}
\item Mots-clés du titre : \eng{student}, \eng{needs}, \eng{budget}. Le mot \eng{why} annonce des arguments, des raisons.
\item Pas d'image ici, mais le mot \eng{budget} évoque un plan, des chiffres, de la discipline.
\item Prédictions possibles : \eng{I think the text will give reasons why students should plan their money.} / \eng{It will probably explain how to make a simple budget and give examples of student expenses.} (« Je pense que le texte donnera des raisons pour lesquelles les élèves devraient planifier leur argent. » / « Il expliquera probablement comment faire un budget simple et donnera des exemples de dépenses d'élèves. »)
\end{enumerate}
Remarque la formulation : on utilise \eng{I think...}, \eng{probably}, \eng{will probably...} pour exprimer une hypothèse, pas une certitude.
\end{exoresolu}

\courssub{Three reading strategies: skimming, scanning, guessing from context}

Observer d'abord : quand tu cherches le numéro de téléphone d'un ami dans ton répertoire, tu ne lis pas tous les noms un par un — tes yeux « sautent » directement à la lettre voulue. Quand tu feuillettes un journal, tu lis les titres pour savoir de quoi parlent les articles. Tu pratiques déjà deux stratégies de lecture sans le savoir ! En anglais, elles ont des noms précis.

\begin{definition}[Skimming — lire pour l'idée générale]
\cle{Skimming} (littéralement « écrémer ») : lire un texte \textbf{rapidement}, sans s'arrêter sur les mots inconnus, pour saisir l'idée générale — \eng{the gist}. On lit le titre, la première et la dernière phrase, et la première phrase de chaque paragraphe. Question typique : \eng{What is the text about?} (« De quoi parle le texte ? »)
\end{definition}

\begin{definition}[Scanning — lire pour une information précise]
\cle{Scanning} : parcourir le texte \textbf{rapidement} pour trouver une information \textbf{précise} : un chiffre, un nom, une date, un lieu. On ne lit pas tout : les yeux cherchent un mot-repère (une majuscule, un nombre, un symbole). Question typique : \eng{How much money did Aminatou receive?} — réponse : \eng{10,000 FCFA}.
\end{definition}

\begin{definition}[Guessing meaning from context — déduire le sens des mots]
Quand tu rencontres un mot inconnu, \textbf{ne cours pas au dictionnaire}. Utilise le contexte : les mots autour, la structure de la phrase, ta logique. Exemple : \eng{Aminatou was \textbf{delighted} when she received the money.} Même sans connaître \eng{delighted}, le contexte (un cadeau d'argent !) indique une émotion positive : « ravie, très contente ».
\end{definition}

\begin{center}
\begin{tabularx}{\linewidth}{|X|X|X|}
\hline
\rowcolor{popblueL} \textbf{Stratégie} & \textbf{Objectif} & \textbf{Question typique} \\
\hline
Skimming & l'idée générale (the gist) & \eng{What is the text about?} \\
\hline
Scanning & un détail précis (chiffre, nom, date) & \eng{When? How much? Who?} \\
\hline
Guessing from context & le sens d'un mot inconnu & \eng{What does “debt” mean?} \\
\hline
\end{tabularx}
\end{center}

\courssub{Key words before reading: five essential words}

Pour lire notre texte sans difficulté, apprenons dès maintenant cinq mots-clés. Observe d'abord ces phrases tirées du texte que nous lirons :

\begin{itemize}
\item \eng{Your \textbf{income} is all the money you receive: pocket money, gifts, or a salary.}
\item \eng{Your \textbf{expenses} are all the things you pay for: food, transport, phone credit.}
\item \eng{A \textbf{budget} is a simple plan: it shows your income and your expenses.}
\item \eng{Try to keep some \textbf{savings} for emergencies — even 100 FCFA a week helps.}
\item \eng{If you spend more than your income, you go into \textbf{debt}.}
\end{itemize}

\begin{definition}[The five key words of personal finance]
\begin{itemize}
\item \cle{income} (nom, souvent indénombrable) : les revenus, l'argent qui \emph{entre}. Prononciation : « in-keum ».
\item \cle{expenses} (nom, presque toujours au pluriel) : les dépenses, l'argent qui \emph{sort}. Prononciation : « èks-pén-siz ».
\item \cle{budget} (nom) : le budget, un plan pour son argent. Prononciation : « beu-djèt ».
\item \cle{savings} (nom, toujours au pluriel) : l'épargne, l'argent mis de côté. Prononciation : « sé-vin-gz ».
\item \cle{debt} (nom) : la dette. Attention à l'orthographe : le \textbf{b} ne se prononce pas ! Prononciation : « dèt ».
\end{itemize}
\end{definition}

\begin{popfigure}
\begin{center}
\begin{tikzpicture}[node distance=0.6cm, every node/.style={font=\small}]
\node[draw=popgreen, fill=popgreenL, rounded corners, minimum width=2.6cm, minimum height=1cm, align=center] (inc) {\textbf{Income}\\ \emph{revenus}};
\node[draw=popblue, fill=popblueL, rounded corners, minimum width=2.6cm, minimum height=1cm, align=center, right=2.2cm of inc] (bud) {\textbf{Budget}\\ \emph{le plan}};
\node[draw=poporange, fill=poporangeL, rounded corners, minimum width=2.8cm, minimum height=1cm, align=center, right=2.2cm of bud, yshift=1.1cm] (exp) {\textbf{Expenses}\\ \emph{dépenses}};
\node[draw=poppurple, fill=poppurpleL, rounded corners, minimum width=2.8cm, minimum height=1cm, align=center, right=2.2cm of bud, yshift=-1.1cm] (sav) {\textbf{Savings}\\ \emph{épargne}};
\draw[-{Stealth[length=3mm]}, very thick, popgreen] (inc) -- (bud);
\draw[-{Stealth[length=3mm]}, very thick, popblue] (bud.east) -- (exp.west);
\draw[-{Stealth[length=3mm]}, very thick, popblue] (bud.east) -- (sav.west);
\node[below=0.9cm of bud, align=center, text=poppink, font=\small\itshape] {No savings left? You go into \textbf{debt}!\\ (Pas d'épargne ? Tu t'endettes !)};
\end{tikzpicture}
\end{center}
\legende{Le circuit de l'argent : les revenus alimentent le budget, qui se répartit entre dépenses et épargne.}
\end{popfigure}

\begin{attention}[Pièges de francophones sur ces cinq mots]
\begin{itemize}
\item \eng{income} ≠ « une rentrée » au sens scolaire ; c'est uniquement l'argent qui entre. Et ne dis pas « my incomes are small » : préfère \eng{My income is small.} (indénombrable).
\item \eng{expenses} : presque toujours au pluriel — \eng{My expenses are too high.} Le singulier \eng{an expense} existe mais est rare.
\item \eng{debt} : le \textbf{b} est muet (« dèt »), comme dans \eng{doubt} (« daoute »). Ne prononce jamais le b !
\item \eng{savings} : toujours avec un \textbf{s}, même pour une petite somme : \eng{My savings are in a tontine.}
\item Faux ami : \eng{to save money} = épargner ; « sauver » (une vie) = \eng{to save a life}. Même verbe, deux sens selon le complément.
\end{itemize}
\end{attention}

\begin{exoresolu}[Matching words and definitions]
Énoncé : associe chaque mot à sa définition anglaise. Mots : \eng{income}, \eng{expenses}, \eng{budget}, \eng{savings}, \eng{debt}. Définitions : (a) \eng{money you keep and do not spend} ; (b) \eng{money you owe to somebody} ; (c) \eng{all the money you receive} ; (d) \eng{a plan for your money} ; (e) \eng{all the money you spend}.
\tcblower
Solution détaillée :
\begin{itemize}
\item \eng{income} → (c) : l'argent reçu (pocket money, salary).
\item \eng{expenses} → (e) : l'argent dépensé (food, transport).
\item \eng{budget} → (d) : le plan qui organise revenus et dépenses.
\item \eng{savings} → (a) : l'argent gardé, non dépensé.
\item \eng{debt} → (b) : l'argent dû à quelqu'un.
\end{itemize}
Vérification par une phrase complète : \eng{If my expenses are higher than my income, my budget fails, I have no savings, and I go into debt.} (« Si mes dépenses dépassent mes revenus, mon budget échoue, je n'ai pas d'épargne et je m'endette. »)
\end{exoresolu}

\begin{savaistu}[Financial education in the United Kingdom]
Au Royaume-Uni, la gestion du budget fait partie du programme scolaire : les élèves suivent des cours de \eng{financial education} (« éducation financière ») où ils apprennent à ouvrir un compte en banque, à comparer des prix, à éviter les dettes et même à comprendre les impôts. Beaucoup d'écoliers britanniques reçoivent leur \eng{pocket money} chaque semaine et tiennent un petit carnet de dépenses. Et au Cameroun ? La \emph{tontine} — cette épargne rotative traditionnelle — est une forme très ancienne d'éducation financière pratiquée par les adultes... et de plus en plus par les élèves !
\end{savaistu}

\begin{aretenir}[What you must remember before reading]
\begin{itemize}
\item Avant de lire : discute du thème, active ton vocabulaire, \textbf{prédis} le contenu (titre + images + premières phrases).
\item Trois stratégies : \cle{skimming} (idée générale), \cle{scanning} (détail précis), \cle{guessing from context} (sens des mots inconnus).
\item Cinq mots-clés : \cle{income} (revenus), \cle{expenses} (dépenses), \cle{budget} (plan), \cle{savings} (épargne), \cle{debt} (dette — b muet).
\item Structures : \eng{spend money \textbf{on} something}, \eng{get money \textbf{from} somebody}, \eng{save money \textbf{for} something}.
\end{itemize}
Tu es maintenant prêt à lire le texte sur la gestion des finances personnelles !
\end{aretenir}

\coursec{Warm-up and pre-reading}

\begin{methode}[Avant la lecture : activer ses connaissances et prédire]
\begin{enumerate}
\item \textbf{Observation d'image.} Regarde l'illustration (un élève qui compte de l'argent, un carnet de budget, une tirelire). Décris-la en anglais : \eng{What can you see? Who is the person? What is he doing?}
\item \textbf{Brainstorming lexical.} En binôme, liste 8 à 10 mots anglais liés à l'argent de poche et à la gestion d'un budget (\eng{pocket money, savings, expenses, budget, income, need, want, debt, bank, to save}).
\item \textbf{Discussion orale (Speaking).} Réponds brièvement : \eng{How much pocket money do you get? Do you save money? What do you spend your money on?}
\item \textbf{Prédiction (Reading strategy).} Lis le titre \eng{« Smart Money for Young People »}. Quels conseils le texte va-t-il donner ? Écris 3 prédictions en anglais.
\end{enumerate}
\end{methode}

\begin{experience}[Activité de classe : « Money Habits Survey »]
En groupes de 4, interroge tes camarades avec ce questionnaire. Note les réponses en anglais.
\begin{enumerate}
\item \eng{Do you receive pocket money regularly?}
\item \eng{Do you have a savings account or a piggy bank?}
\item \eng{What is your biggest expense: food, data, clothes, transport?}
\item \eng{Do you know the 50/30/20 rule?}
\end{enumerate}
Présente les résultats de ton groupe à la classe : \eng{« In our group, 3 students save money, 1 student spends everything on data... »}
\end{experience}

\coursec{Reading text with glossary}

Après notre échauffement sur l'argent de poche (\eng{pocket money}), nous allons maintenant lire un texte complet sur la gestion des finances personnelles. Notre objectif n'est pas de comprendre chaque mot, mais de comprendre le sens général du texte (\eng{skimming}), de repérer des informations précises (\eng{scanning}) et de deviner le sens des mots inconnus grâce au contexte. C'est exactement la compétence qui sera évaluée à l'examen.

\courssub{The reading text: Smart Money for Young People}

\begin{methode}[Avant de lire : trois étapes]
\begin{enumerate}
\item Lis le \textbf{titre} et regarde l'illustration : de quoi va parler le texte ?
\item Lis le texte \textbf{une première fois rapidement} (\eng{skimming}) pour saisir l'idée générale : qui ? où ? quel problème ? quelle solution ?
\item Relis \textbf{lentement} (\eng{scanning}) en soulignant les chiffres et les mots-clés du thème de l'argent.
\end{enumerate}
\end{methode}

\begin{textebox}[Smart Money for Young People]
Tabi is a Form Five student at a secondary school in Bamenda. Every month, he receives 5,000 FCFA as pocket money from his uncle, but after only one week, he is always broke. He buys fried puff-puff and soft drinks every day at break time, he pays for mobile data to watch football videos, and he never knows where his money has gone.

One day, his school counsellor, Mrs Ndifor, calls him to her office. She explains a simple rule of budgeting: “Needs before wants, Tabi.” A need is something you cannot live without, like food, transport to school or exercise books. A want is something nice but not necessary, like new trainers or extra data. “If you pay for your wants first,” she says, “you will never have enough for your needs.”

Then she teaches him the famous 50/30/20 rule. It means that you should spend about 50\% of your money on needs, 30\% on wants, and you must save 20\%. For Tabi, that means 2,500 FCFA for needs, 1,500 FCFA for wants, and 1,000 FCFA for savings every month.

Tabi follows the advice. He opens a savings account at a microfinance bank near Commercial Avenue. He stops borrowing money from his friends, because unnecessary debts are dangerous: when you borrow, you must pay back, and then you have less money for yourself. After six months, Tabi has saved 6,000 FCFA, and he buys a second-hand bicycle to go to school.

Tabi smiles and tells his friends: “It is not how much you earn, but how you manage it.”
\end{textebox}

\courssub{Glossary: understanding difficult words}

Voici les dix mots difficiles du texte, expliqués en anglais simple puis traduits. Au brouillon, entraîne-toi à deviner leur sens \emph{avant} de regarder ce glossaire : c'est la stratégie de \cle{deduction from context}.

\begin{center}
\begin{tabularx}{\linewidth}{|l|X|X|}
\hline
\rowcolor{popblueL} \textbf{Word} & \textbf{Simple English meaning} & \textbf{Français} \\
\hline
pocket money & money parents or family give to a child regularly & argent de poche \\
\hline
broke (adj.) & having no money left & fauché, sans argent \\
\hline
counsellor & a person who gives advice at school & conseiller d'orientation \\
\hline
budget (n.) & a plan for your income and expenses & budget, plan de dépenses \\
\hline
need (n.) & something you must have to survive & un besoin (indispensable) \\
\hline
want (n.) & something you would like to have but can live without & une envie, un désir (superflu) \\
\hline
savings (n. pl.) & money you keep for the future & l'épargne, les économies \\
\hline
to save & to keep money instead of spending it & épargner, économiser \\
\hline
to borrow & to take money from someone with the promise to return it & emprunter \\
\hline
debt (n.) & money that you owe and must pay back & une dette \\
\hline
\end{tabularx}
\end{center}

\begin{definition}[Budget]
A \cle{budget} is a plan that shows how much money you receive (income) and how you plan to spend it (expenses). En français : un budget. Collocations utiles : \eng{to make a budget} (établir un budget), \eng{to stick to a budget} (respecter un budget), \eng{to balance a budget} (équilibrer un budget).
\end{definition}

\begin{definition}[Savings]
\cle{Savings} (pluriel) means money that you keep, usually in a bank, instead of spending it. En français : l'épargne (singulier en français, pluriel en anglais). Le verbe est \eng{to save}. Attention : \eng{a saving} (singulier) = une économie réalisée sur un achat précis (ex. \eng{I made a saving of 500 francs}).
\end{definition}

\begin{attention}[Pièges fréquents pour francophones : Broke / Break ; Borrow / Lend ; Spend on]
\begin{itemize}
\item \textbf{Broke (adj.)} $\neq$ \textbf{Broken (participe passé)}. \eng{He is broke} = Il est fauché. \eng{He has broken the window} = Il a cassé la fenêtre. Le participe passé de \eng{break} est \emph{broken}, jamais \emph{broke} dans le sens « casser ».
\item \textbf{Borrow} (emprunter = prendre) $\neq$ \textbf{Lend} (prêter = donner). Structure : \eng{borrow \textbf{from} sb} / \eng{lend \textbf{to} sb}. \textbf{Jamais} : \emph{Can you borrow me money?} $\rightarrow$ \eng{Can you \textbf{lend} me money?} ou \eng{Can I \textbf{borrow} money \textbf{from} you?}
\item \textbf{Spend} : \eng{spend money \textbf{on} sth} / \eng{spend time \textbf{doing} sth}. Préposition \textbf{on} (pas \emph{in}, pas \emph{for}). Verbe irrégulier : \emph{spend -- spent -- spent}.
\item \textbf{Pay} : \eng{pay \textbf{for} sth} (payer quelque chose), \eng{pay sb} (payer quelqu'un), \eng{pay \textbf{back} sb / sth} (rembourser).
\end{itemize}
\end{attention}

\courssub{Observing the language: key sentences from the text}

Observons deux phrases importantes du texte et leur traduction pour ancrer la grammaire en contexte.

\begin{exemplebox}[Exemples traduits et analysés]
\eng{He saves 1,000 FCFA every month.} « Il épargne 1 000 FCFA chaque mois. » — Remarque le \textbf{-s} de la 3\textsuperscript{e} personne du singulier au \cle{present simple} : \emph{he save\textbf{s}}. Marqueur de fréquence : \eng{every month}.

\eng{She spent all her money on clothes.} « Elle a dépensé tout son argent en vêtements. » — \eng{to spend} est un \cle{verbe irrégulier} : \emph{spend -- spent -- spent}. Structure : \eng{spent money \textbf{on} something}. Préposition \textbf{on} obligatoire.
\end{exemplebox}

\begin{propriete}[Les verbes de l'argent et leurs prépositions (Money verbs \& prepositions)]
\begin{center}
\begin{tabularx}{\linewidth}{|l|X|X|}
\hline
\rowcolor{popgreenL} \textbf{Structure} & \textbf{English example} & \textbf{Français} \\
\hline
to spend money \textbf{on} sth & He spends 500 FCFA on data. & Il dépense 500 F en données. \\
\hline
to save money \textbf{for} sth & She saves money for a bicycle. & Elle épargne pour un vélo. \\
\hline
to borrow money \textbf{from} sb & Tabi borrows money from a friend. & Tabi emprunte de l'argent à un ami. \\
\hline
to lend money \textbf{to} sb & The friend lends money to Tabi. & L'ami prête de l'argent à Tabi. \\
\hline
to pay \textbf{back} money / sb & He pays back 2,000 FCFA. & Il rembourse 2 000 F. \\
\hline
to pay \textbf{for} sth & I pay for my transport. & Je paie mon transport. \\
\hline
\end{tabularx}
\end{center}
\end{propriete}

\courssub{Tabi's budget: the 50/30/20 rule}

La règle 50/30/20 partage le revenu net en trois catégories. Voici le budget mensuel de Tabi (5 000 FCFA) appliqué à la règle.

\begin{popfigure}
\begin{tikzpicture}
% Pie chart: 50% (180deg), 30% (108deg), 20% (72deg)
\fill[popblue] (0,0) -- (90:2.2) arc (90:270:2.2) -- cycle;
\fill[poporange] (0,0) -- (270:2.2) arc (270:378:2.2) -- cycle;
\fill[popgreen] (0,0) -- (378:2.2) arc (378:450:2.2) -- cycle;
% Labels inside slices
\node[text=white, align=center, font=\small\bfseries] at (180:1.1) {Needs\\50\%};
\node[text=white, align=center, font=\small\bfseries] at (324:1.2) {Wants\\30\%};
\node[text=white, align=center, font=\small\bfseries] at (54:1.2) {Savings\\20\%};
% Legend
\node[align=left, font=\small, anchor=west] at (2.6, 1.2) {\textcolor{popblue}{\rule{0.4cm}{0.4cm}} Needs: 2,500 FCFA};
\node[align=left, font=\small, anchor=west] at (2.6, 0.4) {\textcolor{poporange}{\rule{0.4cm}{0.4cm}} Wants: 1,500 FCFA};
\node[align=left, font=\small, anchor=west] at (2.6, -0.4) {\textcolor{popgreen}{\rule{0.4cm}{0.4cm}} Savings: 1,000 FCFA};
\end{tikzpicture}
\legende{Tabi's monthly budget: the 50/30/20 rule applied to 5,000 FCFA.}
\end{popfigure}

\begin{savaistu}[Le saviez-vous ? : Culture financière au Cameroun]
La règle 50/30/20 a été popularisée par la sénatrice américaine Elizabeth Warren comme méthode simple pour les jeunes et les familles. Au Cameroun, notamment dans les régions anglophones (Northwest, Southwest), beaucoup d'élèves et de commerçants utilisent les \eng{njangi} ou \eng{tontines} (appelées \eng{rotating savings and credit associations} ou \eng{ROSCAs} en anglais) : chaque membre cotise une somme fixe chaque mois et reçoit la cagnotte totale à tour de rôle. C'est une forme d'épargne communautaire solidaire très ancrée dans la culture locale.
\end{savaistu}

\coursec{Comprehension activities}

\courssub{Worked example: reading comprehension (modèle de réponse)}

\begin{exoresolu}[Comprehension questions on « Smart Money for Young People »]
\textbf{Questions.}
\begin{enumerate}
\item Why is Tabi always broke after one week?
\item What is the difference between a need and a want according to Mrs Ndifor?
\item How much does Tabi save every month with the 50/30/20 rule?
\item Why does Mrs Ndifor say that debts are dangerous?
\item What does the last sentence of the text mean in your own words?
\end{enumerate}
\tcblower
\textbf{Model Answers.}
\begin{enumerate}
\item Because he spends his pocket money on puff-puff, soft drinks and mobile data without a plan, so his 5,000 FCFA disappear quickly.
\item A need is something you cannot live without (food, transport, exercise books); a want is something nice but not necessary (new trainers, extra data).
\item He saves 20\% of 5,000 FCFA, that is 1,000 FCFA every month.
\item Because when you borrow money, you must pay it back, and then you have less money for your own needs.
\item It means that the amount of money you receive is less important than the way you plan, spend and save it: even a small income is enough if you manage it well.
\end{enumerate}
\end{exoresolu}

\courssub{Practice: deduction from context (vocabulary strategy)}

\begin{exercice}[Your turn: deduce the meaning from context]
Sans dictionnaire, devine le sens des mots en \textbf{gras} grâce au contexte (mots avant/après, logique de la phrase), puis choisis la bonne traduction parmi la liste : \textit{fauché, rembourser, emprunter, épargne, conseiller}.
\begin{enumerate}
\item “After one week, he is always \textbf{broke}.”
\item “He opens a \textbf{savings} account at a microfinance bank.”
\item “His school \textbf{counsellor}, Mrs Ndifor, calls him to her office.”
\item “He stops \textbf{borrowing} money from his friends.”
\item “When you borrow, you must \textbf{pay back} the money.”
\end{enumerate}
\end{exercice}

\begin{corrige}[Exercice : Deduction from context]
\begin{enumerate}
\item \textbf{broke} $\rightarrow$ \textit{fauché} (contexte : « after one week », « money has gone »).
\item \textbf{savings} $\rightarrow$ \textit{épargne} (contexte : « account », « bank », « keep money »).
\item \textbf{counsellor} $\rightarrow$ \textit{conseiller} (contexte : « school », « gives advice », « calls him to her office »).
\item \textbf{borrowing} $\rightarrow$ \textit{emprunter} (contexte : « from his friends », « debts », « pay back »).
\item \textbf{pay back} $\rightarrow$ \textit{rembourser} (contexte : « borrow », « debts », « return money »).
\end{enumerate}
\emph{Remarquez la méthode : le contexte (mots voisins, connecteurs logiques, champ lexical) donne toujours des indices fiables.}
\end{corrige}

\coursec{Thematic vocabulary: Personal Finance}

\courssub{Word families and collocations (Familles de mots et collocations)}

Pour enrichir ton expression écrite et orale, apprends ces mots par familles (nom / verbe / adjectif) et leurs collocations naturelles (mots qui vont ensemble).

\begin{center}
\begin{tabularx}{\linewidth}{|l|l|l|X|}
\hline
\rowcolor{poppurpleL} \textbf{Noun (Nom)} & \textbf{Verb (Verbe)} & \textbf{Adjective (Adjectif)} & \textbf{Collocations utiles / Exemples} \\
\hline
\cle{income} (revenu) & -- & -- & \eng{monthly income, low income, source of income} \\
\hline
\cle{expense} (dépense) & \eng{to spend} & \eng{expensive} & \eng{monthly expenses, cover expenses, unexpected expense} \\
\hline
\cle{budget} (budget) & \eng{to budget} & \eng{budgetary} & \eng{to draw up a budget, tight budget, budget deficit} \\
\hline
\cle{savings} (épargne) & \eng{to save} & -- & \eng{savings account, life savings, to save up for sth} \\
\hline
\cle{debt} (dette) & \eng{to owe} & \eng{indebted} & \eng{to be in debt, to pay off a debt, national debt} \\
\hline
\cle{loan} (prêt bancaire) & \eng{to lend} & -- & \eng{to take out a loan, loan repayment, interest rate} \\
\hline
\cle{account} (compte) & -- & -- & \eng{current account, savings account, bank account} \\
\hline
\cle{interest} (intérêt) & -- & \eng{interesting} & \eng{interest rate, to earn interest, compound interest} \\
\hline
\end{tabularx}
\end{center}

\begin{attention}[Faux amis et pièges de traduction]
\begin{itemize}
\item \textbf{Revenue} (anglais) = \textit{revenu d'entreprise / recettes fiscales}. Pour « salaire / revenu personnel », utilise \eng{income} ou \eng{salary} / \eng{wages}.
\item \textbf{Benefit} = \textit{avantage, allocation (sociale)}. Pas « bénéfice » (qui est \eng{profit}).
\item \textbf{Charge} (verbe) = \textit{faire payer, charger}. « Frais / coûts » = \eng{charges} (nom pluriel) ou \eng{fees} / \eng{costs}.
\item \textbf{Investment} = \textit{investissement}. \eng{To invest in} (préposition \textbf{in}).
\end{itemize}
\end{attention}

\coursec{Reading method and worked example}

\courssub{Method: How to answer comprehension questions (Méthode pas à pas)}

\begin{methode}[Répondre aux questions de compréhension écrite]
\begin{enumerate}
\item \textbf{Lis la question} et identifie le mot-clé (\eng{Who, What, Where, When, Why, How, How much}).
\item \textbf{Scanne le texte} : cherche le mot-clé ou ses synonymes dans le paragraphe concerné.
\item \textbf{Lis la phrase entière} autour du mot trouvé (contexte).
\item \textbf{Formule la réponse en anglais}, en reprenant les mots du texte si possible, ou en reformulant avec tes mots (\eng{paraphrase}). Ne copie pas bêtement de longs passages.
\item \textbf{Vérifie} : la réponse répond-elle exactement à la question ? L'orthographe et la grammaire sont-elles correctes ?
\end{enumerate}
\end{methode}

\begin{exemplebox}[Application de la méthode : Question « Why »]
\textbf{Question :} \eng{Why does Tabi open a savings account?}
\textbf{Étape 1 :} Mot-clé = \eng{Why} (raison) + \eng{open a savings account}.
\textbf{Étape 2 :} Scan $\rightarrow$ Paragraphe 4 : \eng{« He opens a savings account... He stops borrowing money... because unnecessary debts are dangerous »}.
\textbf{Étape 3 :} Contexte : il veut arrêter d'emprunter / il veut épargner.
\textbf{Étape 4 :} \textbf{Answer :} \eng{He opens a savings account to stop borrowing money from his friends and to save money safely.}
\textbf{Étape 5 :} Vérification : répond au « why », anglais correct, pas de copier-coller intégral.
\end{exemplebox}

\coursec{Guided production: summary and opinion}

\courssub{Writing a summary (Résumé en anglais)}

\begin{methode}[Comment rédiger un résumé en 5 phrases maximum (Summary writing)]
\begin{enumerate}
\item Identifie les \textbf{4--5 idées principales} du texte (un par paragraphe).
\item Relie-les avec des \cle{connecteurs logiques} : \eng{First / Firstly, Then, After that, Finally, In the end}.
\item Utilise le \cle{Present Simple} (récit au présent) ou le \cle{Past Simple} (si tu racontes l'histoire comme un fait passé). Sois cohérent.
\item Évite les détails (chiffres exacts, noms propres secondaires) sauf s'ils sont cruciaux.
\item Limite-toi à \textbf{50--60 mots} pour un texte de cette longueur.
\end{enumerate}
\end{methode}

\begin{exemplebox}[Modèle de résumé (environ 55 mots)]
\eng{Tabi, a student in Bamenda, spends his pocket money carelessly and is broke after one week. His counsellor, Mrs Ndifor, teaches him the 50/30/20 budgeting rule: 50\% for needs, 30\% for wants, 20\% for savings. Tabi follows her advice, opens a savings account and stops borrowing. After six months, he buys a bicycle. The story shows that good management matters more than income.}
\end{exemplebox}

\courssub{Giving your opinion (Donner son avis)}

\begin{propriete}[Expressions pour donner son avis (Giving an opinion)]
\begin{center}
\begin{tabularx}{\linewidth}{|l|X|}
\hline
\rowcolor{poppinkL} \textbf{Structure} & \textbf{Exemple} \\
\hline
\eng{In my opinion,} + clause & \eng{In my opinion, the 50/30/20 rule is very useful for students.} \\
\hline
\eng{I think (that)} + clause & \eng{I think that saving money is difficult but necessary.} \\
\hline
\eng{I believe (that)} + clause & \eng{I believe njangi groups help people save together.} \\
\hline
\eng{From my point of view,} + clause & \eng{From my point of view, schools should teach financial literacy.} \\
\hline
\eng{It seems to me that} + clause & \eng{It seems to me that many young people confuse needs and wants.} \\
\hline
\end{tabularx}
\end{center}
\end{propriete}

\begin{exercice}[Guided production: Your summary and opinion]
\begin{enumerate}
\item \textbf{Summary.} Écris un résumé du texte \eng{« Smart Money for Young People »} en \textbf{50--60 mots maximum}, en suivant la méthode en 5 étapes. Utilise le \textbf{Present Simple}.
\item \textbf{Opinion.} Écris \textbf{3 phrases} en anglais pour donner ton avis personnel sur l'importance d'apprendre à gérer son budget quand on est lycéen. Utilise au moins deux structures de la boîte « Giving an opinion ».
\end{enumerate}
\end{exercice}

\begin{corrige}[Exercice : Guided production (Proposition de correction)]
\textbf{1. Summary (54 words) :}
\eng{Tabi, a Form Five student in Bamenda, wastes his 5,000 FCFA pocket money in one week. His school counsellor, Mrs Ndifor, introduces him to the 50/30/20 rule: spend 50\% on needs, 30\% on wants, save 20\%. Tabi applies the rule, opens a savings account and stops borrowing. Six months later, he buys a bicycle. The text proves that managing money well is more important than earning a lot.}

\textbf{2. Opinion (exemples) :}
\eng{In my opinion, learning to budget is essential for students because pocket money is limited. I think that the 50/30/20 rule is a simple tool to avoid debt. From my point of view, schools should include financial education in the curriculum.}
\end{corrige}

\begin{aretenir}[L'essentiel de la Lesson 20 — Managing Personal Finances]
\begin{itemize}
\item \textbf{Stratégies de lecture :} \eng{Skimming} (idée générale), \eng{Scanning} (info précise), \cle{Deduction from context} (deviner le sens).
\item \textbf{Vocabulaire clé :} \eng{pocket money, income, expenses, budget, need, want, savings, debt, loan, account, to borrow, to lend, to spend on, to save for, to pay back, to pay for}.
\item \textbf{Grammaire en contexte :} \eng{Present Simple} (habitudes, règles), \eng{Past Simple} (récit : \emph{spent, borrowed, opened, bought}), Prépositions des verbes d'argent (\eng{on, for, from, to, back}).
\item \textbf{Règle 50/30/20 :} 50\% Needs (Besoins), 30\% Wants (Envies), 20\% Savings (Épargne).
\item \textbf{Production écrite :} Résumer (50--60 mots, connecteurs, temps cohérent) + Donner son avis (\eng{In my opinion, I think that, From my point of view}).
\item \textbf{Pièges francophones :} \eng{broke} (adj.) vs \eng{broken} (pp), \eng{borrow from} vs \eng{lend to}, \eng{spend on} (pas \emph{in}), \eng{revenue} $\neq$ \eng{income}.
\end{itemize}
\end{aretenir}

\coursec{Comprehension activities}

Maintenant que nous avons lu le texte sur Tabi et ses finances personnelles, il est temps de vérifier notre compréhension. Dans cette section, nous allons appliquer les deux grandes stratégies de lecture — le \cle{skimming} (lecture rapide pour l'idée générale) et le \cle{scanning} (lecture ciblée pour les détails précis) — puis répondre aux différents types de questions que l'on rencontre au Probatoire et au Baccalauréat : vrai/faux, QCM, questions ouvertes et déduction du sens des mots.

\courssub{The reading text: a quick reminder}

Pour travailler les activités, relisons le texte support. Il raconte l'histoire de Tabi, un jeune employé de Douala qui apprend à gérer son argent.

\begin{textebox}[Managing personal finances: Tabi's story]
Tabi is a young shop assistant in Douala. Every month, he \textbf{earned} a salary, but by the middle of the month, he was always broke. He had no money left for transport, food or emergencies. His friends laughed at him because he could never afford to buy even a bottle of water at the end of the month.

One day, his uncle, a retired bank clerk from Buea, gave him some wise advice. “My boy,” he said, “money is like water: if you do not control it, it flows away.” He taught Tabi a simple method called the 50/30/20 rule. Every month, Tabi should spend 50\% of his income on needs, such as food, rent and transport; 30\% on wants, like entertainment and new clothes; and he should save 20\% for the future.

Tabi decided to try. His monthly salary was 150,000 FCFA. He started by writing down all his expenses in a small notebook. He was surprised: he spent 5,000 FCFA every month just on snacks and phone credit! He reduced this amount and opened a savings account at a microfinance bank near his shop.

After six months, Tabi had saved 180,000 FCFA. When his mother fell ill, he could pay the hospital bills without borrowing money. Tabi smiled. “Now I understand,” he told his uncle. “Managing money is not about how much you earn, but about how well you plan.”
\end{textebox}

\textbf{Glossaire :}

\begin{center}
\begin{tabularx}{\linewidth}{|l|l|X|}
\hline
\rowcolor{popblueL} \textbf{English} & \textbf{Nature} & \textbf{Français} \\
\hline
broke & adjectif (familier) & fauché, sans le sou \\
\hline
to afford & verbe & avoir les moyens de \\
\hline
wise & adjectif & sage, judicieux \\
\hline
clerk & nom & employé de bureau, guichetier \\
\hline
to flow away & verbe & s'écouler, filer \\
\hline
income & nom & revenu \\
\hline
needs / wants & noms & besoins / envies, désirs \\
\hline
expenses & nom pluriel & dépenses \\
\hline
savings account & nom & compte d'épargne \\
\hline
to borrow & verbe & emprunter (à quelqu'un) \\
\hline
\end{tabularx}
\end{center}

\courssub{First reading: skimming for the main idea}

Le \cle{skimming} consiste à lire rapidement le texte (titre, première phrase de chaque paragraphe, dernière phrase) pour saisir l'idée générale, sans s'arrêter sur les mots inconnus.

\begin{methode}[How to skim a text]
\begin{enumerate}
\item Read the title and the first sentence of each paragraph.
\item Read the last sentence of the text carefully.
\item Ask yourself: “Who? What? What is the general message?”
\item Choose the title that covers the \textbf{whole} text, not just one paragraph.
\end{enumerate}
\end{methode}

\begin{exoresolu}[Skimming: choose the best title]
Énoncé — Read the text quickly and choose the best title:
\begin{enumerate}
\item[a)] The life of a bank clerk in Buea
\item[b)] How Tabi learned to manage his money
\item[c)] Snacks and phone credit in Douala
\end{enumerate}
\tcblower
Solution détaillée — La bonne réponse est \textbf{b) How Tabi learned to manage his money}.
\begin{itemize}
\item L'option a) est un piège : l'oncle est bien un ancien employé de banque de Buea, mais ce n'est qu'un détail du deuxième paragraphe, pas le sujet du texte.
\item L'option c) est aussi un piège : les snacks et le crédit téléphonique ne sont qu'un exemple de dépense inutile (paragraphe 3).
\item L'option b) résume le texte entier : Tabi est fauché (problème), son oncle lui donne un conseil (solution), il épargne et réussit (résultat).
\end{itemize}
Au Bac, un bon titre doit couvrir \textbf{tout} le texte, du début à la fin.
\end{exoresolu}

\courssub{Second reading: scanning for specific details}

Le \cle{scanning} consiste à parcourir le texte rapidement pour repérer une information précise : un chiffre, un nom, une date. On ne lit pas tout : on laisse l'œil « glisser » jusqu'au mot recherché.

\begin{exercice}[Scanning: find the figures]
Scan the text and complete the table with the correct figures.
\begin{enumerate}
\item Tabi's monthly salary: \ldots\ldots\ldots
\item The percentage to spend on needs: \ldots\ldots\ldots
\item The percentage to spend on wants: \ldots\ldots\ldots
\item The percentage to save: \ldots\ldots\ldots
\item The amount Tabi spent on snacks and phone credit every month: \ldots\ldots\ldots
\item The amount Tabi saved after six months: \ldots\ldots\ldots
\end{enumerate}
\end{exercice}

\begin{corrige}[Exercice : scanning]
\begin{enumerate}
\item 150,000 FCFA (paragraphe 3 : “His monthly salary was 150,000 FCFA.”)
\item 50\% (paragraphe 2 : “50\% of his income on needs”)
\item 30\% (paragraphe 2 : “30\% on wants”)
\item 20\% (paragraphe 2 : “he should save 20\%”)
\item 5,000 FCFA (paragraphe 3 : “he spent 5,000 FCFA every month just on snacks and phone credit”)
\item 180,000 FCFA (paragraphe 4 : “Tabi had saved 180,000 FCFA”)
\end{enumerate}
\end{corrige}

\begin{attention}[Chiffres en anglais : attention aux séparateurs]
En anglais, on écrit \eng{150,000} avec une virgule pour les milliers, alors qu'en français on écrit « 150 000 » avec un espace. À l'oral, \eng{150,000 FCFA} se dit “one hundred and fifty thousand francs”. Ne traduisez jamais \eng{150,000} par « 150 » !
\end{attention}

\courssub{True or False: justify with a line from the text}

\begin{attention}[Toujours justifier au vrai/faux]
Au Probatoire et au Bac, une réponse \eng{True} ou \eng{False} \textbf{sans justification tirée du texte perd des points}, même si elle est correcte. Il faut recopier la ligne exacte du texte qui prouve votre réponse. Formule type : \eng{False. The text says: “...” (line ...)}.
\end{attention}

\begin{exoresolu}[True or False — justify your answers]
Énoncé — Say if the following statements are True or False. Justify each answer with a sentence from the text.
\begin{enumerate}
\item Tabi earned money every month.
\item Tabi's uncle worked in a school.
\item The 50/30/20 rule means saving 50\% of your income.
\item Tabi wrote his expenses in a notebook.
\item Tabi borrowed money to pay the hospital bills.
\end{enumerate}
\tcblower
Solution détaillée —
\begin{enumerate}
\item \textbf{True.} The text says: “Every month, he earned a salary.”
\item \textbf{False.} The text says: “his uncle, a retired bank clerk from Buea” — he worked in a bank, not in a school.
\item \textbf{False.} The text says: “he should save 20\% for the future” — the 50\% is for needs, not for savings.
\item \textbf{True.} The text says: “He started by writing down all his expenses in a small notebook.”
\item \textbf{False.} The text says: “he could pay the hospital bills without borrowing money.”
\end{enumerate}
Remarquez la méthode : pour chaque item, on repère les mots-clés (\eng{uncle}, \eng{50\%}, \eng{notebook}, \eng{hospital}), on retrouve la phrase dans le texte, puis on la recopie entre guillemets.
\end{exoresolu}

\courssub{Multiple choice questions (MCQ)}

Le QCM teste la compréhension fine : chaque option semble plausible, il faut donc retourner au texte pour éliminer les distracteurs.

\begin{exercice}[Choose the correct answer]
\begin{enumerate}
\item Why was Tabi always broke in the middle of the month?
\begin{enumerate}
\item[a)] Because he had no job.
\item[b)] Because he spent his money without planning.
\item[c)] Because his salary was too small.
\item[d)] Because he lent money to his friends.
\end{enumerate}
\item Who gave Tabi advice about money?
\begin{enumerate}
\item[a)] His mother.
\item[b)] His friends.
\item[c)] His uncle, a retired bank clerk.
\item[d)] A microfinance banker.
\end{enumerate}
\item What surprised Tabi when he wrote down his expenses?
\begin{enumerate}
\item[a)] He spent 5,000 FCFA a month on snacks and phone credit.
\item[b)] He earned 180,000 FCFA.
\item[c)] His rent was too expensive.
\item[d)] His notebook was too small.
\end{enumerate}
\item What happened when Tabi's mother fell ill?
\begin{enumerate}
\item[a)] He borrowed money from his uncle.
\item[b)] He paid the hospital bills with his savings.
\item[c)] He sold his notebook.
\item[d)] He stopped saving money.
\end{enumerate}
\end{enumerate}
\end{exercice}

\begin{corrige}[Exercice : QCM]
\begin{enumerate}
\item \textbf{b)} — Le texte montre qu'il dépensait sans compter (snacks, crédit téléphonique) ; son salaire existait bien, donc a) et c) sont faux.
\item \textbf{c)} — “his uncle, a retired bank clerk from Buea, gave him some wise advice.”
\item \textbf{a)} — “He was surprised: he spent 5,000 FCFA every month just on snacks and phone credit!”
\item \textbf{b)} — “he could pay the hospital bills without borrowing money”, grâce à ses économies.
\end{enumerate}
\end{corrige}

\courssub{Short answer questions}

\begin{methode}[How to answer a comprehension question]
\begin{enumerate}
\item \textbf{Underline the key words} of the question (soulignez les mots-clés : who, why, what, when + le mot important).
\item \textbf{Locate the paragraph} where the answer is (repérez le paragraphe grâce aux mots-clés).
\item \textbf{Answer with a complete sentence} (répondez par une phrase complète, jamais par un seul mot) : recopiez la phrase du texte ou reformulez-la avec vos propres mots.
\item \textbf{Check} : votre phrase doit avoir un sujet, un verbe et répondre exactement à la question posée.
\end{enumerate}
\end{methode}

\begin{exoresolu}[Short answer questions]
Énoncé — Answer with complete sentences.
\begin{enumerate}
\item Why was Tabi always broke?
\item What advice did he receive from his uncle?
\end{enumerate}
\tcblower
Solution détaillée —
\begin{enumerate}
\item \eng{Tabi was always broke because he spent all his money without planning, and he had no money left by the middle of the month.} (« Tabi était toujours fauché parce qu'il dépensait tout son argent sans planifier. »)
\item \eng{His uncle advised him to follow the 50/30/20 rule: spend 50\% of his income on needs, 30\% on wants, and save 20\% for the future.} (« Son oncle lui a conseillé de suivre la règle des 50/30/20. »)
\end{enumerate}
Piège à éviter : ne répondez pas seulement \eng{Because he had no money} — c'est une réponse incomplète qui ne montre pas que vous avez compris la cause (le manque de planification).
\end{exoresolu}

\courssub{Guessing meaning from context}

Au Bac, le texte contient toujours des mots que vous ne connaissez pas. Ce n'est pas grave : le contexte permet de deviner leur sens.

\begin{methode}[How to guess the meaning of a word]
\begin{enumerate}
\item \textbf{Look at the words around it} (regardez les mots qui l'entourent, avant et après).
\item \textbf{Identify the word class} (identifiez la nature du mot : nom, verbe, adjectif ?).
\item \textbf{Replace it with a synonym} (remplacez-le par un synonyme possible) et vérifiez que la phrase garde un sens.
\end{enumerate}
\end{methode}

\begin{exemplebox}[Worked examples: broke, wise, afford]
\begin{itemize}
\item \eng{Tabi was always broke.} → Le contexte donne la clé : \eng{he had no money left} (« il ne lui restait plus d'argent »). Nature : adjectif. Donc \cle{broke} = « fauché, sans le sou ». Vérification : « Tabi était toujours fauché » — la phrase a un sens.
\item \eng{His uncle gave him some wise advice.} → Le mot précède le nom \eng{advice} : c'est un adjectif. Le conseil (la règle 50/30/20) s'avère excellent. Donc \cle{wise} = « sage, judicieux ». Attention à la prononciation : « ouaïze ».
\item \eng{He could never afford to buy even a bottle of water.} → Nature : verbe suivi de \eng{to + infinitif}. Contexte : il n'avait plus d'argent. Donc \cle{to afford} = « avoir les moyens de, pouvoir se permettre ».
\end{itemize}
\end{exemplebox}

\begin{attention}[Faux amis et pièges de traduction]
\begin{itemize}
\item \eng{Broke} ne veut pas dire « cassé » ici (sens littéral : \eng{a broke chair} est incorrect ; on dirait \eng{a broken chair}). Devant une personne, \eng{broke} = « fauché ».
\item \eng{To afford} ne veut pas dire « offrir » : \eng{I can afford a taxi} = « J'ai les moyens de prendre un taxi ».
\item \eng{To borrow} = « emprunter (à quelqu'un) » ; ne le confondez pas avec \eng{to lend} = « prêter (à quelqu'un) ». \eng{Can I borrow your pen?} / \eng{Can you lend me your pen?}
\end{itemize}
\end{attention}

\begin{center}
\begin{tabularx}{\linewidth}{|X|X|X|}
\hline
\rowcolor{popblueL} \textbf{English} & \textbf{Français} & \textbf{Remarque} \\
\hline
to be broke & être fauché & familier ; adjectif, pas de \eng{a} devant \\
\hline
wise advice & un conseil judicieux & prononciation « ouaïze » \\
\hline
to afford to buy & avoir les moyens d'acheter & suivi de \eng{to + infinitif} \\
\hline
to save money & économiser de l'argent & \eng{save}, pas \eng{spare} ici \\
\hline
to borrow from & emprunter à & \eng{borrow from somebody} \\
\hline
to lend to & prêter à & \eng{lend something to somebody} \\
\hline
\end{tabularx}
\end{center}

\courssub{Interpretation question: giving your opinion}

La dernière question d'une épreuve de compréhension demande souvent votre avis personnel. Il n'y a pas de « bonne » réponse : il faut donner une opinion et la justifier en deux ou trois phrases, en anglais correct.

\begin{exoresolu}[Interpretation: your opinion]
Énoncé — \eng{Do you think the 50/30/20 rule can work in Cameroon? Why?} (Answer in two or three sentences.)
\tcblower
Solution détaillée — Modèle de réponse possible :
\eng{Yes, I think the 50/30/20 rule can work in Cameroon, because many people spend money without planning, like Tabi. Writing down expenses helps families to control their budget. However, for people with very small salaries, saving 20\% may be difficult, so they can start with 5\% or 10\%.}
(« Oui, je pense que la règle peut fonctionner au Cameroun, car beaucoup de gens dépensent sans planifier. Noter ses dépenses aide à contrôler son budget. Cependant, pour les petits salaires, épargner 20\% peut être difficile, donc on peut commencer avec 5\% ou 10\%. »)
Structure à retenir : \textbf{opinion} (\eng{I think...}) + \textbf{raison} (\eng{because...}) + \textbf{nuance ou exemple} (\eng{However... / For example...}).
\end{exoresolu}

\courssub{Writing a summary}

Au-delà de l'opinion, l'épreuve écrite peut demander un \cle{résumé} du texte (souvent 40 à 60 mots). Il faut sélectionner les informations essentielles (qui, quoi, problème, solution, résultat) sans copier-coller de longues phrases, et utiliser ses propres mots (reformulation).

\begin{methode}[How to write a summary]
\begin{enumerate}
\item \textbf{List the key points} (notez les idées principales : 1 par paragraphe).
\item \textbf{Link them with connectors} (reliez-les avec \eng{First, Then, After that, Finally, However, So}).
\item \textbf{Write a draft} (rédigez un brouillon en phrases simples, au passé pour une narration).
\item \textbf{Count words and edit} (comptez les mots ; supprimez les détails inutiles, gardez l'essentiel).
\item \textbf{Final version} (recopiez proprement sans dépasser la limite).
\end{enumerate}
\end{methode}

\begin{exoresolu}[Summary writing]
Énoncé — Summarise Tabi's story in about 50 words.
\tcblower
Solution détaillée —
\eng{Tabi, a shop assistant in Douala, used to spend all his salary before mid-month. His uncle, a retired bank clerk, taught him the 50/30/20 budgeting rule. Tabi tracked his expenses, cut unnecessary costs like snacks, and opened a savings account. After six months, he had saved 180,000 FCFA, which allowed him to pay his mother's hospital bills without borrowing. He learned that planning matters more than income.}
(53 mots)
\end{exoresolu}

\courssub{Summary: the five steps of reading}

\begin{popfigure}
\begin{tikzpicture}[node distance=0.55cm,
etape/.style={rectangle, rounded corners, draw=popblue, fill=popblueL, text width=2.3cm, align=center, minimum height=0.9cm, font=\small\bfseries},
fleche/.style={-{Stealth[length=2.5mm]}, thick, popdark}]
\node[etape, align=center] (p){1. Predict\\ \footnotesize deviner avec le titre};
\node[etape, right=of p, fill=poporangeL, draw=poporange, align=center] (sk){2. Skim\\ \footnotesize idée générale};
\node[etape, right=of sk, fill=popgreenL, draw=popgreen, align=center] (sc){3. Scan\\ \footnotesize détails, chiffres};
\node[etape, right=of sc, fill=poppurpleL, draw=poppurple, align=center] (a){4. Answer\\ \footnotesize phrase complète};
\node[etape, right=of a, fill=popgoldL, draw=popgold, align=center] (c){5. Check\\ \footnotesize justifier, relire};
\draw[fleche] (p) -- (sk);
\draw[fleche] (sk) -- (sc);
\draw[fleche] (sc) -- (a);
\draw[fleche] (a) -- (c);
\end{tikzpicture}
\legende{Les cinq étapes de la lecture efficace : Predict, Skim, Scan, Answer, Check.}
\end{popfigure}

\begin{aretenir}[À retenir]
\begin{itemize}
\item \cle{Skimming} = lire vite pour l'idée générale (choisir un titre, résumer).
\item \cle{Scanning} = chercher une information précise (chiffres, noms, dates).
\item Vrai/Faux : toujours \textbf{justifier en citant une ligne du texte}.
\item Questions ouvertes : répondre par une \textbf{phrase complète} (sujet + verbe).
\item Mot inconnu : contexte → nature du mot → synonyme → vérification.
\item Résumé : idées clés + connecteurs + reformulation + limite de mots.
\end{itemize}
\end{aretenir}

\begin{savaistu}[Reading comprehension at the Bac]
À l'épreuve d'anglais du Probatoire et du Baccalauréat camerounais, la section \eng{Reading Comprehension} vaut généralement \textbf{10 points sur 20}. Les questions suivent presque toujours l'ordre du texte : la réponse à la question 1 se trouve avant la réponse à la question 2. Utilisez ce repère pour gagner du temps le jour de l'examen !
\end{savaistu}

\coursec{Thematic vocabulary}

Dans la section précédente, nous avons lu le texte sur la gestion des finances personnelles et répondu aux questions de compréhension. Vous avez remarqué que ce texte regorge de mots liés à l'argent : \eng{income}, \eng{expenses}, \eng{savings}, \eng{debt}\ldots{} Cette quatrième section a pour but de \textbf{systématiser ce vocabulaire} : nous allons l'organiser en champs lexicaux clairs, apprendre les collocations essentielles, déjouer les faux amis et construire les familles de mots. Un bon vocabulaire thématique, c'est la clé pour parler et écrire naturellement sur l'argent.

\courssub{Overview: a lexical mind map}

Avant d'entrer dans le détail, observons la carte mentale ci-dessous. Elle organise tout le vocabulaire de la leçon autour de quatre branches : l'argent qui \textbf{entre} (\eng{money in}), l'argent qui \textbf{sort} (\eng{money out}), l'épargne et l'emprunt (\eng{saving and borrowing}) et les adjectifs qui qualifient les habitudes (\eng{habits}).

\begin{popfigure}
\begin{tikzpicture}[
  mindmap,
  grow cyclic,
  every node/.style={concept, font=\small, align=center},
  concept color=popblue,
  root concept/.append style={concept color=popdark, font=\small\bfseries, text=white},
  level 1/.append style={level distance=3.4cm, sibling angle=90},
  level 2/.append style={level distance=2.4cm, font=\scriptsize},
]
\node[root concept]{PERSONAL FINANCES}
  child[concept color=popgreen]{ node {Money in}
    child { node {income, salary} }
    child { node {to earn, pocket money} }
  }
  child[concept color=poporange]{ node {Money out}
    child { node {expenses, bills, fees} }
    child { node {to spend, to pay, to buy} }
  }
  child[concept color=popblue]{ node {Saving and borrowing}
    child { node {savings, bank account} }
    child { node {to borrow, to lend, loan, debt} }
  }
  child[concept color=poppurple]{ node {Habits}
    child { node {wise, careless, broke} }
    child { node {affordable, expensive} }
  };
\end{tikzpicture}
\legende{Carte mentale du vocabulaire des finances personnelles : quatre champs lexicaux à mémoriser ensemble.}
\end{popfigure}

\begin{methode}[Comment apprendre un champ lexical]
\begin{enumerate}
\item \textbf{Observez} les mots dans le texte : ne les apprenez jamais isolément, mais dans leur phrase.
\item \textbf{Classez-les} par famille de sens (argent qui entre / qui sort / qu'on garde).
\item \textbf{Notez les collocations} : on apprend \eng{to earn money} et non le seul mot \eng{earn}.
\item \textbf{Fabriquez vos propres phrases} sur votre vie réelle (ton argent de poche, les dépenses de ta famille).
\item \textbf{Réutilisez} les mots à l'oral et à l'écrit dans la semaine : un mot utilisé trois fois est un mot acquis.
\end{enumerate}
\end{methode}

\courssub{Field 1 — Money in: the money that comes in}

Observons d'abord ces phrases tirées de situations réelles :

\eng{My father earns a salary of 150,000 francs a month.} « Mon père gagne un salaire de 150 000 francs par mois. »

\eng{My aunt gives me pocket money every Sunday.} « Ma tante me donne de l'argent de poche chaque dimanche. »

\eng{She received her first income from her small business.} « Elle a reçu ses premiers revenus de sa petite entreprise. »

\begin{definition}[Money in — l'argent qui entre]
\begin{itemize}
\item \cle{income} (nom, indénombrable) : le \textbf{revenu}, c'est-à-dire l'ensemble de l'argent qu'une personne reçoit régulièrement (salaire, bénéfices, loyers). \eng{His monthly income is not very high.} « Son revenu mensuel n'est pas très élevé. »
\item \cle{salary} (nom) : le \textbf{salaire}, la somme fixe versée chaque mois à un employé. \eng{Teachers receive their salary at the end of the month.} « Les enseignants reçoivent leur salaire à la fin du mois. »
\item \cle{pocket money} (nom, indénombrable) : l'\textbf{argent de poche} donné aux enfants ou adolescents. \eng{I save half of my pocket money.} « J'économise la moitié de mon argent de poche. »
\item \cle{to earn} (verbe) : \textbf{gagner} de l'argent par son travail. \eng{She earns money by selling doughnuts.} « Elle gagne de l'argent en vendant des beignets. »
\item \cle{to receive} (verbe) : \textbf{recevoir} (plus général, sans idée de travail). \eng{He received 10,000 francs from his uncle.} « Il a reçu 10 000 francs de son oncle. »
\end{itemize}
\end{definition}

\begin{attention}[To earn ou to win ?]
En français, « gagner » a deux sens que l'anglais distingue strictement :
\begin{itemize}
\item \eng{to earn} = gagner de l'argent \textbf{par son travail} : \eng{He earns 50,000 francs a month.} « Il gagne 50 000 francs par mois. »
\item \eng{to win} = gagner \textbf{à un jeu, un concours, un match} : \eng{She won 100,000 francs in the lottery.} « Elle a gagné 100 000 francs à la loterie. »
\end{itemize}
Ne dites jamais \eng{I win a salary} : c'est une erreur typique de francophone.
\end{attention}

\courssub{Field 2 — Money out: the money that goes out}

\eng{Every month, my mother pays the electricity bill and the school fees.} « Chaque mois, ma mère paie la facture d'électricité et les frais de scolarité. »

\eng{Teenagers often spend too much money on phone credit.} « Les adolescents dépensent souvent trop d'argent en crédit téléphonique. »

\begin{definition}[Money out — l'argent qui sort]
\begin{itemize}
\item \cle{expenses} (nom, au pluriel) : les \textbf{dépenses}. \eng{Food and transport are my main expenses.} « La nourriture et le transport sont mes principales dépenses. »
\item \cle{to spend} (verbe irrégulier : spend, spent, spent) : \textbf{dépenser}. \eng{He spent all his money at the market.} « Il a dépensé tout son argent au marché. »
\item \cle{to pay} (verbe irrégulier : pay, paid, paid) : \textbf{payer}. \eng{She paid the taxi driver.} « Elle a payé le chauffeur de taxi. »
\item \cle{to buy} (verbe irrégulier : buy, bought, bought) : \textbf{acheter}. \eng{They bought a new television.} « Ils ont acheté une nouvelle télévision. »
\item \cle{bill} (nom) : la \textbf{facture} (électricité, eau, téléphone, restaurant). \eng{The water bill is very high this month.} « La facture d'eau est très élevée ce mois-ci. »
\item \cle{fees} (nom, au pluriel) : les \textbf{frais} (scolarité, inscription, services). \eng{School fees must be paid in September.} « Les frais de scolarité doivent être payés en septembre. »
\end{itemize}
\end{definition}

\begin{propriete}[Les constructions de to spend et to pay]
Ces deux verbes ont des constructions différentes qu'il faut mémoriser :
\begin{itemize}
\item \eng{to spend money \textbf{on} something} : \eng{I spend 500 francs on transport every day.} « Je dépense 500 francs en transport chaque jour. »
\item \eng{to spend time/money \textbf{doing} something} : \eng{She spent two hours shopping.} « Elle a passé deux heures à faire les courses. »
\item \eng{to pay \textbf{for} something} : \eng{He paid for the books.} « Il a payé les livres. »
\item \eng{to pay somebody} : \eng{They paid the mechanic.} « Ils ont payé le mécanicien. »
\end{itemize}
\end{propriete}

\courssub{Field 3 — Saving and borrowing: keeping and lending money}

\eng{If you save a little money every week, you will never need to borrow.} « Si tu épargnes un peu d'argent chaque semaine, tu n'auras jamais besoin d'emprunter. »

\begin{definition}[Saving and borrowing — épargner et emprunter]
\begin{itemize}
\item \cle{savings} (nom, au pluriel) : l'\textbf{épargne}, l'argent mis de côté. \eng{She keeps her savings in a bank account.} « Elle garde son épargne sur un compte bancaire. »
\item \cle{to save} (verbe) : \textbf{épargner, économiser}. \eng{I save 2,000 francs every month.} « J'économise 2 000 francs chaque mois. »
\item \cle{a bank account} (nom) : un \textbf{compte bancaire}. \eng{My brother opened a bank account in Douala.} « Mon frère a ouvert un compte bancaire à Douala. »
\item \cle{to borrow} (verbe) : \textbf{emprunter} (prendre quelque chose à quelqu'un pour le rendre). \eng{Can I borrow your bicycle?} « Puis-je emprunter ton vélo ? »
\item \cle{to lend} (verbe irrégulier : lend, lent, lent) : \textbf{prêter} (donner temporairement à quelqu'un). \eng{Can you lend me your pen?} « Peux-tu me prêter ton stylo ? »
\item \cle{a loan} (nom) : un \textbf{prêt} (somme prêtée par une banque). \eng{The bank gave him a loan to start his shop.} « La banque lui a accordé un prêt pour démarrer sa boutique. »
\item \cle{debt} (nom) : la \textbf{dette} (attention, le \emph{b} ne se prononce pas : « dèt »). \eng{He is in debt because he borrowed too much.} « Il est endetté parce qu'il a trop emprunté. »
\end{itemize}
\end{definition}

\begin{attention}[To borrow ou to lend ? Le piège classique]
Ces deux verbes décrivent la même action, mais vue de deux côtés opposés :
\begin{itemize}
\item \eng{to borrow something \textbf{from} somebody} = emprunter quelque chose \textbf{à} quelqu'un (je prends). \eng{Can I borrow your pen?} « Puis-je emprunter ton stylo ? »
\item \eng{to lend something \textbf{to} somebody} = prêter quelque chose \textbf{à} quelqu'un (je donne). \eng{Can you lend me your pen?} « Peux-tu me prêter ton stylo ? »
\end{itemize}
Erreur fréquente : \eng{Please borrow me your book.} est faux. Dites : \eng{Please lend me your book.} ou \eng{Can I borrow your book?}
\end{attention}

\courssub{Field 4 — Qualifying money habits}

\eng{A wise student saves money; a careless one is always broke.} « Un élève avisé épargne ; un élève insouciant est toujours fauché. »

\begin{definition}[Adjectifs pour qualifier les habitudes financières]
\begin{itemize}
\item \cle{wise} (adj.) : \textbf{sage, avisé, prudent}. \eng{It is wise to keep some money for emergencies.} « Il est sage de garder de l'argent pour les urgences. »
\item \cle{careless} (adj.) : \textbf{insouciant, négligent}. \eng{He is careless with his money.} « Il est insouciant avec son argent. »
\item \cle{broke} (adj., familier) : \textbf{fauché}, sans argent. \eng{I can't go out; I'm broke!} « Je ne peux pas sortir ; je suis fauché ! »
\item \cle{rich} (adj.) : \textbf{riche}. \eng{She became rich thanks to her cocoa business.} « Elle est devenue riche grâce à son commerce de cacao. »
\item \cle{poor} (adj.) : \textbf{pauvre}. \eng{The family is poor but honest.} « La famille est pauvre mais honnête. »
\item \cle{affordable} (adj.) : \textbf{abordable} (qu'on peut se permettre d'acheter). \eng{This phone is affordable for students.} « Ce téléphone est abordable pour les étudiants. »
\item \cle{expensive} (adj.) : \textbf{cher, coûteux}. \eng{Bread is more expensive than last year.} « Le pain est plus cher que l'an dernier. »
\end{itemize}
\end{definition}

\begin{attention}[Cher = expensive ou dear ?]
Pour dire qu'une chose « coûte cher », utilisez \eng{expensive} : \eng{This bag is very expensive.} L'adjectif \eng{dear} signifie « cher » au sens affectif (\eng{my dear friend} = « mon cher ami ») et ne s'emploie presque jamais pour les prix en anglais moderne. À l'inverse, « pas cher » se dit \eng{cheap} ou \eng{affordable}.
\end{attention}

\courssub{Summary table: the essential words}

\begin{center}
\begin{tabularx}{\linewidth}{|l|X|X|}
\hline
\rowcolor{popblueL} English & Français & Exemple \\
\hline
income & revenu & \eng{Her income is 200,000 francs a year.} \\
\hline
salary & salaire & \eng{He receives his salary monthly.} \\
\hline
pocket money & argent de poche & \eng{I get pocket money on Saturdays.} \\
\hline
expenses & dépenses & \eng{Transport is a big expense.} \\
\hline
budget & budget & \eng{We must respect our budget.} \\
\hline
savings & épargne & \eng{Her savings are in the bank.} \\
\hline
debt & dette & \eng{He paid all his debts.} \\
\hline
loan & prêt & \eng{She asked the bank for a loan.} \\
\hline
bill & facture & \eng{The electricity bill arrived today.} \\
\hline
fees & frais & \eng{School fees are expensive.} \\
\hline
\end{tabularx}
\end{center}

\courssub{Collocations: words that go together}

En anglais, certains verbes s'associent naturellement avec \eng{money}. Ces combinaisons s'appellent des \textbf{collocations} : il faut les apprendre comme des blocs.

\begin{center}
\begin{tabularx}{\linewidth}{|X|X|}
\hline
\rowcolor{popgreenL} Collocation & Traduction et exemple \\
\hline
\eng{to make money} & gagner de l'argent (se dit aussi \eng{to earn money}) — \eng{She makes money by selling clothes.} \\
\hline
\eng{to save money} & économiser de l'argent — \eng{We save money for the holidays.} \\
\hline
\eng{to waste money} & gaspiller de l'argent — \eng{Don't waste your money on useless things.} « Ne gaspille pas ton argent en choses inutiles. » \\
\hline
\eng{to manage a budget} & gérer un budget — \eng{My father manages the family budget carefully.} \\
\hline
\eng{to run into debt} & s'endetter — \eng{He ran into debt after buying a car.} \\
\hline
\eng{to earn money} & gagner de l'argent — \eng{She earns 50,000 francs a month.} « Elle gagne 50 000 francs par mois. » \\
\hline
\eng{to spend money on} & dépenser de l'argent pour — \eng{He spends money on video games.} \\
\hline
\eng{to borrow money from} & emprunter de l'argent à — \eng{I borrowed money from my sister.} \\
\hline
\end{tabularx}
\end{center}

\courssub{Word families: building words from one root}

Connaître les familles de mots permet de multiplier son vocabulaire sans effort.

\begin{center}
\begin{tabular}{|l|l|l|}
\hline
\rowcolor{poppurpleL} Verbe & Nom(s) & Adjectif / participe \\
\hline
to save & savings, a saver & saved \\
\hline
to spend (spent, spent) & spending & spent \\
\hline
to borrow & borrowing & borrowed \\
\hline
to lend (lent, lent) & a loan, a lender & lent \\
\hline
to earn & earnings & earned \\
\hline
to pay (paid, paid) & payment & paid \\
\hline
\end{tabular}
\end{center}

\eng{She is a good saver: her savings grow every month.} « C'est une bonne épargnante : son épargne augmente chaque mois. »

\eng{My spending is too high this term.} « Mes dépenses sont trop élevées ce trimestre. »

\begin{attention}[Money est indénombrable]
En anglais, \eng{money} est \textbf{indénombrable} : il n'a jamais de pluriel et s'accorde au singulier.
\begin{itemize}
\item Correct : \eng{Money is important.} « L'argent est important. »
\item Faux : \eng{Moneys are important.}
\item Pour compter, utilisez une unité : \eng{a coin} (une pièce), \eng{a note / a banknote} (un billet), \eng{some money} (de l'argent).
\end{itemize}
Même remarque pour \eng{income} et \eng{pocket money}, indénombrables eux aussi.
\end{attention}

\begin{attention}[Faux amis : économie et dépenser]
\begin{itemize}
\item « L'économie » du pays se dit \eng{the economy} : \eng{The economy of Cameroon is growing.} Mais « l'économie » comme matière scolaire se dit \eng{economics} : \eng{She studies economics at university.}
\item « Économiser » (de l'argent) se dit \eng{to save (money)}, jamais \eng{to economize money}.
\item « Dépenser » se dit \eng{to spend} ; le verbe \eng{to dépense} n'existe pas. Le nom « dépense » se dit \eng{expense} ou \eng{spending}.
\end{itemize}
\end{attention}

\begin{savaistu}[Pocket money or allowance?]
En anglais \textbf{britannique}, l'argent de poche se dit \eng{pocket money}. En anglais \textbf{américain}, on dit \eng{allowance}. Les deux expressions sont correctes : à l'examen, utilisez celle que vous avez apprise, mais soyez capable de comprendre les deux. Autre différence célèbre : un billet de banque est \eng{a note} en Grande-Bretagne et \eng{a bill} aux États-Unis — alors qu'en Grande-Bretagne, \eng{a bill} désigne une facture !
\end{savaistu}

\courssub{Practice text: recycling the vocabulary}

Lisons maintenant un court texte qui réutilise presque tout le vocabulaire de cette section. Soulignez mentalement chaque mot de la carte mentale que vous reconnaissez.

\begin{textebox}[A Young Saver from Bamenda]
Manka is a seventeen-year-old student in Bamenda. Every month, she receives pocket money from her parents and earns a little extra income by plaiting hair for her neighbours on Saturdays. Unlike many teenagers, Manka is very wise with money. She divides everything she receives into three parts. The first part is for her daily expenses: transport to school, food and phone credit. The second part goes into her savings. She opened a bank account last year, and she deposits money there every month. The third part is for emergencies.

Last term, her best friend Linda wanted to borrow 5,000 francs from her to buy new shoes. Manka refused politely. “I can lend you my old shoes,” she said with a smile, “but I never lend money, because I don't want to run into debt myself.” Linda was not angry; she understood that Manka was managing her budget carefully.

At the end of the year, Manka had saved 60,000 francs. With this money, she paid her school fees for the following year. Her father was very proud. “You are not rich,” he told her, “but you are richer than many careless adults, because you know how to manage money.” Manka smiled. She knew that being wise with money is better than being rich.
\end{textebox}

\textbf{Glossaire :} \emph{to plait hair} (v.) : tresser les cheveux ; \emph{unlike} (prép.) : contrairement à ; \emph{to deposit} (v.) : déposer (de l'argent) ; \emph{emergency} (n.) : urgence, imprévu ; \emph{proud} (adj.) : fier.

\textbf{Questions de compréhension :}
\begin{enumerate}
\item How does Manka earn extra income?
\item Into how many parts does she divide her money? Name them.
\item Why did Manka refuse to lend money to Linda?
\item What did she do with her savings at the end of the year?
\item What lesson does her father teach at the end of the text?
\end{enumerate}

\courssub{Worked exercise}

\begin{exoresolu}[Complete with the correct word]
Complétez les phrases avec le mot qui convient, choisi dans la liste : \eng{borrow — lend — earn — save — debt — expenses — broke — affordable}.

\begin{enumerate}
\item Can I \ldots\ldots\ldots your dictionary? Mine is at home.
\item My father works hard to \ldots\ldots\ldots money for the family.
\item Food and transport are my biggest \ldots\ldots\ldots every month.
\item If you \ldots\ldots\ldots 1,000 francs every week, you will have 52,000 francs at the end of the year.
\item He bought an expensive car and now he is in \ldots\ldots\ldots.
\item Sorry, I can't pay for your drink: I'm completely \ldots\ldots\ldots!
\item This school bag is \ldots\ldots\ldots: it costs only 3,000 francs.
\item Banks \ldots\ldots\ldots money to people who want to start a business.
\end{enumerate}

\tcblower

\textbf{Solution détaillée :}
\begin{enumerate}
\item \eng{borrow} — le sujet « prend » le dictionnaire : \eng{Can I borrow your dictionary?} (« Puis-je emprunter ton dictionnaire ? »)
\item \eng{earn} — gagner de l'argent par le travail : \eng{to earn money}.
\item \eng{expenses} — nourriture et transport sont des \textbf{dépenses}.
\item \eng{save} — mettre de l'argent de côté chaque semaine : \eng{to save}.
\item \eng{debt} — la collocation est \eng{to be in debt} (« être endetté »).
\item \eng{broke} — « fauché », sans argent (adjectif familier).
\item \eng{affordable} — un prix bas qu'on peut se permettre : « abordable ».
\item \eng{lend} — la banque « donne » temporairement l'argent : \eng{Banks lend money to people}. (On pourrait aussi dire \eng{People borrow money from banks}.)
\end{enumerate}
\end{exoresolu}

\begin{exercice}[Your turn: build your own sentences]
Écrivez une phrase complète en anglais avec chacun des mots suivants, en parlant de votre propre vie ou de votre famille : \eng{pocket money}, \eng{to save}, \eng{bill}, \eng{wise}, \eng{to waste money}. Puis traduisez chaque phrase en français.
\end{exercice}

\begin{corrige}[Exercice « Your turn »]
Corrigé indicatif (vos phrases peuvent différer ; vérifiez la grammaire et la collocation) :
\begin{enumerate}
\item \eng{I receive 2,000 francs of pocket money every month.} « Je reçois 2 000 francs d'argent de poche chaque mois. »
\item \eng{I save some money to buy a school bag.} « J'économise de l'argent pour acheter un cartable. »
\item \eng{My mother pays the electricity bill every month.} « Ma mère paie la facture d'électricité chaque mois. »
\item \eng{It is wise to keep money for emergencies.} « Il est sage de garder de l'argent pour les imprévus. »
\item \eng{I never waste money on useless things.} « Je ne gaspille jamais mon argent en choses inutiles. »
\end{enumerate}
\end{corrige}

\begin{aretenir}[Ce qu'il faut retenir de cette section]
\begin{itemize}
\item Quatre champs lexicaux : \eng{money in} (\eng{income, salary, pocket money, to earn}), \eng{money out} (\eng{expenses, bills, fees, to spend, to pay}), \eng{saving and borrowing} (\eng{savings, bank account, to borrow, to lend, loan, debt}) et les adjectifs (\eng{wise, careless, broke, affordable, expensive}).
\item \eng{to borrow from} = emprunter à ; \eng{to lend to} = prêter à. Ne les inversez jamais.
\item \eng{money} est indénombrable : \eng{money is}, jamais \eng{moneys are}.
\item « Économie » (pays) = \eng{economy} ; (matière) = \eng{economics} ; « dépenser » = \eng{to spend}.
\item Apprenez les collocations en blocs : \eng{to make money, to save money, to waste money, to manage a budget, to run into debt}.
\end{itemize}
\end{aretenir}

\coursec{Reading method and worked example}

Dans la section précédente, nous avons enrichi notre vocabulaire thématique sur la gestion des finances personnelles : \eng{income}, \eng{expenses}, \eng{savings}, \eng{budget}, \eng{to afford}... Mais connaître les mots ne suffit pas. Au baccalauréat comme à l'examen de fin d'année, la réussite en \cle{reading comprehension} repose sur une \cle{méthode}. Un bon lecteur ne lit pas un texte mot à mot : il applique une stratégie en cinq étapes. C'est cette stratégie que nous allons maintenant apprendre, puis appliquer pas à pas sur un mini-texte.

\courssub{The five steps of efficient reading}

\begin{definition}[Reading strategies]
La \cle{compréhension écrite} efficace repose sur cinq étapes ordonnées : \eng{predict} (prédire le contenu à partir du titre), \eng{skim} (survoler pour saisir l'idée générale), \eng{scan} (repérer une information précise), \eng{answer} (rédiger la réponse en phrase complète) et \eng{check} (vérifier la grammaire et la pertinence de la réponse).
\end{definition}

\begin{propriete}[The five steps — form and use]
\begin{enumerate}
\item \textbf{Predict} (environ 1 min) : lire le titre, regarder les éventuelles illustrations, et se demander : \eng{What is this text about? What do I already know about this topic?}
\item \textbf{Skim} (environ 2 min) : lire rapidement le texte une première fois, sans s'arrêter sur les mots inconnus, pour dégager l'\cle{idée générale} en une phrase.
\item \textbf{Scan} (environ 3 min) : relire en cherchant des \cle{informations précises} (chiffres, noms, dates, lieux) qui répondent aux questions. L'œil saute de mot-clé en mot-clé.
\item \textbf{Answer} : rédiger chaque réponse en \cle{phrase complète}, courte et précise, en reprenant les mots de la question.
\item \textbf{Check} (environ 2 min) : relire chaque réponse et vérifier l'orthographe, l'accord sujet-verbe et la cohérence avec le texte.
\end{enumerate}
\end{propriete}

\begin{popfigure}
\begin{tikzpicture}[font=\small]
\node[draw=popblue, fill=popblueL, rounded corners, minimum width=2.2cm, minimum height=1cm, align=center] (a) at (0,0) {\textbf{Predict}\\ 1 min};
\node[draw=poporange, fill=poporangeL, rounded corners, minimum width=2.2cm, minimum height=1cm, align=center] (b) at (3,0) {\textbf{Skim}\\ 2 min};
\node[draw=popgreen, fill=popgreenL, rounded corners, minimum width=2.2cm, minimum height=1cm, align=center] (c) at (6,0) {\textbf{Scan}\\ 3 min};
\node[draw=poppurple, fill=poppurpleL, rounded corners, minimum width=2.2cm, minimum height=1cm, align=center] (d) at (9,0) {\textbf{Answer}\\ 3 min};
\node[draw=popgold, fill=popgoldL, rounded corners, minimum width=2.2cm, minimum height=1cm, align=center] (e) at (12,0) {\textbf{Check}\\ 2 min};
\draw[-{Stealth}, thick, popdark] (a) -- (b);
\draw[-{Stealth}, thick, popdark] (b) -- (c);
\draw[-{Stealth}, thick, popdark] (c) -- (d);
\draw[-{Stealth}, thick, popdark] (d) -- (e);
\end{tikzpicture}
\legende{Frise des cinq étapes de lecture avec durées indicatives pour un texte court.}
\end{popfigure}

\begin{methode}[Réussir la compréhension écrite au Bac]
Voici la démarche complète à adopter le jour de l'examen :
\begin{enumerate}
\item \textbf{Lisez d'abord les questions, avant le texte.} Ainsi, vous savez exactement quoi chercher : un chiffre ? un lieu ? une opinion ? Votre lecture devient ciblée et vous gagnez un temps précieux.
\item Lisez ensuite le titre et prédisez le contenu (\eng{predict}).
\item Survolez le texte pour l'idée générale (\eng{skim}) : ne traduisez pas mentalement chaque mot.
\item Pour chaque question, repérez dans le texte le passage qui contient la réponse (\eng{scan}) et soulignez-le.
\item Rédigez une réponse courte, en phrase complète, avec vos propres mots si possible.
\item Relisez tout : grammaire, orthographe, logique.
\end{enumerate}
\end{methode}

\courssub{Worked example: Mary's market budget}

Appliquons maintenant les cinq étapes à un mini-texte. Lisez-le d'abord en entier.

\begin{textebox}[Mary's market budget]
Mary is a student in Bamenda. Every Saturday, she goes to the market with a small budget of 5,000 francs. She spends 2,000 francs on vegetables, 1,500 francs on fish and 1,000 francs on fruit. She keeps the remaining 500 francs in her savings box at home. Mary says that managing her money helps her to avoid waste. At the end of each month, she uses her savings to buy exercise books for school.
\end{textebox}

\textbf{Glossaire :} \eng{budget} (n.m.) : budget ; \eng{to spend (spent, spent)} (v.) : dépenser ; \eng{remaining} (adj.) : restant ; \eng{savings box} (n.) : tirelire, boîte d'épargne ; \eng{waste} (n.m.) : gaspillage ; \eng{exercise books} (n.) : cahiers ; \eng{to keep (kept, kept)} (v.) : garder, conserver ; \eng{to avoid} (v.) : éviter.

\textbf{Étape 1 — Predict.}\par
Le titre est \eng{Mary's market budget}. Avant même de lire, nous pouvons prédire : le texte parlera probablement d'une personne nommée Mary, de ses achats au marché, et de la manière dont elle répartit son argent. Cette prédiction active nos connaissances (le lexique du marché et de l'argent vu dans la section précédente) et facilite la lecture.

\textbf{Étape 2 — Skim.}\par
Une lecture rapide suffit pour dégager l'idée générale en une phrase : \eng{The text explains how Mary, a student in Bamenda, manages her small weekly budget at the market.} « Le texte explique comment Mary, une élève de Bamenda, gère son petit budget hebdomadaire au marché. » Remarquez que nous n'avons pas cherché à comprendre chaque détail : seulement le sujet global.

\textbf{Étape 3 — Scan.}\par
Prenons la question : \eng{How much does Mary spend on vegetables?} Les mots-clés sont \eng{how much} et \eng{vegetables}. Notre œil parcourt le texte à la recherche du mot \eng{vegetables} et d'un chiffre associé. Nous trouvons : \eng{She spends 2,000 francs on vegetables}. L'information cherchée est : \cle{2,000 francs}.

\textbf{Étape 4 — Answer.}\par
On rédige une phrase complète qui reprend les éléments de la question : \eng{Mary spends 2,000 francs on vegetables.} « Mary dépense 2 000 francs en légumes. »

\textbf{Étape 5 — Check.}\par
Relisons la réponse : le sujet \eng{Mary} est à la troisième personne du singulier, donc le verbe prend un \textbf{-s} : \eng{spends}. Le chiffre correspond bien au texte. La phrase est courte et correcte. Réponse validée.

\begin{exoresolu}[Applying the five steps]
Énoncé — Using the text \eng{Mary's market budget}, answer the question: \eng{How much does Mary spend on vegetables?}
\tcblower
Solution détaillée —
\begin{enumerate}
\item \textbf{Predict} : le titre annonce un texte sur le budget de Mary au marché.
\item \textbf{Skim} : idée générale = Mary gère un budget de 5 000 francs chaque samedi.
\item \textbf{Scan} : je repère le mot \eng{vegetables} dans la phrase \eng{She spends 2,000 francs on vegetables}. Le chiffre cherché est \textbf{2,000 francs}.
\item \textbf{Answer} : \eng{Mary spends 2,000 francs on vegetables.}
\item \textbf{Check} : sujet singulier (\eng{Mary}) donc \eng{spends} avec -s ; chiffre conforme au texte ; phrase complète et courte. Correct.
\end{enumerate}
\end{exoresolu}

\courssub{Answering comprehension questions: rules and pitfalls}

\begin{propriete}[How to build a good answer]
Une bonne réponse de compréhension écrite est une \cle{phrase complète} construite en trois temps :
\begin{enumerate}
\item reprendre le \cle{sujet} de la question (\eng{How much does Mary spend...?} devient \eng{Mary spends...}) ;
\item insérer l'\cle{information} trouvée par scanning (\eng{2,000 francs}) ;
\item terminer par le \cle{complément} de la question (\eng{on vegetables}).
\end{enumerate}
La réponse doit être \textbf{courte} (une seule phrase suffit généralement) et \textbf{précise} (ni plus, ni moins que ce que demande la question).
\end{propriete}

\begin{center}
\begin{tabularx}{\linewidth}{|X|X|X|}
\hline
\rowcolor{popblueL} Question & Réponse correcte & Réponse à éviter \\
\hline
\eng{How much does Mary spend on vegetables?} & \eng{Mary spends 2,000 francs on vegetables.} & \eng{2,000} (trop fragmentaire) ou recopie de tout le paragraphe \\
\hline
\eng{Where does Mary keep her savings?} & \eng{She keeps her savings in a box at home.} & \eng{In her savings box at home she keeps the remaining 500 francs every Saturday at the market...} \\
\hline
\eng{What does Mary buy with her savings?} & \eng{She buys exercise books for school.} & \eng{Exercise books} (phrase incomplète) \\
\hline
\end{tabularx}
\end{center}

\begin{attention}[Ne recopiez jamais tout un paragraphe]
Beaucoup d'élèves, par peur de se tromper, recopient trois ou quatre lignes du texte comme réponse. C'est une \cle{erreur grave} : le correcteur ne peut pas savoir si vous avez vraiment compris, et il peut sanctionner la réponse. La bonne réponse est \textbf{courte et précise} : une phrase qui contient exactement l'information demandée, rien de plus. Mauvais : recopier \eng{Every Saturday, she goes to the market with a small budget of 5,000 francs. She spends 2,000 francs on vegetables...} Bon : \eng{Mary spends 2,000 francs on vegetables.}
\end{attention}

\begin{attention}[Accordez le verbe avec le sujet]
L'erreur la plus fréquente des francophones est l'oubli du \textbf{-s} de la troisième personne du singulier au présent simple. On écrit \eng{Mary spends}, jamais \eng{Mary spend}. De même : \eng{She keeps her savings...}, \eng{He earns money...}, \eng{It costs 500 francs}. Rappel : au présent simple, avec \eng{he}, \eng{she}, \eng{it} ou un nom singulier, le verbe prend un \textbf{-s} (ou \textbf{-es} après -s, -sh, -ch, -o : \eng{she watches}, \eng{he goes}). À l'étape \eng{check}, vérifiez systématiquement cet accord : c'est le premier détail que le correcteur remarque.
\end{attention}

\begin{exemplebox}[From question to answer — more practice]
\begin{itemize}
\item \eng{How much does Mary spend on fish?} → \eng{She spends 1,500 francs on fish.} « Elle dépense 1 500 francs en poisson. »
\item \eng{When does Mary go to the market?} → \eng{She goes to the market every Saturday.} « Elle va au marché chaque samedi. »
\item \eng{Why does managing her money help Mary?} → \eng{Managing her money helps her to avoid waste.} « Gérer son argent l'aide à éviter le gaspillage. »
\item \eng{Where does Mary live?} → \eng{She lives in Bamenda.} « Elle habite à Bamenda. »
\end{itemize}
Observez : chaque réponse reprend le sujet, conjugue le verbe correctement et apporte uniquement l'information demandée.
\end{exemplebox}

\begin{savaistu}[Les correcteurs du Bac acceptent la reformulation]
Bonne nouvelle : au baccalauréat, les correcteurs \cle{acceptent les réponses reformulées} avec vos propres mots, à condition que le sens soit correct et que l'anglais soit acceptable. Vous n'êtes donc pas obligé de recopier la phrase exacte du texte. Par exemple, pour \eng{Why does managing her money help Mary?}, les réponses \eng{Managing her money helps her to avoid waste}, \eng{Because it helps her not to waste money} et \eng{Because she does not waste money thanks to managing her money} sont toutes les trois valables. Reformuler montre même au correcteur que vous avez réellement compris le texte : c'est valorisé.
\end{savaistu}

\begin{exercice}[Your turn — apply the five steps]
En appliquant les cinq étapes au texte \eng{Mary's market budget}, répondez en phrases complètes :
\begin{enumerate}
\item \eng{What is Mary's total budget every Saturday?}
\item \eng{How much money does Mary save each week?}
\item \eng{What does Mary buy at the end of each month?}
\end{enumerate}
\end{exercice}

\begin{corrige}[Exercice : Your turn]
\begin{enumerate}
\item \eng{Mary's total budget is 5,000 francs every Saturday.} (Scanning : \eng{a small budget of 5,000 francs}.)
\item \eng{She saves 500 francs each week.} (Scanning : \eng{the remaining 500 francs}. Attention : \eng{saves} avec -s.)
\item \eng{At the end of each month, she buys exercise books for school.} (Scanning : dernière phrase du texte.)
\end{enumerate}
Chaque réponse est courte, complète, et le verbe est accordé avec le sujet.
\end{corrige}

\begin{aretenir}[L'essentiel de la méthode]
\begin{itemize}
\item Cinq étapes, dans l'ordre : \eng{Predict} → \eng{Skim} → \eng{Scan} → \eng{Answer} → \eng{Check}.
\item À l'examen, lisez les \cle{questions avant le texte} pour cibler votre lecture.
\item \eng{Skimming} = idée générale ; \eng{scanning} = information précise (chiffre, nom, date).
\item Réponse = phrase \textbf{courte, complète et précise} ; jamais de recopie d'un paragraphe entier.
\item À la relecture, vérifiez l'accord sujet-verbe : \eng{Mary spends}, pas \eng{Mary spend}.
\item La reformulation est acceptée si le sens est juste et l'anglais correct.
\end{itemize}
\end{aretenir}

\coursec{Guided production: summary and opinion}

Dans la section précédente, nous avons appris une méthode rigoureuse pour lire un texte sur les finances personnelles : le \cle{skimming} pour saisir l'idée générale, le \cle{scanning} pour repérer une information précise, et la déduction du sens des mots grâce au contexte. Nous allons maintenant passer de la \emph{réception} à la \emph{production} : résumer le texte avec nos propres mots, donner notre opinion, puis rédiger un court texte personnel intitulé \eng{My monthly budget}. C'est la dernière étape de la leçon, celle où vous montrez que vous avez vraiment compris.

\courssub{How to summarize the text in three sentences}

\begin{methode}[Résumer un texte en trois phrases]
\begin{enumerate}
\item \textbf{Relis le titre et la première phrase de chaque paragraphe} : c'est là que se trouvent les idées principales.
\item \textbf{Retiens une seule idée par paragraphe} : ne recopie pas les exemples ni les détails.
\item \textbf{Reformule avec tes propres mots} : change le vocabulaire et la structure des phrases (par exemple, transforme une phrase active en phrase plus simple).
\item \textbf{Utilise la trame} : \eng{The text is about...} (idée générale) / \eng{First, the writer explains that...} (première idée) / \eng{Finally, he or she advises...} (dernière idée).
\item \textbf{Ne donne jamais ton opinion dans un résumé} : un résumé rapporte les idées de l'auteur, pas les tiennes. Pas de \eng{I think} dans un résumé !
\end{enumerate}
\end{methode}

\begin{exemplebox}[Modèle de résumé en trois phrases]
\eng{The text is about managing personal finances. First, the writer explains that a budget helps us to control our income and our expenses. Finally, he advises young people to save a little money every month.}

« Le texte parle de la gestion des finances personnelles. D'abord, l'auteur explique qu'un budget nous aide à contrôler nos revenus et nos dépenses. Enfin, il conseille aux jeunes d'épargner un peu d'argent chaque mois. »
\end{exemplebox}

\begin{attention}[Résumé ou opinion ?]
Beaucoup d'élèves mélangent les deux. Retiens bien :
\begin{itemize}
\item Dans un \cle{résumé} : \eng{The writer says that...}, \eng{According to the text...} — jamais \eng{I think}.
\item Dans une \cle{opinion} : \eng{I think that...}, \eng{In my opinion...} — jamais \eng{The writer says}.
\end{itemize}
\end{attention}

\courssub{Giving your opinion: useful expressions}

Observe d'abord ces exemples, tirés de discussions entre élèves :

\begin{exemplebox}[Expressions d'opinion — exemples traduits]
\eng{In my opinion, every student should have a budget.} = « À mon avis, chaque élève devrait avoir un budget. »

\eng{I agree that saving money is important.} = « Je suis d'accord qu'épargner de l'argent est important. »

\eng{I believe that young people spend too much money on phone credit.} = « Je crois que les jeunes dépensent trop d'argent en crédit téléphonique. »

\eng{For me, a tontine is a good way to save money.} = « Pour moi, la tontine est un bon moyen d'épargner. »

\eng{It is important to plan your expenses before the end of the month.} = « Il est important de planifier ses dépenses avant la fin du mois. »
\end{exemplebox}

\begin{center}
\begin{tabularx}{\linewidth}{|X|X|X|}\hline
\rowcolor{popblueL} \textbf{English} & \textbf{Français} & \textbf{Emploi} \\ \hline
I think (that)... & Je pense que... & opinion neutre, très courant \\ \hline
In my opinion,... & À mon avis,... & opinion un peu plus formelle \\ \hline
I believe (that)... & Je crois que... & conviction forte \\ \hline
For me,... & Pour moi,... & point de vue personnel \\ \hline
I agree (that)... / I agree with... & Je suis d'accord que... / avec... & accord \\ \hline
I disagree with... & Je ne suis pas d'accord avec... & désaccord \\ \hline
It is important to... & Il est important de... & conseil général \\ \hline
In my view,... & Selon moi,... & formel, à l'écrit \\ \hline
\end{tabularx}
\end{center}

\begin{attention}[Le piège n° 1 : « I am agree »]
Ne dis \textbf{jamais} \eng{I am agree}. En anglais, \cle{agree} est un \textbf{verbe}, pas un adjectif : on dit simplement \eng{I agree} (« je suis d'accord »). De même : \eng{I disagree} (et non \eng{I am disagree}). Le français « je suis d'accord » utilise l'adjectif « d'accord », d'où l'erreur des francophones qui traduisent mot à mot.

Autre piège : après \eng{I think that}, on n'utilise pas le subjonctif comme en français. \eng{I think that saving is important} (et non une forme bizarre imitant « je pense qu'épargner \emph{soit} important »).
\end{attention}

\begin{propriete}[Conseiller avec \cle{should} + base verbale]
Pour donner un conseil, on utilise l'auxiliaire modal \cle{should} suivi de la \textbf{base verbale} (l'infinitif sans \eng{to}) :

\eng{You should save some money every month.} = « Tu devrais épargner de l'argent chaque mois. »

\begin{itemize}
\item \textbf{Forme} : sujet + should + base verbale. Should est \textbf{invariable} : \eng{he should save} (jamais \eng{he shoulds}).
\item \textbf{Négation} : should not / shouldn't + base verbale : \eng{You shouldn't spend all your pocket money.} = « Tu ne devrais pas dépenser tout ton argent de poche. »
\item \textbf{Question} : Should + sujet + base verbale ? \eng{Should I open a savings account?} = « Devrais-je ouvrir un compte d'épargne ? »
\item Pas de \eng{to} après should : on dit \eng{you should save}, jamais \eng{you should to save}.
\end{itemize}
\end{propriete}

\courssub{Guided writing: “My monthly budget”}

Avant de rédiger, lis ce modèle de production écrit par une élève de Première à Yaoundé.

\begin{textebox}[Model paragraph — “My monthly budget”, by Amina, Form 5]
Every month, I receive about 10,000 CFA francs from my parents. It is my pocket money, and sometimes I earn a little more by helping my aunt in her shop at Mokolo market. First, I pay for my real needs: my transport to school by taxi or bendskin, my lunch, and my exercise books. These expenses cost me about 6,000 CFA francs. Then I have my wants. I like buying phone credit and sometimes a snack with my friends after school. I try to spend only 2,000 CFA francs on these things, because they are not necessary. Finally, I save the rest of my money, about 2,000 CFA francs, in a small box at home. At the end of the year, I will use my savings to buy a new school bag and some clothes. In my opinion, every student should have a simple budget like mine. When you write down your income and your expenses, you control your money; your money does not control you. I believe that saving money, even a small amount, is a very good habit for the future.
\end{textebox}

\textbf{Glossaire} : \emph{pocket money} (n.) argent de poche ; \emph{to earn} (v.) gagner ; \emph{bendskin} (n.) moto-taxi ; \emph{want} (n.) envie, désir (chose non essentielle) ; \emph{snack} (n.) casse-croûte ; \emph{to save} (v.) épargner ; \emph{savings} (n. plur.) économies ; \emph{amount} (n.) somme, montant ; \emph{habit} (n.) habitude.

\begin{methode}[Rédiger “My monthly budget” en 6 à 8 lignes]
Suis le plan imposé en quatre étapes :
\begin{enumerate}
\item \textbf{Mes revenus} (Income) : \eng{Every month, I receive... / I earn... by...}
\item \textbf{Mes besoins} (Needs) : \eng{First, I pay for my needs: transport, food, school materials...}
\item \textbf{Mes envies} (Wants) : \eng{Then, I sometimes buy... because I like...}
\item \textbf{Mon épargne} (Savings) : \eng{Finally, I save... CFA francs. I will use my savings to...}
\end{enumerate}
Termine par une phrase d'opinion : \eng{In my opinion, every student should...}
\end{methode}

\begin{popfigure}
\begin{center}
\begin{tikzpicture}[node distance=0.5cm]
\node[draw=popblue, fill=popblueL, rounded corners, minimum width=2.7cm, minimum height=1.5cm, align=center] (inc) {\textbf{1. Income}\\ revenus};
\node[draw=popgreen, fill=popgreenL, rounded corners, minimum width=2.7cm, minimum height=1.5cm, align=center, right=of inc] (need) {\textbf{2. Needs}\\ besoins};
\node[draw=poporange, fill=poporangeL, rounded corners, minimum width=2.7cm, minimum height=1.5cm, align=center, right=of need] (want) {\textbf{3. Wants}\\ envies};
\node[draw=poppurple, fill=poppurpleL, rounded corners, minimum width=2.7cm, minimum height=1.5cm, align=center, right=of want] (save) {\textbf{4. Savings}\\ épargne};
\draw[-{Stealth}, very thick, popdark] (inc) -- (need);
\draw[-{Stealth}, very thick, popdark] (need) -- (want);
\draw[-{Stealth}, very thick, popdark] (want) -- (save);
\end{tikzpicture}
\end{center}
\legende{Le plan de rédaction en quatre cases : de tes revenus vers ton épargne.}
\end{popfigure}

\begin{exoresolu}[Rédaction guidée corrigée]
\textbf{Énoncé.} Rédige un paragraphe de 6 à 8 lignes intitulé \eng{My monthly budget} en respectant le plan : revenus $\rightarrow$ besoins $\rightarrow$ envies $\rightarrow$ épargne. Utilise au moins une expression d'opinion et un conseil avec \eng{should}.
\tcblower
\textbf{Solution détaillée (modèle de copie).}

\eng{Every month, I receive 8,000 CFA francs from my uncle in Douala, and I earn 2,000 CFA francs by selling groundnuts after school. First, I pay for my needs: my transport, my lunch and my pens. It costs about 5,000 CFA francs. Then, I spend 2,000 CFA francs on my wants, like phone credit and biscuits. Finally, I save 3,000 CFA francs in our class savings club. I will use my savings to buy textbooks next year. In my opinion, a budget is very useful. I think every student should save some money every month.}

Points de réussite : le plan en quatre étapes est respecté ; le lexique de la leçon est présent (\eng{receive, earn, needs, wants, save, savings}) ; une opinion (\eng{In my opinion}) et un conseil avec \eng{should} + base verbale (\eng{should save}) terminent le texte.
\end{exoresolu}

\begin{savaistu}[La tontine : a savings club]
La \cle{tontine} est une pratique d'épargne très courante au Cameroun : les membres d'un groupe cotisent régulièrement, et chacun reçoit le total à tour de rôle. En anglais, on peut la décrire comme \eng{a rotating savings group} ou plus simplement \eng{a savings club}. Exemple : \eng{My mother belongs to a savings club with her colleagues; every month, one member receives all the money.} = « Ma mère fait partie d'une tontine avec ses collègues ; chaque mois, un membre reçoit tout l'argent. »
\end{savaistu}

\courssub{Peer correction with a simple grid}

\begin{experience}[Correction par les pairs]
Échange ta production avec ton voisin et corrige-la avec cette grille simple (1 point par critère) :
\begin{enumerate}
\item \textbf{Plan} : les quatre étapes (income, needs, wants, savings) sont-elles présentes et dans l'ordre ?
\item \textbf{Lexique} : au moins cinq mots de la leçon (\eng{budget, income, expenses, save, savings, pocket money...}) sont-ils utilisés correctement ?
\item \textbf{Grammaire} : verbes conjugués au présent simple, \eng{should} + base verbale, pas de \eng{I am agree} ?
\item \textbf{Opinion} : le texte se termine-t-il par une phrase d'opinion claire ?
\item \textbf{Longueur} : 6 à 8 lignes ?
\end{enumerate}
Rends la copie avec une remarque positive et un conseil d'amélioration, en anglais : \eng{Good vocabulary! Check the verb “to save”.}
\end{experience}

\begin{exercice}[À toi de jouer]
\begin{enumerate}
\item Résume le texte modèle d'Amina en trois phrases avec la trame \eng{The text is about... / First... / Finally...}.
\item Donne ton opinion sur cette affirmation en 3 ou 4 phrases : \eng{Students should not receive pocket money.} Utilise \eng{I agree} ou \eng{I disagree} et \eng{because}.
\item Rédige ton propre texte \eng{My monthly budget} (6 à 8 lignes) en suivant le plan imposé.
\end{enumerate}
\end{exercice}

\begin{corrige}[Exercice — Guided production]
\textbf{1.} \eng{The text is about Amina's monthly budget. First, she explains her income and her needs, like transport and food. Finally, she saves 2,000 CFA francs every month and advises students to have a budget.}

\textbf{2.} Exemple : \eng{I disagree with this idea because pocket money teaches us to manage our finances. In my opinion, students should receive a small amount every month. It is important to learn how to save money when we are young.}

\textbf{3.} Vérification avec la grille de correction par les pairs : plan respecté, lexique de la leçon, présent simple correct, \eng{should} + base verbale, phrase d'opinion finale.
\end{corrige}

\begin{aretenir}[L'essentiel de la section]
\begin{itemize}
\item Un \cle{résumé} rapporte les idées de l'auteur en trois phrases, sans opinion personnelle.
\item Une \cle{opinion} s'exprime avec \eng{I think}, \eng{In my opinion}, \eng{I believe that}, \eng{For me}.
\item On dit \eng{I agree}, jamais \eng{I am agree}.
\item Le conseil : \cle{should} + base verbale (\eng{You should save money.}).
\item La rédaction \eng{My monthly budget} suit toujours le plan : income $\rightarrow$ needs $\rightarrow$ wants $\rightarrow$ savings.
\end{itemize}
\end{aretenir}

\newpage

\coursec{Activité d'intégration}

\coursec{Lesson 20 — Reading: Texts on managing personal finances}

\courssub{Objectifs de l'activité}

Cette activité d'intégration permet à l'élève de :
\begin{itemize}
\item \cle{réinvestir} le lexique financier étudié dans la leçon (\eng{budget, save, income, expenses, needs, wants, debt, interest, afford}) dans une production écrite personnelle et créative ;
\item \cle{mobiliser} les stratégies de lecture (\eng{skimming} et \eng{scanning}) pour exploiter rapidement un document d'information ;
\item \cle{utiliser} les auxiliaires modaux \eng{should} et \eng{must} pour donner des conseils et exprimer la recommandation ;
\item \cle{produire} une affiche collective en anglais correct et la présenter oralement en deux minutes ;
\item \cle{interagir} en posant et en répondant à des questions de compréhension en anglais.
\end{itemize}

\courssub{Situation}

\textbf{Contexte.} Le club d'anglais (\eng{English Club}) du Lycée de Nkongsamba prépare la \emph{Financial Literacy Week}, une semaine dédiée à la gestion de l'argent de poche. Beaucoup d'élèves dépensent tout leur argent de poche (\eng{pocket money}) en quelques jours : beignets, crédit téléphonique, transport en \emph{bendskin}, jeux en ligne. Le directeur du club a demandé à chaque groupe de quatre élèves de créer une affiche intitulée \emph{“Smart Money Tips for Students”} pour aider les élèves à mieux gérer leur argent. La meilleure affiche sera affichée au mur de la salle et publiée sur le panneau du club.

Pour vous inspirer, le président du club a affiché ce petit texte d'information :

\begin{textebox}[The 50/30/20 rule — a simple way to budget]
Financial advisers often recommend a simple rule to manage your money: the 50/30/20 rule. It divides your income into three parts.

First, spend about 50\% of your money on your \textbf{needs}: things you cannot live without, such as food, transport to school and school materials. Second, use about 30\% for your \textbf{wants}: things you enjoy but do not really need, like snacks, phone credit or entertainment. Finally, \textbf{save} 20\% of your income. Put this money aside for future projects or emergencies.

Many students say they have nothing to save, but even 100 francs a week becomes more than 5,000 francs in a year. The secret is simple: save first, spend after. If you borrow money you cannot repay, you fall into \textbf{debt}, and debt grows because of \textbf{interest}. Smart students plan a \textbf{budget}, write down their \textbf{expenses}, and never spend more than they earn.
\end{textebox}

\textbf{Glossaire :} \eng{advisers} (conseillers), \eng{to divide} (diviser), \eng{to put aside} (mettre de côté), \eng{emergencies} (urgences), \eng{to borrow} (emprunter), \eng{to repay} (rembourser), \eng{to grow} (augmenter).

\courssub{Consignes}

\begin{enumerate}
\item \textbf{Read and understand.} En groupe, lisez le texte \emph{The 50/30/20 rule}. D'abord, faites un \eng{skimming} (lecture rapide) et répondez en une phrase : \emph{What is the main idea of the text?} Ensuite, faites un \eng{scanning} (recherche d'informations précises) et complétez : \emph{50\% = \ldots ; 30\% = \ldots ; 20\% = \ldots}
\item \textbf{Find the vocabulary.} Relevez dans le texte cinq mots ou expressions du lexique financier et donnez leur équivalent français. Déduisez le sens de \eng{to fall into debt} à partir du contexte, sans dictionnaire.
\item \textbf{Write your five tips.} Rédigez cinq conseils en anglais pour les élèves du lycée. Utilisez \eng{should} pour les recommandations et \eng{must} pour les règles fortes. Chaque conseil doit contenir au moins un mot du lexique de la leçon (\eng{save, budget, needs, wants, avoid debt, expenses, afford}). Exemple de structure : \eng{You should save 20\% of your pocket money every week.}
\item \textbf{Draw the pie chart.} Illustrez la règle 50/30/20 par un camembert (\eng{pie chart}) avec trois parts étiquetées en anglais : \eng{Needs 50\%}, \eng{Wants 30\%}, \eng{Savings 20\%}.
\item \textbf{Design the poster.} Composez l'affiche : titre \emph{“Smart Money Tips for Students”}, les cinq conseils numérotés, le camembert, et une phrase d'accroche finale (\eng{slogan}), par exemple : \eng{Save first, spend after!}
\item \textbf{Present your poster.} Chaque groupe présente son affiche à la classe en deux minutes : chaque membre lit au moins un conseil. Les autres groupes posent deux questions de compréhension, par exemple : \eng{Why should we save first? What is the difference between a need and a want?}
\item \textbf{Vote.} La classe vote pour la meilleure affiche (clarté de l'anglais, exactitude du lexique, qualité de la présentation). L'affiche gagnante est affichée au mur de la salle.
\end{enumerate}

\courssub{Ressources à mobiliser}

\begin{propriete}[Should et must pour conseiller]
\begin{itemize}
\item \eng{should} + base verbale : conseil, recommandation. \eng{You should write down your expenses.} « Tu devrais noter tes dépenses. » Négation : \eng{shouldn't} (\eng{You shouldn't borrow money for snacks.}).
\item \eng{must} + base verbale : obligation forte, règle. \eng{You must never spend more than you earn.} « Tu ne dois jamais dépenser plus que tu ne gagnes. » Négation (interdiction) : \eng{must not} / \eng{mustn't} (\eng{You mustn't spend money you don't have.}).
\item Après \eng{should} et \eng{must}, le verbe reste à la base verbale, sans \eng{to} et sans \eng{-s} à la troisième personne : \eng{He should save}, et non \emph{he should saves}.
\end{itemize}
\end{propriete}

\begin{center}
\begin{tabularx}{\linewidth}{|X|X|X|}\hline
\rowcolor{popblueL} English & Français & Exemple utile \\ \hline
income & revenu & \eng{My income is my pocket money.} \\ \hline
expenses & dépenses & \eng{Write down all your expenses.} \\ \hline
needs & besoins & \eng{Food and transport are needs.} \\ \hline
wants & envies & \eng{Phone credit is a want, not a need.} \\ \hline
to save / savings & épargner / épargne & \eng{Save 20\% every week.} \\ \hline
budget & budget & \eng{Plan a weekly budget.} \\ \hline
debt / to avoid debt & dette / éviter les dettes & \eng{Avoid debt: do not borrow for wants.} \\ \hline
interest & intérêts & \eng{Debt grows because of interest.} \\ \hline
to afford & avoir les moyens de & \eng{I cannot afford new shoes this month.} \\ \hline
\end{tabularx}
\end{center}

\begin{attention}[Erreurs fréquentes des francophones]
\begin{itemize}
\item Ne dites pas \emph{to make economies} : le faux ami « économies » se dit \eng{savings} ou \eng{to save money}.
\item \eng{money} est indénombrable : \eng{money is}, jamais \emph{monies} ici ; on dit \eng{some money}, pas \emph{a money}.
\item Pas de \eng{to} après les modaux : \eng{you should save}, pas \emph{you should to save}.
\item \eng{to afford} signifie « avoir les moyens de », pas « se permettre » au sens de « s'autoriser ».
\end{itemize}
\end{attention}

\courssub{Méthode : Stratégies de lecture — Skimming et Scanning}

\begin{methode}[Comment lire efficacement un texte d'information]
\begin{enumerate}
\item \textbf{Skimming (lecture rapide pour l'idée générale).} Lisez le titre, la première et la dernière phrase de chaque paragraphe, repérez les mots-clés en gras ou en italique. Ne lisez pas chaque mot. Objectif : répondre à \emph{What is the main idea?}
\item \textbf{Scanning (recherche ciblée pour des détails précis).} Balayez le texte des yeux à la recherche de chiffres, pourcentages, noms propres, dates ou mots-clés précis (ex. \eng{50\%}, \eng{debt}, \eng{interest}). Objectif : compléter un tableau ou répondre à des questions factuelles.
\item \textbf{Deviner le sens des mots inconnus.} Identifiez la nature du mot (verbe, nom, adjectif) grâce à sa position. Regardez le contexte immédiat (mots avant/après, connecteurs logiques comme \eng{because}, \eng{but}, \eng{for example}). Proposez un sens cohérent en français.
\end{enumerate}
\end{methode}

\courssub{Production guidée et corrigé commenté}

\begin{exoresolu}[My Smart Budget Poster — modèle de production]
Rédigez le contenu de l'affiche \emph{“Smart Money Tips for Students”} : cinq conseils avec \eng{should}/\eng{must}, le camembert 50/30/20 et un slogan final.
\tcblower
\textbf{Smart Money Tips for Students}

\begin{enumerate}
\item \eng{You should plan a budget every week and write down your expenses.}
\item \eng{You must pay for your needs first: food, transport and school materials.}
\item \eng{You should save 20\% of your pocket money before you spend anything.}
\item \eng{You shouldn't borrow money for your wants: avoid debt and interest.}
\item \eng{You must never spend more money than you earn.}
\end{enumerate}

\eng{Pie chart: Needs 50\% — Wants 30\% — Savings 20\%.}

\eng{Slogan: “Save first, spend after — smart students, smart money!”}

\textbf{Commentaire des choix de langue.}
\begin{itemize}
\item \textbf{Modaux bien différenciés} : \eng{should} et \eng{shouldn't} pour les recommandations (conseils 1, 3, 4), \eng{must} pour les règles fortes (conseils 2 et 5). Le verbe qui suit est toujours à la base verbale (\eng{plan, pay, save, borrow, spend}).
\item \textbf{Lexique de la leçon réinvesti} : \eng{budget, expenses, needs, wants, save, pocket money, borrow, debt, interest, earn} — dix mots du champ lexical financier en cinq phrases.
\item \textbf{Structures simples et correctes} (niveau A2+/B1) : phrases courtes au présent, connecteur \eng{before} + proposition (\eng{before you spend anything}).
\item \textbf{Attention aux pièges évités} : \eng{20\% of your pocket money} (pas \emph{of pocket money} sans déterminant), \eng{more money than you earn} (comparatif correct avec \eng{than}).
\item \textbf{Slogan} : parallélisme \eng{Save first, spend after} facile à mémoriser, et jeu de mots \eng{smart students, smart money} qui reprend le titre de l'affiche.
\end{itemize}
\end{exoresolu}

\courssub{À retenir}

\begin{aretenir}[Synthèse lexicale et grammaticale — Lesson 20]
\begin{itemize}
\item \textbf{Vocabulaire clé} : \eng{income, expenses, needs, wants, budget, to save / savings, debt, interest, to afford, pocket money}.
\item \textbf{Distinction \eng{needs} vs \eng{wants}} : \eng{needs} = indispensable (survie, scolarité) ; \eng{wants} = plaisir, superflu.
\item \textbf{Grammaire} : \eng{should/shouldn't} (conseil) vs \eng{must/mustn't} (obligation/interdiction) + \textbf{base verbale sans \eng{to}}.
\item \textbf{Règle 50/30/20} : 50\% \eng{needs} — 30\% \eng{wants} — 20\% \eng{savings}.
\item \textbf{Stratégies} : \eng{skimming} (idée générale), \eng{scanning} (info précise), déduction contextuelle.
\end{itemize}
\end{aretenir}

\courssub{Bilan de l'activité}

À l'issue de cette activité, l'élève a :
\begin{itemize}
\item lu et exploité un texte d'information financière grâce au \eng{skimming} (idée générale) et au \eng{scanning} (chiffres de la règle 50/30/20) ;
\item déduit le sens de mots nouveaux à partir du contexte (\eng{to fall into debt}, \eng{to put aside}) ;
\item réutilisé activement au moins dix mots du lexique financier dans une production écrite authentique et créative ;
\item employé correctement \eng{should} et \eng{must} pour formuler des conseils et des règles ;
\item pris la parole en public pendant deux minutes et interagi avec ses camarades en posant et en répondant à des questions en anglais.
\end{itemize}

\begin{experience}[Pour aller plus loin : The Pocket Money Challenge]
En prolongement, chaque élève tient un journal de ses dépenses en anglais pendant une semaine (\eng{My Expenses Diary}) : date, objet acheté, montant, catégorie \eng{need} ou \eng{want}. En fin de semaine, chacun écrit trois phrases de bilan : \eng{This week, I spent \ldots on needs and \ldots on wants. Next week, I should \ldots} Le groupe compare les résultats et désigne le \emph{best saver} de la classe.
\end{experience}

\begin{savaistu}[Did you know?]
La règle 50/30/20 a été popularisée aux États-Unis dans les années 2000 et elle est aujourd'hui enseignée dans de nombreuses écoles anglophones dans le cadre de la \eng{financial literacy} (littératie financière), c'est-à-dire l'éducation à la gestion de l'argent. Au Cameroun, les tontines jouent un rôle comparable : chaque membre cotise régulièrement et reçoit à son tour une somme qui lui permet de financer un projet. Dans les deux cas, le principe est le même : \eng{save regularly, spend wisely!} (« épargne régulièrement, dépense sagement ! »).
\end{savaistu}

\newpage

\coursec{Méthodes et exercices résolus}

\coursec{Lesson 20 — Reading: Managing Personal Finances}

\courssub{Warm-up and Pre-reading}

\begin{experience}[Discussion : Argent de poche et premiers revenus]
\begin{enumerate}
\item En binôme, discutez en anglais : \eng{How do you get money? (pocket money, small jobs, gifts)} — \eng{What do you spend it on? (phone credit, snacks, clothes, school supplies)} — \eng{Do you save any? How?}
\item L'enseignant note au tableau les mots-clés qui émergent : \eng{income, expenses, savings, budget, need, want, emergency}.
\item Objectif : activer le schéma mental (\eng{schemata}) et le lexique du thème avant la lecture.
\end{enumerate}
\end{experience}

\begin{methode}[Stratégies avant lecture (Pre-reading)]
\begin{enumerate}
\item \textbf{Observez le titre, la source, les images} : ils orientent le thème (conseils bancaires, témoignage d'élève, article de blog).
\item \textbf{Prédisez le contenu} : à partir des mots du titre, écrivez 3 idées que vous attendez de trouver dans le texte.
\item \textbf{Repérez la structure} : nombre de paragraphes, présence de sous-titres, de listes à puces (\eng{bullet points}), de chiffres (dates, pourcentages).
\item \textbf{Fixez un but de lecture} : \eng{“I read to find three tips for saving money.”}
\end{enumerate}
\end{methode}

\courssub{Reading Text: “Smart Money Habits for Young Cameroonians”}

\begin{textebox}[Article original — OnBuch+ / Contexte camerounais]
\eng{Smart Money Habits for Young Cameroonians}

\eng{Many secondary school students in Yaoundé, Douala or Buea receive pocket money from their parents or earn a little cash during holidays. Yet, by the middle of the month, most of them wonder where their money went. Learning to manage personal finances early is not only for adults; it is a survival skill.}

\eng{The first step is to track every single franc. Write down your income (pocket money, gifts, small wages) and your expenses (phone credit, snacks, transport, school materials). You can use a simple notebook, a spreadsheet on your phone, or a budgeting app. “If you don’t tell your money where to go, you will wonder where it went,” says Jean-Paul, a financial coach in Douala.}

\eng{The second habit is to pay yourself first. Before buying airtime or new shoes, put aside at least 10\% of your income into a savings envelope or a Mobile Money savings wallet (like MoMo Savings or Orange Money Épargne). This emergency fund protects you when unexpected costs arise — a torn school uniform, a sudden illness, or a rise in transport fares.}

\eng{Third, distinguish between needs and wants. Needs are essential: food, rent (if you live alone), transport to school, medical care. Wants are pleasant but not vital: the latest video game, fancy clothes, daily soft drinks. A good rule is the 50/30/20 rule: 50\% for needs, 30\% for wants, 20\% for savings and debt repayment.}

\eng{Finally, avoid unnecessary debt. Borrowing from friends or using “buy now, pay later” offers for non-essential items traps you in a cycle of repayment. If you must borrow — for example, to buy a laptop for your studies — compare interest rates, read the contract carefully, and have a clear repayment plan.}

\eng{Managing money is a habit, not a talent. Start today with a small notebook. Your future self will thank you.}
\end{textebox}

\begin{definition}[Glossaire — Mots et expressions clés du texte]
\begin{itemize}
\item \eng{pocket money} (n.) : argent de poche.
\item \eng{earn} (v., irr. \eng{earned, earned}) : gagner (de l'argent par son travail).
\item \eng{cash} (n.) : espèces, argent liquide.
\item \eng{track} (v.) : suivre, noter, surveiller (\eng{track expenses}).
\item \eng{income} (n.) : revenus, rentrées d'argent.
\item \eng{expenses} (n. pl.) : dépenses, sorties d'argent.
\item \eng{spreadsheet} (n.) : tableau de calcul (Excel, Google Sheets).
\item \eng{pay yourself first} (expr.) : s payer en premier (épargner avant de dépenser).
\item \eng{set aside} (phr. v.) : mettre de côté, réserver.
\item \eng{emergency fund} (n.) : fonds d'urgence, épargne de précaution.
\item \eng{arise} (v., irr. \eng{arose, arisen}) : survenir, se présenter (problème, besoin).
\item \eng{distinguish} (v.) : faire la distinction, différencier.
\item \eng{needs vs. wants} : besoins vs désirs (envies).
\item \eng{debt repayment} (n.) : remboursement de dettes.
\item \eng{unnecessary debt} (n.) : dette inutile / superflue.
\item \eng{interest rate} (n.) : taux d'intérêt.
\item \eng{repayment plan} (n.) : plan de remboursement.
\end{itemize}
\end{definition}

\courssub{Reading Strategies: Skimming, Scanning, Guessing Meaning}

\begin{methode}[Lire efficacement : Skimming, Scanning, Inférence lexicale]
\begin{enumerate}
\item \textbf{Skimming (lecture rapide globale)} : Lisez le titre, l'introduction, la conclusion et la \textbf{première phrase de chaque paragraphe} (souvent la \eng{topic sentence}). Objectif : saisir l'\eng{main idea} (idée générale) en 1 à 2 minutes. Ne lisez pas chaque mot.
\item \textbf{Scanning (recherche ciblée)} : Pour une question précise (chiffre, nom, date, mot-clé), balayez le texte des yeux \textbf{sans lire linéairement}. Repérez le mot-clé ou ses synonymes.
\item \textbf{Inférence du sens (Guessing meaning from context)} : Ne sortez pas le dictionnaire tout de suite. Utilisez :
      \begin{itemize}
      \item La \textbf{définition explicite} : \eng{“that is to say”, “in other words”, “which means”}.
      \item L'\textbf{exemple} : \eng{“such as”, “for example”, “like”}.
      \item L'\textbf{opposition} : \eng{“but”, “however”, “unlike”, “whereas”, “on the contrary”}.
      \item La \textbf{cause/conséquence} : \eng{“because”, “so”, “therefore”, “as a result”}.
      \item La \textbf{famille de mots} : \eng{save (v.) $\rightarrow$ savings (n.), saver (n.)} ; \eng{spend $\rightarrow$ spending, spender}.
      \item Les \textbf{préfixes/suffixes} : \eng{un-necessary, ir-responsible, finance-ial, budget-ing}.
      \end{itemize}
\item \textbf{Mots-outils logiques (Connectors)} : Surlignez \eng{however (cependant), moreover (de plus), therefore (par conséquent), for instance (par exemple), in conclusion (en conclusion)}. Ils structurent l'argumentation.
\end{enumerate}
\end{methode}

\begin{exoresolu}[Application du Skimming et Scanning sur le texte “Smart Money Habits…”]
\textbf{Énoncé :} 1) Quel est le thème principal du texte ? (Skimming) 2) Quel pourcentage du revenu faut-il épargner selon le texte ? (Scanning) 3) Donnez le sens de \eng{“arise”} (paragraphe 3) grâce au contexte. 4) Quel connecteur introduit le dernier conseil du texte (paragraphe 5) ?
\tcblower
\textbf{1) Skimming :} Lecture des premières phrases : §1 (problème : argent disparu), §2 (étape 1 : tracker), §3 (étape 2 : payer soi-même), §4 (étape 3 : besoins/désirs), §5 (étape 4 : éviter dette), conclusion. \textbf{Réponse :} \eng{The text gives four smart habits to help young Cameroonians manage their money well.} (Le texte donne 4 habitudes pour bien gérer son argent.)

\textbf{2) Scanning :} Chercher chiffre / \%\ / \eng{percent} / \eng{10}. Trouvé §3 : \eng{“put aside at least 10\% of your income”}. \textbf{Réponse :} \eng{At least 10\% of your income.}

\textbf{3) Inférence \eng{arise} (§3) :} Contexte : \eng{“unexpected costs arise”}. \eng{Unexpected costs} = coûts imprévus. Idée de quelque chose qui arrive soudain. \eng{Arise} $\approx$ \eng{happen / occur}. \textbf{Réponse :} \eng{Arise means “happen” or “occur” (survenir).}

\textbf{4) Connecteur §5 :} Le paragraphe commence par \eng{“Finally,”}. \textbf{Réponse :} \eng{Finally} (en dernier lieu).
\end{exoresolu}

\courssub{Comprehension Activities}

\begin{methode}[Répondre aux questions de compréhension — Règles d'or]
\begin{enumerate}
\item \textbf{Type de question} : \eng{Wh-} (fait), \eng{True/False Justify} (citation obligatoire), \eng{In your own words} (reformulation), \eng{Vocabulary} (synonyme/contexte), \eng{Interpretation} (\eng{Why/How} + opinion étayée).
\item \textbf{Localisez} : Soulignez la zone de la réponse dans le texte avant d'écrire.
\item \textbf{Phrase complète} : Sujet + verbe + complément. Reprenez le sujet de la question. \eng{Question: Why does Jean-Paul recommend tracking expenses? $\rightarrow$ Answer: Jean-Paul recommends tracking expenses because...}
\item \textbf{Concordance des temps} : Question au \eng{Past Simple} $\rightarrow$ Réponse au \eng{Past Simple}. Question au \eng{Present Perfect} $\rightarrow$ Réponse au \eng{Present Perfect}.
\item \textbf{True/False} : \eng{True} ou \eng{False} + \textbf{Citation exacte} entre guillemets anglais “ ”. \eng{False. The text states: “Many students wonder where their money went.”}
\item \textbf{Reformulation (\eng{In your own words})} : Changez la structure (voix passive/active), remplacez le vocabulaire (synonymes), gardez le sens. Interdit de copier-coller.
\item \textbf{Relecture} : \eng{-s} à la 3e pers. sg (\eng{he saves}), majuscule \eng{I}, point final, orthographe \eng{money/advice/information} (invariables).
\end{enumerate}
\end{methode}

\begin{exoresolu}[Compréhension détaillée — Type Examen / Bac]
\textbf{Questions sur le texte “Smart Money Habits…” :}
\begin{enumerate}
\item \eng{True or False? Justify with a quote.} “Most students in Cameroon know exactly how they spend their pocket money.”
\item \eng{According to the text, what are the two tools mentioned to track expenses?}
\item \eng{Find in the text a synonym for “essential” (paragraph 4).}
\item \eng{In your own words, explain the “50/30/20 rule”.}
\item \eng{Why does the author advise against “buy now, pay later” for non-essential items?}
\end{enumerate}
\tcblower
\textbf{1) False.} \eng{The text says: “by the middle of the month, most of them wonder where their money went.”} (Ne pas savoir où l'argent est passé = ne pas savoir comment il a été dépensé).

\textbf{2) Factuelle (§2) :} \eng{A simple notebook, a spreadsheet on your phone, or a budgeting app.} (Accepter deux sur trois). \textbf{Réponse :} \eng{The text mentions a simple notebook, a spreadsheet on your phone, or a budgeting app.}

\textbf{3) Vocabulaire (§4) :} \eng{Needs are essential...} $\rightarrow$ \eng{essential} = \eng{vital} (donné dans la phrase : \eng{“Needs are essential: food... medical care. Wants are pleasant but not vital.”}). \textbf{Réponse :} \eng{Vital.}

\textbf{4) Reformulation (§4) :} \textbf{Réponse :} \eng{The rule suggests spending half of your income on necessary things, 30\% on things you like but don't need, and saving the remaining 20\% or using it to pay back debts.}

\textbf{5) Interprétation (§5) :} \textbf{Réponse :} \eng{Because it traps you in a cycle of repayment for things you don't really need, which creates financial stress.}
\end{exoresolu}

\begin{exercice}[Entraînement — Compréhension écrite]
\textbf{Relisez le texte et répondez en phrases complètes :}
\begin{enumerate}
\item Who is Jean-Paul and what is his advice?
\item What does “pay yourself first” mean practically for a student?
\item Find two examples of “unexpected costs” given in paragraph 3.
\item What is the risk of borrowing for non-essential items?
\item \eng{Critical thinking :} Do you think the 50/30/20 rule is realistic for a student in Bamenda? Why / Why not? (3 phrases min.)
\end{enumerate}
\end{exercice}

\courssub{Thematic Vocabulary: Personal Finance}

\begin{methode}[Apprendre le lexique par familles et collocations]
\begin{enumerate}
\item \textbf{Classez par champ sémantique} (voir tableau ci-dessous).
\item \textbf{Mémorisez les collocations} (associations mots fixes) : \eng{make a budget, track expenses, set aside money, pay school fees, run into debt, take out a loan, pay back / repay a loan, earn a living, save for a rainy day}.
\item \textbf{Attention aux paires pièges} :
      \begin{itemize}
      \item \eng{Lend (to sb) / Borrow (from sb)} : \eng{I \textbf{lend} my book \textbf{to} Paul.} / \eng{Paul \textbf{borrows} my book \textbf{from} me.}
      \item \eng{Salary (mensuel, cadre) / Wages (hebdo/horaire, ouvrier)}.
      \item \eng{Expense (dépense) / Expensive (adj. cher)}.
      \item \eng{Afford (avoir les moyens) / Effort}.
      \end{itemize}
\item \textbf{Prépositions obligatoires} : \eng{spend money \textbf{on}}, \eng{save \textbf{for}}, \eng{borrow \textbf{from}}, \eng{lend \textbf{to}}, \eng{pay \textbf{for}}, \eng{invest \textbf{in}}, \eng{charge \textbf{for}}.
\item \textbf{Noms indénombrables} : \eng{money, advice, information, cash, furniture, luggage} — jamais de \eng{-s}, pas de \eng{a/an}. \eng{Some money, a piece of advice}.
\end{enumerate}
\end{methode}

\begin{propriete}[Familles de mots et Tableau de conjugaison utiles]
\begin{center}
\begin{tabularx}{\linewidth}{|X|X|X|X|}
\hline
\rowcolor{popblueL} \textbf{Verbe} & \textbf{Nom (action/résultat)} & \textbf{Nom (agent)} & \textbf{Adjectif / Participe} \\
\hline
\eng{save} & \eng{savings} & \eng{saver} & \eng{saving (adj: \eng{a saving account})} \\
\hline
\eng{spend} (\eng{spent, spent}) & \eng{spending} & \eng{spender} & \eng{spending (adj: \eng{spending power})} \\
\hline
\eng{earn} & \eng{earnings} & \eng{earner} & \eng{earned (income)} \\
\hline
\eng{budget} & \eng{budget, budgeting} & --- & \eng{budgetary} \\
\hline
\eng{invest} & \eng{investment} & \eng{investor} & \eng{investment (adj)} \\
\hline
\eng{borrow} & \eng{loan (n.), borrowing} & \eng{borrower} & \eng{borrowed} \\
\hline
\eng{lend} (\eng{lent, lent}) & \eng{loan, lending} & \eng{lender} & \eng{lent} \\
\hline
\eng{owe} & \eng{debt} & \eng{debtor} & \eng{owing} \\
\hline
\eng{repay} (\eng{repaid, repaid}) & \eng{repayment} & --- & \eng{repayable} \\
\hline
\end{tabularx}
\end{center}

\textbf{Verbes irréguliers clés :} \eng{spend $\rightarrow$ spent $\rightarrow$ spent} | \eng{lend $\rightarrow$ lent $\rightarrow$ lent} | \eng{earn $\rightarrow$ earned $\rightarrow$ earned} (régulier) | \eng{pay $\rightarrow$ paid $\rightarrow$ paid} | \eng{owe $\rightarrow$ owed $\rightarrow$ owed} | \eng{arise $\rightarrow$ arose $\rightarrow$ arisen}.
\end{propriete}

\begin{exoresolu}[Vocabulary in Context — Complétion]
\textbf{Complétez avec le mot juste (forme correcte) :} \eng{afford — budget — borrow — expenses — income — lend — savings — track — debt — repay}
\begin{enumerate}
\item My monthly \ldots\ is 50,000 FCFA, but my \ldots\ are higher.
\item I cannot \ldots\ a new laptop right now; I must \ldots\ money each month.
\item If you \ldots\ money from a bank, you must \ldots\ it with interest.
\item Please \ldots\ me 1,000 FCFA for the bus; I will give it back tomorrow.
\item Wise students \ldots\ their spending daily to stick to their \ldots\.
\item He is in \ldots\ because he bought a motorbike on credit without a \ldots\ plan.
\end{enumerate}
\tcblower
\textbf{1)} \eng{income — expenses} (sujet pluriel \eng{are}).
\textbf{2)} \eng{afford — save} (\eng{cannot afford} + nom ; \eng{save money} collocation).
\textbf{3)} \eng{borrow — repay} (\eng{borrow from}, \eng{repay} = rembourser).
\textbf{4)} \eng{lend} (\eng{lend to sb} = prêter à qqn ; \eng{borrow from} = emprunter à qqn).
\textbf{5)} \eng{track — budget} (\eng{track expenses/spending}, \eng{stick to a budget}).
\textbf{6)} \eng{debt — repayment} (\eng{be in debt}, \eng{repayment plan}).
\end{exoresolu}

\begin{exercice}[Entraînement — Vocabulaire]
\begin{enumerate}
\item Traduisez en anglais : \eng{1. Je ne peux pas me permettre ce voyage. 2. Elle a emprunté de l'argent à sa sœur. 3. Nous devons rembourser le prêt. 4. Il suit ses dépenses sur une application. 5. Épargner est une bonne habitude.}
\item Complétez avec la préposition correcte : \eng{on / for / from / to / in / back}
      \begin{enumerate}
      \item She spends a lot of money \ldots\ clothes.
      \item I am saving \ldots\ a new phone.
      \item He borrowed 10,000 FCFA \ldots\ his brother.
      \item The bank lent money \ldots\ the company.
      \item You must pay \ldots\ the goods before leaving.
      \item Invest \ldots\ your education.
      \end{enumerate}
\item \textbf{Odd one out (Intrus)} : Circle the word that does not belong.
      \begin{enumerate}
      \item income — salary — wages — \textbf{expenses}
      \item lend — borrow — \textbf{afford} — loan
      \item thrifty — frugal — economical — \textbf{extravagant}
      \item debt — loan — credit — \textbf{savings}
      \end{enumerate}
\end{enumerate}
\end{exercice}

\courssub{Grammar Focus: Modals of Advice and Obligation}

\begin{propriete}[Modaux de conseil et d'obligation : Should / Ought to / Must / Have to]
\textbf{Contexte :} Dans les textes de conseils financiers (\eng{“You should save...”, “You must avoid debt...”}), ces modaux sont omniprésents.

\begin{center}
\begin{tabularx}{\linewidth}{|X|X|X|X|}
\hline
\rowcolor{popblueL} \textbf{Modal} & \textbf{Force / Nuance} & \textbf{Structure} & \textbf{Exemple (Finances)} \\
\hline
\eng{should} & Conseil fort, avis personnel & \eng{should + BV} & \eng{You \textbf{should} make a budget.} \\
\hline
\eng{ought to} & Conseil moral, devoir (plus formel) & \eng{ought to + BV} & \eng{Students \textbf{ought to} learn finance early.} \\
\hline
\eng{must} & Obligation interne, forte nécessité (locuteur) & \eng{must + BV} & \eng{I \textbf{must} pay my fees tomorrow.} \\
\hline
\eng{have to} & Obligation externe, règle, loi, contrainte & \eng{have/has to + BV} & \eng{You \textbf{have to} show ID to open an account.} \\
\hline
\eng{had to} & Passé de \eng{must/have to} (obligation passée) & \eng{had to + BV} & \eng{Last year, I \textbf{had to} borrow for books.} \\
\hline
\end{tabularx}
\end{center}

\textbf{Forme négative :}
\begin{itemize}
\item \eng{should not / shouldn't} = déconseil (\eng{You shouldn't spend all your cash.}).
\item \eng{must not / mustn't} = interdiction stricte (\eng{You mustn't share your PIN.}).
\item \eng{don't/doesn't have to} = \textbf{absence d'obligation} (pas « ne pas devoir » !) (\eng{You don't have to pay now = Vous n'êtes pas obligé de payer maintenant.}).
\item \eng{ought not to / oughtn't to} (rare, formel).
\end{itemize}

\textbf{Passé (regret / critique) :} \eng{should / ought to + have + V3 (participe passé)}.
\eng{You \textbf{should have saved} money. (Tu aurais dû épargner — mais tu ne l'as pas fait.)}
\eng{He \textbf{ought to have compared} interest rates.}
\end{propriete}

\begin{exemplebox}[Exemples authentiques]
\eng{“You \textbf{should} track every franc.”} (Conseil du texte) \\
\eng{“You \textbf{have to} read the contract before signing.”} (Obligation légale/banque) \\
\eng{“I \textbf{must} stop buying useless things.”} (Résolution personnelle) \\
\eng{“We \textbf{don't have to} be rich to start saving.”} (Absence d'obligation) \\
\eng{“You \textbf{mustn't} ignore your debts.”} (Interdiction) \\
\eng{“I \textbf{should have bought} the cheaper phone.”} (Regret passé)
\end{exemplebox}

\begin{exercice}[Grammaire — Modaux]
\textbf{Complétez avec :} \eng{should / shouldn't / must / mustn't / have to / don't have to / had to}
\begin{enumerate}
\item You \ldots\ create a budget if you want to control your money. (Conseil fort)
\item Students \ldots\ pay school fees at the start of the year. (Règle école)
\item You \ldots\ share your Mobile Money PIN with anyone. (Interdiction)
\item I \ldots\ work during holidays to pay for my books. (Obligation passée)
\item You \ldots\ be a millionaire to open a savings account. (Absence d'obligation)
\item \ldots\ I \ldots\ declare my small income to taxes? (Question obligation externe)
\end{enumerate}
\end{exercice}

\courssub{Writing and Speaking: Summary and Opinion}

\begin{methode}[Rédiger un résumé (Summary) et une opinion (Opinion Essay court)]
\begin{enumerate}
\item \textbf{Résumé (50--80 mots)} :
      \begin{itemize}
      \item Identifiez 3--4 idées forces (une par §).
      \item Reformulez : changez le vocabulaire (synonymes), la structure (actif $\leftrightarrow$ passif), combinez les phrases.
      \item Utilisez des connecteurs : \eng{First(ly), Secondly, Moreover, Furthermore, Finally, In short / To sum up}.
      \item Temps : \textbf{Present Simple} (vérités générales, idées du texte). Pas de \eng{I}, pas d'opinion.
      \item Comptez les mots (\eng{Word count}).
      \end{itemize}
\item \textbf{Opinion personnelle (3--5 phrases)} :
      \begin{itemize}
      \item Marqueur : \eng{In my opinion, / I believe that / To my mind, / As far as I am concerned, / I agree / disagree that...}
      \item Argument + Exemple personnel (\eng{For example, / For instance, / Take my case:}).
      \item Structure : \eng{Opinion $\rightarrow$ Reason $\rightarrow$ Example}.
      \end{itemize}
\end{enumerate}
\end{methode}

\begin{exoresolu}[Summary \& Opinion — Modèle Bac]
\textbf{Tâche :} \eng{Summarise the text “Smart Money Habits…” in about 60 words. Then give your opinion on the importance of saving for students.}
\tcblower
\textbf{Summary (62 words) :}
\eng{The article advises young Cameroonians to adopt four smart financial habits. First, tracking all income and expenses reveals where money goes. Second, paying yourself first by saving at least 10\% builds an emergency fund. Third, distinguishing needs from wants using the 50/30/20 rule balances spending. Finally, avoiding unnecessary debt prevents stressful repayment cycles. Developing these habits early ensures future financial stability.}

\textbf{Opinion :}
\eng{In my opinion, saving is crucial for students, even with small amounts. I believe that regular saving teaches discipline and prepares us for emergencies. For instance, last term I saved 500 FCFA weekly from my pocket money and could buy my physics textbook without asking my parents.}
\end{exoresolu}

\begin{experience}[Expression orale : Jeu de rôle — Conseiller financier / Étudiant]
\textbf{Consigne :} Par groupes de 3. Rôles : \eng{A: Financial Advisor (banque/MoMo), B: Student (revenu 15.000 FCFA/mois), C: Observer (grille d'évaluation).}
\textbf{Scénario :} L'étudiant veut acheter un smartphone à 120.000 FCFA (crédit possible). Le conseiller doit analyser son budget (besoins/désirs), proposer un plan d'épargne ou un crédit responsable, expliquer les taux.
\textbf{Durée :} 5 min préparation, 3 min jeu, 2 min feedback (Observer note : vocabulaire finance, modaux \eng{should/must}, politesse, prononciation).
\textbf{Phrases utiles :} \eng{“Have you considered...?”, “You should...”, “If I were you, I would...”, “What is your monthly income?”, “The interest rate is...”, “I cannot afford the monthly repayments.”}
\end{experience}

\begin{exercice}[Production écrite — Devoir / Examen]
\textbf{Sujet :} \eng{“Many young people in Cameroon run into debt because they confuse needs and wants.” Write a short article (120--150 words) for your school magazine explaining the difference, giving three practical tips to avoid debt, and encouraging your peers to start saving now.}
\textbf{Indications :} Structure : Titre accrocheur -- Introduction (problème) -- Corps (3 conseils + connecteurs) -- Conclusion (appel à l'action). Vocabulaire imposé : \eng{track expenses, emergency fund, needs vs wants, repay, budget, afford}. Modaux : \eng{should, must, have to}.
\end{exercice}

\courssub{Common Errors to Avoid (Pièges du Bac)}

\begin{attention}[Les 7 erreurs qui coûtent des points]
\begin{enumerate}
\item \textbf{Faux amis (False Friends)} :
      \begin{itemize}
      \item \eng{Actually} = « en fait / en réalité » (≠ \eng{currently} = « actuellement »).
      \item \eng{Assist} = « aider » (≠ \eng{attend} = « assister à »).
      \item \eng{Economic} = « relatif à l'économie » (≠ \eng{economical} = « économe, rentable »).
      \item \eng{Eventually} = « finalement / à la fin » (≠ \eng{possibly} = « éventuellement »).
      \item \eng{Demand} = « exiger » (≠ \eng{ask for} = « demander »).
      \end{itemize}
\item \textbf{Confusion \eng{Lend} / \eng{Borrow}} : \eng{I \textbf{borrow} from the bank} (j'emprunte) vs \eng{The bank \textbf{lends} to me} (la banque prête). \eng{Can you \textbf{lend} me your pen?} (Me prêter) / \eng{Can I \textbf{borrow} your pen?} (Emprunter).
\item \textbf{Prépositions verbes argent} : \eng{spend \textbf{on}}, \eng{save \textbf{for}}, \eng{borrow \textbf{from}}, \eng{lend \textbf{to}}, \eng{pay \textbf{for}}, \eng{invest \textbf{in}}. Jamais \eng{spend for}, \eng{borrow to}.
\item \textbf{Noms indénombrables} : \eng{advice} (pas \eng{advices}), \eng{information} (pas \eng{informations}), \eng{money} (pas \eng{moneys}), \eng{furniture}, \eng{luggage}, \eng{equipment}. $\rightarrow$ \eng{a piece of advice, some information}.
\item \textbf{Concordance des temps (Backshifting)} : Question au prétérit $\rightarrow$ Réponse au prétérit. \eng{Why \textbf{did} he save? $\rightarrow$ He \textbf{saved} because...} (Pas \eng{save}).
\item \textbf{3e personne singulier Present Simple} : \eng{He/She/It \textbf{saves} / \textbf{earns} / \textbf{has} / \textbf{does}}. Oublier le \eng{-s} = faute grave.
\item \textbf{Orthographe / Ponctuation} : \eng{Which} (pas \eng{wich}), \eng{Because} (pas \eng{becaus}), \eng{Receive} (règle \eng{i} avant \eng{e} sauf après \eng{c}), \eng{Believe}. Majuscule \eng{I} toujours. Point final obligatoire.
\end{enumerate}
\end{attention}

\courssub{Cultural Insight: Mobile Money in Cameroon}

\begin{savaistu}[Mobile Money : La banque dans la poche]
Au Cameroun, \eng{Mobile Money} (MoMo par MTN, Orange Money) a révolutionné l'inclusion financière.
\begin{itemize}
\item \textbf{Services :} Dépôt/retrait (\eng{cash in/out}), transfert, paiement factures (Eneo, Camwater), \textbf{épargne} (\eng{MoMo Savings, Orange Épargne}) avec intérêts, micro-crédit (\eng{MoMo Kash, Orange Prêt}).
\item \textbf{Vocabulaire spécifique :} \eng{Agent} (le kiosque), \eng{PIN} (code secret), \eng{Transaction ID}, \eng{Fees/Charges} (frais), \eng{Limit} (plafond journalier).
\item \textbf{Conseil sécurité :} Ne jamais communiquer son \eng{PIN} ni son \eng{OTP} (One-Time Password reçu par SMS). Vérifier le nom du destinataire avant de valider.
\item \textbf{Statistique indicative :} Plus de 10 millions de comptes actifs au Cameroun (données 2023). C'est le premier "compte bancaire" pour beaucoup de lycéens et étudiants.
\end{itemize}
\textbf{Activité :} Comparez les frais d'un virement MoMo vs virement bancaire classique pour 10.000 FCFA. Discutez : \eng{“Is Mobile Money a good tool for a student's first budget?”}
\end{savaistu}

\begin{corrige}[Exercice Compréhension (p. 4)]
\begin{enumerate}
\item Jean-Paul is a financial coach in Douala. He advises tracking every franc to know where money goes.
\item It means saving money (e.g., 10\% of income) immediately after receiving it, before spending on anything else (airtime, snacks).
\item A torn school uniform, a sudden illness, a rise in transport fares.
\item It traps you in a cycle of repayment for non-essential items.
\item \eng{Open answer. Ex: Yes, because 15.000 FCFA allows 7.500 for needs, 4.500 for wants, 3.000 savings. / No, because needs (transport, food) often exceed 50\%.}
\end{enumerate}
\end{corrige}

\begin{corrige}[Exercice Vocabulaire (p. 5)]
\begin{enumerate}
\item 1. \eng{I cannot afford this trip.} 2. \eng{She borrowed money from her sister.} 3. \eng{We must repay the loan.} 4. \eng{He tracks his expenses on an app.} 5. \eng{Saving is a good habit.}
\item \eng{a) on, b) for, c) from, d) to, e) for, f) in.}
\item \eng{a) expenses (autres = revenus), b) afford (autres = prêt/emprunt), c) extravagant (autres = économe), d) savings (autres = dettes).}
\end{enumerate}
\end{corrige}

\begin{corrige}[Exercice Grammaire Modaux (p. 6)]
\begin{enumerate}
\item \eng{should} (conseil)
\item \eng{have to} (règle externe)
\item \eng{mustn't} (interdiction)
\item \eng{had to} (passé)
\item \eng{don't have to} (absence obligation)
\item \eng{Do / have to} (question obligation externe)
\end{enumerate}
\end{corrige}

\newpage

\coursec{Exercices}

\coursec{Lesson 20 — Reading: Managing Personal Finances}

\courssub{Warm-up and Pre-reading}

\begin{experience}[Speaking: Money Habits in Cameroon]
Travaillez en binôme. Discutez des questions suivantes en anglais. Prenez des notes pour un retour en classe.
\begin{enumerate}
\item How do you usually receive your pocket money? (Cash? Mobile money? Bank transfer?)
\item What are your three main \textbf{expenses} each week? (Food, transport, phone credit, snacks, clothes…)
\item Do you \textbf{save} money? If yes, where do you keep it? (Tontine, mobile money wallet, piggy bank, bank account?)
\item Have you ever \textbf{borrowed} money from a friend or a family member? Did you \textbf{pay it back} on time?
\item What is \textbf{Mobile Money} (Orange Money, MTN MoMo)? How has it changed daily life in your neighbourhood?
\end{enumerate}
\end{experience}

\begin{methode}[Pre-reading Strategy: Predicting Content]
Avant de lire un texte, observez le titre, les images, les mots en gras et le premier paragraphe.
\begin{enumerate}
\item \textbf{Title:} “Managing Your Money: A Guide for Young Cameroonians”. What do you expect to read about?
\item \textbf{Keywords:} \cle{budget}, \cle{income}, \cle{expenses}, \cle{savings}, \cle{debt}, \cle{tontine}, \cle{mobile money}, \cle{withdraw}. Which ones do you know? Guess the others.
\item \textbf{Questions:} Write two questions you hope the text will answer.
\end{enumerate}
\end{methode}

\courssub{Reading Text with Glossary}

\begin{textebox}[Managing Your Money: A Guide for Young Cameroonians]
Managing personal finances is not only for adults. In Cameroon, whether you are a student in Yaoundé, an apprentice in Douala or a young trader in Buea, knowing how to handle money is a vital life skill. Many young people receive \textbf{pocket money} from their parents or earn a small \textbf{income} from holiday jobs, like selling \textbf{groundnuts} or \textbf{puff-puff} at the market. The first rule is simple: you must know where your money goes.

Start by making a \textbf{monthly budget}. Write down your total \textbf{income} (salary, pocket money, gifts). Then list all your \textbf{expenses}: \textbf{school fees}, transport (bendskin or taxi fares), food, phone credit and entertainment. If your expenses are higher than your income, you have a \textbf{deficit}. You must reduce unnecessary spending — do you really need that new pair of sneakers or that extra data bundle?

\textbf{Saving} is the key to financial security. Even 500 francs CFA put aside every week grows into a useful sum. In many communities, the \textbf{tontine} (or \emph{njangi}) remains a popular way to save. Members contribute a fixed amount regularly, and each member receives the total pot in turn. It requires discipline and trust. Alternatively, opening a \textbf{bank account} at a microfinance institution or using a \textbf{mobile money} wallet (MoMo, Orange Money) allows you to \textbf{deposit} and \textbf{withdraw} cash safely, pay \textbf{bills} (electricity, water) and \textbf{transfer} money to family in Bamenda or Bertoua instantly.

Beware of \textbf{debt}. Borrowing money to buy \textbf{assets} (like a motorbike for a taxi-moto business) can be an investment. But borrowing for \textbf{consumption} (clothes, parties, gambling) leads to a debt trap. Never borrow from \textbf{loan sharks} who charge exorbitant interest rates. Finally, protect your \textbf{PIN} and \textbf{OTP} (One-Time Password). Scammers send fake messages pretending to be your operator. Remember: your bank or mobile money provider will never ask for your secret code by SMS.
\end{textebox}

\begin{definition}[Glossary / Lexique du texte]
\begin{center}
\begin{tabularx}{\linewidth}{|X|X|X|}
\hline
\rowcolor{popblueL} \textbf{English} & \textbf{Part of Speech} & \textbf{Français} \\
\hline
pocket money & noun (U) & argent de poche \\
income & noun (C/U) & revenu, rentrée d'argent \\
expenses & noun (pl) & dépenses \\
school fees & noun (pl) & frais de scolarité \\
budget & noun (C) & budget \\
deficit & noun (C) & déficit \\
savings & noun (pl) & économies \\
tontine / njangi & noun (C) & tontine (épargne tournante) \\
mobile money & noun (U) & mobile money (argent mobile) \\
deposit (v.) & verb & déposer (de l'argent) \\
withdraw (v.) & verb & retirer (de l'argent) \\
transfer (v.) & verb & transférer, virer \\
bill & noun (C) & facture \\
debt & noun (C/U) & dette \\
asset & noun (C) & bien, actif (qui rapporte) \\
consumption & noun (U) & consommation (biens non durables) \\
loan shark & noun (C) & usurier \\
interest rate & noun (C) & taux d'intérêt \\
PIN / OTP & noun (C) & code secret / code à usage unique \\
scammer & noun (C) & arnaqueur, escroc \\
\hline
\end{tabularx}
\end{center}
\end{definition}

\courssub{Comprehension Activities}

\courssub{Skimming (Reading for Gist)}
\begin{exercice}[Main Idea]
Lisez le texte rapidement (1 minute). Choisissez la meilleure réponse.
\begin{enumerate}
\item The main purpose of the text is to:
\begin{itemize}
\item[A.] explain the history of banking in Cameroon.
\item[B.] advise young people on how to manage their money.
\item[C.] advertise a specific mobile money company.
\item[D.] describe the dangers of tontines.
\end{itemize}
\item The tone of the text is:
\begin{itemize}
\item[A.] humorous and sarcastic.
\item[B.] formal and academic.
\item[C.] practical and advisory.
\item[D.] critical and negative.
\end{itemize}
\end{enumerate}
\end{exercice}

\courssub{Scanning (Reading for Specific Information)}
\begin{exercice}[Find the Information]
Retrouvez dans le texte les informations exactes (copiez la phrase ou l'expression).
\begin{enumerate}
\item Two examples of holiday jobs for students.
\item The definition of a \textbf{deficit}.
\item Two examples of \textbf{bills} you can pay via mobile money.
\item The difference between borrowing for an \textbf{asset} and borrowing for \textbf{consumption}.
\item What a \textbf{scammer} might do to steal your money.
\end{enumerate}
\end{exercice}

\courssub{Detailed Comprehension}
\begin{exercice}[True, False or Not Mentioned?]
Lisez attentivement. Cochez \textbf{T} (True), \textbf{F} (False) ou \textbf{NM} (Not Mentioned). Justifiez par une citation du texte.
\begin{enumerate}
\item The text says most students in Yaoundé have a bank account.
\item A tontine allows members to receive a lump sum periodically.
\item Mobile money is only useful for people who live in big cities.
\item The text advises against borrowing money to start a small business.
\item You should give your OTP to a customer service agent if they call you.
\end{enumerate}
\end{exercice}

\courssub{Thematic Vocabulary}

\begin{definition}[Core Vocabulary: Personal Finances]
\begin{center}
\begin{tabularx}{\linewidth}{|X|X|X|}
\hline
\rowcolor{popblueL} \textbf{Category} & \textbf{English} & \textbf{Français} \\
\hline
\textbf{Money Verbs} & earn, spend, save, waste, borrow, lend, owe, pay back, deposit, withdraw, transfer, invest, budget (v.) & gagner, dépenser, économiser, gaspiller, emprunter, prêter, devoir, rembourser, déposer, retirer, virer, investir, budgétiser \\
\hline
\textbf{Banking \& Budget Nouns} & income, salary, wages, pocket money, expenses, bills, budget, savings, account, balance, statement, receipt, transaction, deficit, surplus & revenu, salaire, salaire horaire, argent de poche, dépenses, factures, budget, économies, compte, solde, relevé, reçu, transaction, déficit, excédent \\
\hline
\textbf{Debt \& Loans} & debt, loan, borrower, lender, interest (rate), instalment, collateral, guarantor, loan shark, credit score & dette, prêt, emprunteur, prêteur, intérêt (taux), mensualité, garantie, caution, usurier, cote de crédit \\
\hline
\textbf{Cameroon Context} & tontine / njangi, mobile money (MoMo, Orange Money), agent / kiosk, bendskin, school fees, groundnuts, puff-puff, CSC (Caisse Sociale) & tontine, mobile money, agent/kiosque, moto-taxi, frais de scolarité, arachides, beignets, caisse sociale \\
\hline
\textbf{Adjectives} & thrifty / frugal, careless / reckless, broke, solvent, bankrupt, affordable, expensive, priceless & économe, insouciant, fauché, solvable, en faillite, abordable, cher, inestimable \\
\hline
\end{tabularx}
\end{center}
\end{definition}

\begin{attention}[Faux Amis \& Pièges Fréquents (Francophones)]
\begin{itemize}
\item \textbf{Borrow vs Lend:} \eng{I borrow from you} (j'emprunte à toi) / \eng{You lend to me} (tu me prêtes). \textbf{Pas} “I lend from you”.
\item \textbf{Spend vs Waste:} \eng{Spend} = dépenser (neutre). \eng{Waste} = gaspiller (négatif, inutile).
\item \textbf{Salary vs Wages:} \eng{Salary} = salaire mensuel fixe. \eng{Wages} = salaire horaire/hebdomadaire (souvent manuel).
\item \textbf{Coin vs Change:} \eng{Coin} = pièce de monnaie (métal). \eng{Change} = monnaie (rendu) / petite monnaie.
\item \textbf{Benefit vs Profit:} \eng{Benefit} = avantage, prestation (sociale). \eng{Profit} = bénéfice (commercial).
\item \textbf{Prepositions:} \eng{spend money \textbf{on} sth}, \eng{pay \textbf{for} sth}, \eng{borrow \textbf{from} sb}, \eng{lend \textbf{to} sb}, \eng{save \textbf{for} sth}, \eng{invest \textbf{in} sth}.
\item \textbf{State Verbs (pas de continuous):} \eng{owe, own, possess, cost, worth, need, want, mean, understand, know, believe}. \eng{He owes me money} (pas “is owing”).
\end{itemize}
\end{attention}

\begin{exemplebox}[Collocations Utiles]
\begin{itemize}
\item \eng{make a budget} (établir un budget) — \textbf{pas} “do a budget”.
\item \eng{keep track of expenses} (suivre ses dépenses).
\item \eng{live within one's means} (vivre selon ses moyens).
\item \eng{run out of money} (ne plus avoir d'argent).
\item \eng{make ends meet} (joindre les deux bouts).
\item \eng{a waste of money} (une perte d'argent / gaspillage).
\end{itemize}
\end{exemplebox}

\courssub{Reading Method and Worked Example}

\begin{methode}[Reading Strategies: Skimming \& Scanning]
\textbf{Skimming (Lecture rapide / globale)} — Objectif : saisir l'idée générale.
\begin{enumerate}
\item Lisez le \textbf{titre}, le \textbf{premier} et le \textbf{dernier paragraphe}.
\item Repérez les \textbf{mots-clés} en gras, les noms propres, les chiffres.
\item Ne vous arrêtez \textbf{pas} sur les mots inconnus.
\item Répondez : “De quoi parle le texte ? Quel est le but de l'auteur ?”
\end{enumerate}

\textbf{Scanning (Recherche ciblée)} — Objectif : trouver une info précise (date, nom, chiffre, mot).
\begin{enumerate}
\item Identifiez le \textbf{mot-clé} de la question (ex: “school fees”, “tontine”).
\item Faites glisser les yeux verticalement sur le texte pour localiser ce mot.
\item Lisez \textbf{uniquement la phrase concernée} pour extraire la réponse.
\end{enumerate}

\textbf{Guessing Meaning from Context (Déduction)} — Objectif : comprendre un mot inconnu sans dictionnaire.
\begin{enumerate}
\item Regardez la \textbf{classe de mots} (nom, verbe, adjectif ?).
\item Cherchez des \textbf{indices} : définition (“is a system where…”), exemple (“such as…”), contraste (“unlike…”, “but”), cause/effet (“because”, “so”).
\item Remplacez le mot par votre hypothèse : la phrase a-t-elle du sens ?
\end{enumerate}
\end{methode}

\begin{exemplebox}[Worked Example: Guessing “Thrifty” and “Loan Shark”]
\textbf{Contexte (inventé pour l'exemple) :} “Amina is \textbf{thrifty}; she never buys things she does not need. Unlike her brother who is a \textbf{spendthrift}, she saves 500 francs weekly.” \\
\textbf{Analyse :} “Thrifty” est un adjectif (après \textit{is}). Contraste marqué par “Unlike… spendthrift” (qui dépense tout). \textbf{Déduction :} \textit{thrifty} = économe, qui dépense peu.

\textbf{Contexte (texte principal) :} “Never borrow from \textbf{loan sharks} who charge exorbitant interest rates.” \\
\textbf{Analyse :} “Loan shark” est un nom (après \textit{from}). Relatif “who charge exorbitant interest rates” = qui exigent des taux exorbitants. \textbf{Déduction :} \textit{loan shark} = usurier (prêteur illégal/abusif).
\end{exemplebox}

\courssub{Guided Production: Summary and Opinion}

\begin{methode}[Writing a Summary (Résumer un texte)]
\begin{enumerate}
\item Identifiez l'\textbf{idée principale} de chaque paragraphe (1 phrase max).
\item Sélectionnez les \textbf{informations essentielles} (qui, quoi, comment, pourquoi). Supprimez les exemples et détails.
\item Reliez les idées avec des \textbf{connecteurs logiques} (\eng{First, Then, However, Finally, Moreover, Consequently}).
\item Réécrivez \textbf{avec vos propres mots} (paraphrase), sans copier-coller.
\item Vérifiez : nombre de mots imposé, temps (souvent \textbf{present simple}), objectivité (pas d'opinion personnelle).
\end{enumerate}
\end{methode}

\begin{exercice}[Guided Summary]
Complétez le résumé du texte \emph{Managing Your Money} en utilisant les connecteurs donnés.
\begin{center}
\begin{tabularx}{\linewidth}{|X|}
\hline
\rowcolor{popblueL} \textbf{Summary} \\
\hline
\textbf{First}, managing money is a vital skill for young Cameroonians. \textbf{Then}, you should create a monthly budget listing your income and expenses to avoid a deficit. \textbf{Moreover}, saving regularly, whether through a tontine or a mobile money wallet, builds financial security. \textbf{However}, beware of debt: borrowing for consumption leads to a debt trap, unlike borrowing for an asset. \textbf{Finally}, protect your PIN and OTP from scammers who send fake messages. \\
\hline
\end{tabularx}
\end{center}
\textbf{Votre tour :} Réécrivez ce résumé en \textbf{4 phrases maximum} (environ 50 mots) en remplaçant le vocabulaire par des synonymes (ex: \textit{vital skill} $\to$ \textit{essential ability}; \textit{builds} $\to$ \textit{ensures}).
\end{exercice}

\begin{exercice}[Personal Opinion / Expression Écrite Guidée]
Répondez en anglais (3--4 phrases) à l'une des questions. Utilisez \eng{I think / In my opinion / Because / For example}.
\begin{enumerate}
\item Is a tontine better than a bank account for a student? Why?
\item “Mobile money makes spending too easy.” Do you agree?
\item What is the most difficult thing about saving money for you?
\end{enumerate}
\end{exercice}

\courssub{Grammar Focus: Present Simple vs Present Continuous}

\begin{propriete}[Present Simple vs Present Continuous]
\begin{center}
\begin{tabularx}{\linewidth}{|X|X|}
\hline
\rowcolor{popblueL} \textbf{Present Simple} & \textbf{Present Continuous} \\
\hline
\textbf{Forme:} Base verbale (V) / V+\textbf{s} (3\textsuperscript{e} pers. sg.) & \textbf{Forme:} \textit{be} (am/is/are) + V-\textbf{ing} \\
\hline
\textbf{Emplois :} & \textbf{Emplois :} \\
1. Habitudes, routines (\textit{every day, usually, often}). & 1. Action en cours \textbf{maintenant} (\textit{now, at the moment, listen!}). \\
2. Vérités générales, faits permanents. & 2. Situation temporaire \textbf{cette période} (\textit{this week, this year, currently}). \\
3. Verbes d'état (\textit{know, like, need, want, owe, cost, mean, understand, believe, own, possess}). & 3. Projet futur planifié (\textit{tomorrow, next Monday} + verbe de mouvement/action). \\
4. Programmes, horaires (\textit{The train leaves at 6}). & 4. Tendances, changements (\textit{More and more people are using...}). \\
\hline
\textbf{Mots-signaux :} always, usually, often, sometimes, never, every day/week/month, on Mondays. & \textbf{Mots-signaux :} now, right now, at the moment, currently, this week/year, today, Look! Listen! \\
\hline
\textbf{Attention :} Verbes d'état $\to$ \textbf{Jamais} de \textit{-ing}. & \textbf{Attention :} \textit{Have} (posséder) = state verb. \textit{Have} (action: eat, take a shower) = action verb. \\
\eng{I have a phone.} / \eng{I'm having lunch.} & \\
\hline
\end{tabularx}
\end{center}
\end{propriete}

\begin{exemplebox}[Exemples Contrastés]
\begin{itemize}
\item \eng{My father \textbf{works} in a bank.} (Fait permanent) — \eng{This month, he \textbf{is working} in the Buea branch.} (Temporaire)
\item \eng{Students usually \textbf{receive} pocket money on Mondays.} (Habitude) — \eng{Listen! The teacher \textbf{is explaining} the budget rule.} (Maintenant)
\item \eng{He \textbf{owes} me 5,000 francs.} (État — \textbf{pas} “is owing”) — \eng{He \textbf{is paying} me back today.} (Action en cours)
\item \eng{Mobile money \textbf{changes} lives.} (Tendance générale) — \eng{More traders \textbf{are using} it \textbf{nowadays}.} (Changement en cours)
\end{itemize}
\end{exemplebox}

\courssub{Je vérifie mes connaissances}

\begin{exercice}[QCM : Le bon mot financier]
Pour chaque phrase, choisis la bonne réponse (A, B ou C).
\begin{enumerate}
\item I have to \textbf{......} money from the bank to pay my school fees.\\
A. lend \hspace{1cm} B. borrow \hspace{1cm} C. spend
\item My father \textbf{......} part of his salary every month.\\
A. saves \hspace{1cm} B. wastes \hspace{1cm} C. owes
\item Transport and food are my main \textbf{......}.\\
A. income \hspace{1cm} B. expenses \hspace{1cm} C. profit
\item She earns a monthly \textbf{......} of 150,000 francs.\\
A. debt \hspace{1cm} B. income \hspace{1cm} C. bill
\item Can you \textbf{......} me 2,000 francs until Friday? I will pay you back.\\
A. lend \hspace{1cm} B. borrow \hspace{1cm} C. steal
\item A \textbf{......} helps you plan how to spend your money.\\
A. budget \hspace{1cm} B. wallet \hspace{1cm} C. price
\item If you spend more than you earn, you will be in \textbf{......}.\\
A. savings \hspace{1cm} B. debt \hspace{1cm} C. cash
\item Mobile money allows people to \textbf{......} money by phone.\\
A. transfer \hspace{1cm} B. print \hspace{1cm} C. burn
\end{enumerate}
\end{exercice}

\begin{exercice}[Vrai ou Faux ?]
Lis les phrases suivantes et écris \textbf{True} ou \textbf{False}. Justifie chaque réponse par une courte phrase en anglais.
\begin{enumerate}
\item “To borrow” means to give money to somebody for a short time.
\item Your “income” is the money you receive regularly, for example a salary.
\item “Expenses” are the things you spend money on.
\item A “tontine” is a system where members save money together and take turns to receive the total.
\item If you “save” money, you spend it immediately.
\item “To pay back” means to return money that you borrowed.
\end{enumerate}
\end{exercice}

\begin{exercice}[Vocabulaire : Traduction]
Traduis en français les mots et expressions suivants, extraits du texte de la leçon.
\begin{enumerate}
\item pocket money
\item to manage money
\item a bank account
\item to withdraw cash
\item school fees
\item to waste money
\end{enumerate}
\end{exercice}

\begin{exercice}[Grammaire : Complète avec le bon temps]
Complète les phrases avec le verbe entre parenthèses au \textbf{present simple} ou au \textbf{present continuous}.
\begin{enumerate}
\item Every month, my mother ...... (put) some money in the tontine.
\item Listen! The banker ...... (explain) how to open an account.
\item Students usually ...... (receive) their pocket money on Mondays.
\item I ...... (save) for a new phone this year, so I don't buy snacks anymore.
\item Right now, many traders in Douala ...... (use) mobile money to pay their suppliers.
\item He never ...... (spend) more than he earns.
\end{enumerate}
\end{exercice}


\newpage

\coursec{Corrigés des exercices}

\begin{corrige}[Exercice 1 — Main Idea (Skimming)]
\textbf{1. Réponse : B} — \eng{The main purpose of the text is to advise young people on how to manage their money.} \\
\textbf{Justification :} le titre \eng{A Guide for Young Cameroonians} annonce un guide de conseils ; chaque paragraphe donne une règle pratique (\eng{Start by making a monthly budget}, \eng{Saving is the key}, \eng{Beware of debt}). Les options A (histoire des banques), C (publicité) et D (dangers des tontines) ne correspondent à aucun paragraphe.

\textbf{2. Réponse : C} — \eng{The tone of the text is practical and advisory.} \\
\textbf{Justification :} le texte utilise l'impératif (\eng{Start by…}, \eng{Write down…}, \eng{Never borrow…}, \eng{Remember…}) et s'adresse directement au lecteur (\eng{you}), ce qui est typique d'un texte de conseils pratiques. Il n'est ni humoristique (A), ni académique (B), ni négatif (D).

\textbf{Point de vigilance :} en skimming, ne cherchez pas à tout comprendre : le titre, le premier paragraphe et les impératifs suffisent ici à identifier le but et le ton. Relire ensuite le premier et le dernier paragraphe permet de confirmer l'hypothèse avant de valider la réponse.
\end{corrige}

\begin{corrige}[Exercice 2 — Find the Information (Scanning)]
\textbf{1. Two examples of holiday jobs:} \eng{“selling groundnuts or puff-puff at the market”} (vendre des arachides ou des beignets au marché).

\textbf{2. Definition of a deficit:} \eng{“If your expenses are higher than your income, you have a deficit.”} (Si vos dépenses sont supérieures à vos revenus, vous êtes en déficit.)

\textbf{3. Two examples of bills paid via mobile money:} \eng{“electricity, water”} (l'électricité, l'eau) — phrase complète : \eng{“…pay bills (electricity, water)…”}.

\textbf{4. Asset vs consumption:} \eng{“Borrowing money to buy assets (like a motorbike for a taxi-moto business) can be an investment. But borrowing for consumption (clothes, parties, gambling) leads to a debt trap.”} (Emprunter pour un bien qui rapporte est un investissement ; emprunter pour consommer mène au piège de l'endettement.)

\textbf{5. What a scammer might do:} \eng{“Scammers send fake messages pretending to be your operator.”} (Les arnaqueurs envoient de faux messages en se faisant passer pour votre opérateur.)

\textbf{Point de vigilance :} en scanning, copiez la phrase \emph{exactement} comme dans le texte, sans la reformuler : c'est ce qui prouve que vous avez localisé l'information. Repérez d'abord les mots-repères de la question (noms propres, chiffres, mots techniques), puis laissez votre œil glisser sur le texte jusqu'à les retrouver.
\end{corrige}

\begin{corrige}[Exercice 3 — True, False or Not Mentioned?]
\textbf{1. NM (Not Mentioned).} Le texte dit \eng{“whether you are a student in Yaoundé…”} mais ne précise jamais que la plupart des élèves ont un compte en banque. Attention : ne confondez pas « mentionné dans le texte » et « vrai dans la réalité ».

\textbf{2. T (True).} \eng{“Members contribute a fixed amount regularly, and each member receives the total pot in turn.”} — \emph{a lump sum periodically} correspond exactement à \emph{receives the total pot in turn}.

\textbf{3. F (False).} Le texte indique que le mobile money permet de \eng{“transfer money to family in Bamenda or Bertoua instantly”} ; il ne limite jamais son utilité aux grandes villes — bien au contraire, il rapproche les zones éloignées des banques.

\textbf{4. F (False).} \eng{“Borrowing money to buy assets (like a motorbike for a taxi-moto business) can be an investment.”} Le texte déconseille d'emprunter pour la \emph{consommation}, pas pour créer une activité.

\textbf{5. F (False).} \eng{“Your bank or mobile money provider will never ask for your secret code by SMS.”} Il ne faut donc jamais communiquer son OTP, même à un soi-disant agent.

\textbf{Méthode :} pour chaque item, citez la ligne du texte qui justifie votre choix ; sans citation, la réponse n'est pas validée en examen. Rappel : \emph{Not Mentioned} signifie que le texte ne dit ni oui ni non — même si l'affirmation semble vraie dans la vie réelle.
\end{corrige}

\begin{corrige}[Exercice 4 — Guided Summary]
\textbf{Proposition de résumé en 4 phrases (environ 50 mots) :}

\eng{Managing money is an essential ability for young people in Cameroon. First, drawing up a monthly budget of earnings and spending prevents a deficit. Moreover, putting money aside regularly, in a tontine or a mobile wallet, ensures financial security. Finally, avoid borrowing for consumption and never share your PIN or OTP with anyone.}

(52 mots — « Gérer son argent est une compétence essentielle pour les jeunes du Cameroun. D'abord, établir un budget mensuel des revenus et des dépenses évite le déficit. De plus, mettre de l'argent de côté régulièrement, dans une tontine ou un portefeuille mobile, garantit la sécurité financière. Enfin, évitez d'emprunter pour consommer et ne communiquez jamais votre code secret à personne. »)

\textbf{Synonymes utilisés :} \emph{vital skill} $\to$ \emph{essential ability} ; \emph{income} $\to$ \emph{earnings} ; \emph{saving} $\to$ \emph{putting money aside} ; \emph{builds} $\to$ \emph{ensures} ; \emph{beware of} $\to$ \emph{avoid}.

\textbf{Barème indicatif (type Bac, /5) :} respect de la limite de mots (1 pt) ; paraphrase sans copier-coller (1 pt) ; idées essentielles des 4 paragraphes (2 pts) ; correction grammaticale et connecteurs (1 pt).

\textbf{Points de vigilance :} présent simple partout ; pas d'opinion personnelle dans un résumé ; ne pas écrire \eng{“a 50 words summary”} mais \eng{“a 50-word summary”} (adjectif composé sans \emph{s}).
\end{corrige}

\begin{corrige}[Exercice 5 — Personal Opinion]
\textbf{Proposition de rédaction modèle (question 1) :}

\eng{In my opinion, a tontine is better than a bank account for a student because it is simple and free. For example, in my school, each member of our njangi contributes 500 francs every week, and we receive the total pot in turn. I think this system teaches us discipline, whereas a bank account often has fees. However, a bank account is safer for large sums of money.}

(« À mon avis, une tontine est mieux qu'un compte en banque pour un élève car c'est simple et gratuit. Par exemple, dans mon école, chaque membre de notre njangi cotise 500 francs par semaine et nous recevons la cagnotte à tour de rôle. Je pense que ce système nous apprend la discipline, alors qu'un compte en banque a souvent des frais. Cependant, un compte en banque est plus sûr pour les grosses sommes. »)

\textbf{Modèle court (question 2) :} \eng{I agree that mobile money makes spending too easy. Because the money is on the phone, we buy credit or data bundles without thinking. In my opinion, we should transfer only a fixed amount to our wallet every week.}

(« Je suis d'accord pour dire que le mobile money rend les dépenses trop faciles. Comme l'argent est sur le téléphone, nous achetons du crédit ou des forfaits internet sans réfléchir. À mon avis, nous devrions ne transférer qu'un montant fixe sur notre portefeuille chaque semaine. »)

\textbf{Barème indicatif (/5) :} prise de position claire avec \eng{I think / In my opinion} (1 pt) ; justification avec \eng{because} (1 pt) ; exemple avec \eng{for example} (1 pt) ; grammaire et vocabulaire du thème (2 pts).

\textbf{Points de vigilance :} \eng{I agree with you} mais \eng{I agree that…} (pas de \emph{with} devant une proposition) ; \eng{money} est indénombrable : \eng{too much money}, jamais “too much moneys” ; \eng{advice} est également indénombrable : \eng{a piece of advice}, jamais “an advice”.
\end{corrige}

\begin{corrige}[Exercice 6 — QCM : Le bon mot financier]
\textbf{1. B — borrow.} \eng{I have to borrow money from the bank…} On \emph{emprunte à} la banque : \eng{borrow from}. \emph{Lend} (prêter) aurait le sens inverse ; \emph{spend} (dépenser) ne convient pas car on paie ensuite les frais.

\textbf{2. A — saves.} \eng{My father saves part of his salary every month.} \emph{Every month} = habitude, present simple avec \emph{-s} à la 3\textsuperscript{e} personne. \emph{Wastes} (gaspille) et \emph{owes} (doit) contredisent l'idée positive.

\textbf{3. B — expenses.} \eng{Transport and food are my main expenses.} Transport et nourriture sont des \emph{dépenses}, pas des revenus (\emph{income}) ni un bénéfice (\emph{profit}).

\textbf{4. B — income.} \eng{She earns a monthly income of 150,000 francs.} Le verbe \emph{earns} (gagne) impose \emph{income} (revenu). On ne « gagne » ni une dette ni une facture.

\textbf{5. A — lend.} \eng{Can you lend me 2,000 francs until Friday?} La personne interrogée \emph{donne} l'argent : \eng{lend somebody something}. \emph{Borrow} serait : \eng{Can I borrow 2,000 francs from you?}

\textbf{6. A — budget.} \eng{A budget helps you plan how to spend your money.} C'est la définition même du budget.

\textbf{7. B — debt.} \eng{If you spend more than you earn, you will be in debt.} Expression figée : \eng{to be in debt} (être endetté). \emph{Savings} et \emph{cash} sont positifs.

\textbf{8. A — transfer.} \eng{Mobile money allows people to transfer money by phone.} Structure : \eng{allow somebody to do something}.

\textbf{Point de vigilance :} la paire \emph{borrow / lend} est le piège classique des francophones : \eng{borrow something from somebody} (emprunter à) et \eng{lend something to somebody} (prêter à). Le français « prêter » se traduit par \emph{lend}, jamais par \emph{borrow}.
\end{corrige}

\begin{corrige}[Exercice 7 — Vrai ou Faux ?]
\textbf{1. False.} \eng{To borrow means to take or receive money from somebody; to lend means to give money for a short time.} (Emprunter = recevoir ; prêter = donner.)

\textbf{2. True.} \eng{Your income is the money you receive regularly, for example a salary, wages or pocket money.}

\textbf{3. True.} \eng{Expenses are the things you spend money on, such as food, transport and bills.}

\textbf{4. True.} \eng{In a tontine, members contribute money regularly and receive the total pot in turn.}

\textbf{5. False.} \eng{To save money means to keep it for the future, not to spend it immediately.} (Save = mettre de côté ; spend = dépenser.)

\textbf{6. True.} \eng{To pay back means to return the money that you borrowed.} Attention : \emph{pay back} est un phrasal verb séparable : \eng{I will pay you back} / \eng{I will pay the money back}.

\textbf{Point de vigilance :} la justification doit être une phrase complète en anglais, pas un seul mot. Avec un pronom, le phrasal verb séparable impose l'ordre verbe + pronom + particule : \eng{I will pay it back}, jamais “I will pay back it”.
\end{corrige}

\begin{corrige}[Exercice 8 — Vocabulaire : Traduction]
\begin{enumerate}
\item \eng{pocket money} : argent de poche.
\item \eng{to manage money} : gérer son argent. (\emph{Manage} = gérer, administrer ; faux ami : \emph{manager} = le directeur, le gérant — et non « le manager » au sens français d'entraîneur sportif, qui se dit \emph{coach}.)
\item \eng{a bank account} : un compte en banque.
\item \eng{to withdraw cash} : retirer de l'argent liquide. (Verbe irrégulier : \emph{withdraw – withdrew – withdrawn}.)
\item \eng{school fees} : les frais de scolarité. (Toujours au pluriel en anglais dans ce sens.)
\item \eng{to waste money} : gaspiller de l'argent. (Ne pas confondre avec \emph{to spend} = dépenser, neutre.)
\end{enumerate}
\end{corrige}

\begin{corrige}[Exercice 9 — Grammaire : Present Simple ou Present Continuous]
\textbf{1. puts.} \eng{Every month, my mother puts some money in the tontine.} \emph{Every month} = habitude $\to$ present simple ; 3\textsuperscript{e} personne du singulier $\to$ \emph{put + s} (attention, \emph{put} a la même forme à la base verbale et au prétérit, mais prend bien le \emph{s} au présent : \emph{he puts}).

\textbf{2. is explaining.} \eng{Listen! The banker is explaining how to open an account.} \emph{Listen!} = action en cours maintenant $\to$ present continuous : \emph{is + explain-ing}.

\textbf{3. receive.} \eng{Students usually receive their pocket money on Mondays.} \emph{Usually / on Mondays} = routine $\to$ present simple ; sujet pluriel $\to$ base verbale sans \emph{s}.

\textbf{4. am saving.} \eng{I am saving for a new phone this year, so I don't buy snacks any more.} \emph{This year} = période temporaire en cours $\to$ present continuous. Orthographe : \emph{save} $\to$ \emph{saving} (le \emph{e} final muet tombe devant \emph{-ing}).

\textbf{5. are using.} \eng{Right now, many traders in Douala are using mobile money to pay their suppliers.} \emph{Right now} = action en cours $\to$ \emph{are + using} (\emph{use} $\to$ \emph{using}, même règle du \emph{e} muet).

\textbf{6. spends.} \eng{He never spends more than he earns.} \emph{Never} = habitude $\to$ present simple ; 3\textsuperscript{e} personne $\to$ \emph{spends}. Noter la structure comparative : \eng{more than he earns}.

\textbf{Points de vigilance :} (a) ne jamais mettre les verbes d'état au continuous (\eng{He owes me money}, pas “is owing”) ; (b) le \emph{-s} de la 3\textsuperscript{e} personne disparaît dans les questions et négations : \eng{Does he save? He doesn't save.} ; (c) doubler la consonne finale après voyelle courte accentuée : \emph{run} $\to$ \emph{running}, \emph{sit} $\to$ \emph{sitting} ; (d) mots-repères à mémoriser : \emph{always, usually, every week, never} $\to$ present simple ; \emph{now, at the moment, right now, Listen!, Look!} $\to$ present continuous.
\end{corrige}

\newpage

\coursec{Fiche bilan}

\begin{popfigure}
\begin{tikzpicture}[mindmap, grow cyclic, every node/.style={concept, text=white, font=\small},
  root concept/.append style={concept color=popblue, font=\bfseries\small},
  level 1/.append style={level distance=3.4cm, sibling angle=51.4},
  level 1 concept/.append style={font=\scriptsize}]
\node[root concept]{Managing Personal Finances}
  child[concept color=poporange]{ node[align=center]{Key vocabulary\\budget, income, expenses, savings, debt} }
  child[concept color=popgreen]{ node[align=center]{Reading strategies\\skimming, scanning, guessing meaning} }
  child[concept color=poppurple]{ node[align=center]{Grammar focus\\modals: should, must, have to} }
  child[concept color=poppink]{ node[align=center]{Comprehension\\true/false, wh-questions, references} }
  child[concept color=popgold]{ node[align=center]{Text types\\article, advice column, report} }
  child[concept color=popblue!70]{ node[align=center]{Production\\summary + personal opinion} }
  child[concept color=popmuted]{ node[align=center]{Culture\\money habits in Cameroon and abroad} };
\end{tikzpicture}
\legende{Carte mentale de la leçon : lire un texte sur la gestion des finances personnelles.}
\end{popfigure}

\begin{bilan}
\textbf{1. Lexique clé (à connaître par cœur)}
\begin{itemize}
  \item \eng{income} : revenu ; \eng{expenses / expenditure} : dépenses ; \eng{a budget} : un budget ; \eng{to budget} : budgétiser.
  \item \eng{savings} : économies ; \eng{to save (money)} : économiser ; \eng{to spend (money) on} : dépenser (pour) ; \eng{to waste money} : gaspiller de l'argent.
  \item \eng{a loan} : un prêt ; \eng{to borrow (from)} : emprunter (à) ; \eng{to lend (to)} : prêter (à) ; \eng{a debt} : une dette ; \eng{to be in debt} : être endetté.
  \item \eng{a bank account} : un compte bancaire ; \eng{to withdraw money} : retirer de l'argent ; \eng{to deposit money} : déposer de l'argent.
  \item \eng{a bill} : une facture ; \eng{to pay in cash / by mobile money} : payer en liquide / par mobile money ; \eng{interest} : intérêts.
  \item \eng{to afford} : avoir les moyens de ; \eng{to live within one's means} : vivre selon ses moyens ; \eng{to make ends meet} : joindre les deux bouts.
\end{itemize}

\textbf{2. Stratégies de lecture}
\begin{center}
\begin{tabularx}{\linewidth}{|l|X|X|}\hline
\rowcolor{popblueL} Stratégie & English & À faire \\ \hline
\cle{Skimming} & read quickly for the general idea & lire titre, première phrase de chaque paragraphe \\ \hline
\cle{Scanning} & look for specific information & repérer chiffres, noms, dates, mots-clés \\ \hline
\cle{Guessing meaning} & deduce meaning from context & type de mot, préfixes/suffixes, mots voisins \\ \hline
\end{tabularx}
\end{center}

\textbf{3. Grammaire à revoir}
\begin{itemize}
  \item Modaux de conseil et d'obligation : \eng{You should save ten per cent of your income.} « Tu devrais économiser 10 \% de tes revenus. » ; \eng{You must pay your bills on time.} « Tu dois payer tes factures à temps. »
  \item \eng{should} + base verbale (conseil) ; \eng{mustn't} = interdiction, \eng{don't have to} = absence d'obligation.
  \item Prépositions piéges : \eng{spend money ON something}, \eng{borrow FROM}, \eng{lend TO}, \eng{pay FOR}.
\end{itemize}

\textbf{4. Méthode express pour résumer et donner son opinion}
\begin{enumerate}
  \item Identifier le thème et l'idée principale de chaque paragraphe.
  \item Reformuler avec ses propres mots en 3 à 5 phrases : \eng{The text deals with... The writer explains that...}
  \item Donner son avis : \eng{In my opinion... / I think that... / I agree because...}
\end{enumerate}
\end{bilan}

\begin{aretenir}[Checklist avant le Bac]
\begin{itemize}
  \item[$\square$] Je connais le lexique de l'argent : \eng{income, expenses, budget, savings, loan, debt, bill}.
  \item[$\square$] Je distingue \eng{borrow from} (emprunter à) et \eng{lend to} (prêter à).
  \item[$\square$] Je sais faire un \emph{skimming} pour trouver l'idée générale d'un texte.
  \item[$\square$] Je sais faire un \emph{scanning} pour repérer une information précise.
  \item[$\square$] Je sais deviner le sens d'un mot inconnu grâce au contexte.
  \item[$\square$] Je réponds aux questions avec des phrases complètes, pas un seul mot.
  \item[$\square$] Je justifie mes réponses \emph{true/false} par une citation du texte.
  \item[$\square$] J'utilise correctement \eng{should}, \eng{must}, \eng{have to} pour donner un conseil.
  \item[$\square$] Je n'oublie pas les prépositions : \eng{spend on}, \eng{pay for}, \eng{interested in}.
  \item[$\square$] Je sais résumer un texte en 3 à 5 phrases avec mes propres mots.
  \item[$\square$] Je sais exprimer mon opinion : \eng{In my opinion... / As far as I am concerned...}
  \item[$\square$] Je relis ma production pour vérifier l'accord sujet-verbe et le temps.
\end{itemize}
\end{aretenir}

\findecours

\end{document}
